% 牌理（はいり）：手牌の速度と価値を最大化させるための理論
% 牌効率（はいこうりつ）：四面子一雀頭を効率よく完成させるための理論

\documentclass[../nankiru_main.tex]{subfiles}

% override main tex render option
\renewcommand{\renderOption}{0}

\begin{document}

\subfilepreamble{牌理}

「牌理」とは順位を最善させる打牌準則\footnote{「牌理」と「牌効率」の定義には諸説ある。戦術書でもネットでもこの二つの言葉の定義は一致しない。ですので、本ノートでは先頭でこれらの言葉を定義し、これらの定義も他の場所に必ずしも適用するわけではない。}。その準則は以下の複数の要素で決まる：

\begin{itemize}
  \item 手牌の期待値
  \begin{itemize}
    \item 和了までの速度
    \begin{itemize}
      \item 受け入れの枚数
      \begin{itemize}
        \item 切られている枚数
        \item 山にありそう枚数
      \end{itemize}
      \item 良形へ変化できる枚数
      \item 仕掛けできるか否か
      \item 待ちの良さ
    \end{itemize}
    
    \item 打点
    \begin{itemize}
      \item 役の作成
      \item ドラの受け入れ
      \item 立直
      \item 副露
    \end{itemize}
  \end{itemize}
  
  \item 点棒状況
  \begin{itemize}
    \item 供託・本場
    \item 祝儀
  \end{itemize}
  \item 他家の癖・運など
\end{itemize}

単なる聴牌までの速度を最大化させる打牌は「牌効率」\footnote{手なりともいう。}。麻雀では、状況により一つの要素を他より優先することはよくある。断ラスなので速度を犠牲しても満貫や跳満を目指すとか、対面がすごくついているので序盤から安全牌を複数持つとか。

この節の構成は色々な形から、まず門前の前提で受け入れしか考えない牌効率を検討し、徐々に打点・手役・点棒状況などの要素を加えて打牌の選択を探検する。

%===================================================
\subsection{単純搭子}

「搭子」は、三つの数牌が揃ってできた順子になる前にあと一つの数牌が必要する形。こうすると、搭子の枚数は二つ限らないので、この文にまた「単純搭子」と「複合搭子」を定義する。「単純搭子」とは、2枚の数牌だけでできた搭子。一方、「複合搭子」は3枚以上の数牌でできた搭子を指す。例えば、\mj{12m} は単純搭子、\mj{135m} は複合搭子。単純搭子の種類は表~\ref{tab:simple-tatsu} に載っている。この小節は単純搭子と単純搭子同士ができた形を説明する。

\begin{table}[ht]
  \centering
  \begin{tabular}{ll}
    \mj{12m} & 辺張 \\
    \mj{13m} & 嵌張 \\
    \mj{23m} & 両面
  \end{tabular}
  \caption{単純搭子の種類}
  \label{tab:simple-tatsu}
\end{table}

%--------------------------------------------
\subsubsection{辺張}

\begin{table}[ht]
  \centering
  \begin{tabular}{ll}
    強さ & 辺張搭子 \\
    \hline
    孤立 \mj{3m} 以下 & \mj{1224m} \\
    & \mj{1246m} \\
    $\downarrow$ & \mj{1244m} \\
    & \mj{1245m} \\
    単独 \mj{12m} 以下 & \mj{1255m} \\
    & \mj{1256m} \\
    $\downarrow$ & \mj{1266m} \\
    & \mj{1267m} \\
    単独 \mj{12m} 以上 & \mj{12567m} \\
    $\downarrow$ & \mj{12456m} \\
    & \mj{11223m}
    \end{tabular}
    \caption{辺張と他の形との複合形の強さの比較。下に行けば行くほど強い}
  \label{tab:pench}
\end{table}

辺張は搭子（辺張・嵌張・両面）の中にもっとも両面になりにくい搭子。けれど、状況によって辺張は強くなる場合もある。辺張の強さは表~\ref{tab:pench} の通りそのフォローに左右されている\citep{gendai-penchan}。


\subsubsection*{受けにフォローがある場合}

\mj{1246m} の場合、\mj{1m} 切っても受けの \mj{3m} はリャンカンの受けになるので、あまりロスにならない\citep{gendai-penchan}。

\mj{1244m} の場合、\mj{3m} 引いたら頭＋1面子ができるが、\mj{1m} 切っても \mj{3m} がまだ \mj{244m} の受けに含んでいるので、もしくは \mj{3m} 引いても1面子＋孤立の \mj{4m} として働ける\footnote{要注意のは聴牌まで頭ができないとき振聴になるリスク。}\citep{gendai-penchan}。


\subsubsection*{変化にフォロー搭子がある場合}

\mj{1255m} と \mj{1256m} の場合、嵌張に変化する時が必要な \mj{4m} は他の搭子とさらに強い形になるので、搭子オーバーフローの時に辺張を払ってもあまり構わない。\mj{1266m} の場合も同じ、\mj{466m} の形が \mj{12m} より少し強い\citep{gendai-penchan}。

\mj{1267m} の場合は、変化の結果の \mj{24m} がさらに \mj{45m} に変化するその受けの \mj{5m} が既に \mj{67m} の受けに含んでいる。だが、この程度は辺張をただ少し弱くするだけ、大きな差は出ない\citep{gendai-penchan}。


\subsubsection*{変化が複合形になる場合}

\mj{12567m} の場合、\mj{4m} 引いたら \mj{24567m} の形になり、\mj{3m} 引いたら既に2面子になるので単純の \mj{4567m} 四連形よりやや強い。\mj{12456m} の場合に \mj{7m} 引いたら上記と同じ\citep{gendai-penchan}。


\subsubsection*{辺張外し}

\textbf{基本的には辺張は孤立牌より残したい}。ツモ2回で面子ができる孤立牌に対して、辺張は一回のツモ抽選で面子ができるので、巡目が深くて聴牌まで余裕がない時、打点上昇の幅が大きくない時、もしくはドラ2枚以上持っている場合は、素直に聴牌に向かうべき\citep[][pp. 48--55]{tetsusenryaku}。辺張外しの有効な場面は下記の通り：

\begin{itemize}
  \item 両向聴\citep{yuusee-389}
  \begin{itemize}
    \item 孤立の37が2種類ある
  \end{itemize}
  \item 一向聴\citep{yuusee-126}
  \begin{itemize}
    \item 辺張が外せばタンヤオがつく
    \item ドラなしで他の部分に中ぶくれや四連形がある
  \end{itemize}
\end{itemize}

辺張と孤立の37の比較は表~\ref{tab:penchan-vs-37} の通り\citep{yuusee-389}。両向聴など聴牌までまた遠い階段なら、愚形を外してより良い形や待を狙う。ただし、一向聴の時、立直に近いので聴牌を優先する。よって、手牌\ref{mj:penchan-1}と\ref{mj:penchan-2}は似ているが、\ref{mj:penchan-1}は一向聴なので打 \mj{3p7s} 推奨、\ref{mj:penchan-2}は両向聴なので打 \mj{89p} 推奨。

\begin{table}[ht]
  \centering
  \begin{tabular}{lp{17ex}p{17ex}}
    & メリット & デメリット \\
    \hline
    辺張外し & 良い待ちになりやすい；安全牌や打点上昇の牌に空間 & カン4や6になる可能性も残っている \\
    孤立牌切り & 即リーができる；辺張の縦引き & 辺張待ちはそこそこ
  \end{tabular}
  \caption{辺張と孤立の37の比較}
  \label{tab:penchan-vs-37}
\end{table}

\tehai{234666m 389p 237s 33z}[東1局　北家　4巡目　ドラ \mj{3s}][penchan-1]

\tehai{335666m 389p 237s 33z}[東1局　北家　4巡目　ドラ \mj{3s}][penchan-2]

手牌\ref{mj:penchan-3}も両向聴なので \mj{89p} を外す\footnote{\citet[pp. 48--55]{tetsusenryaku}は「より良い待ちを作ろう」だけを言及した。「両向聴なので孤立牌より辺張を落とす」のは\citet{yuusee-389}のセオリー。}。

\tehai{4678m 3589p 110677s}[東1局　西家　3巡目　ドラ \mj{2m}][penchan-3]

手牌\ref{mj:penchan-4}はドラ3で打点はもう十分である場合は、素直に聴牌に向かうべき。ですが、ソウズの部分が \mj{114067s} になったら伸びやすいので、その時は辺張を外す。

\tehai{4678m 3589p 110677s}[東1局　西家　3巡目　ドラ \mj{1s}][penchan-4]

どっちでもいい状況はドラなしリーのみ一向聴の手\citep{yuusee-126}。例え、\ref{mj:penchan-5}のような浮き牌に両面ができればピンフがつく例や、\ref{mj:penchan-6}のようなピンフの可能性が一切ないの手。\mj{6m} を先に切れば、待ちの \mj{3m} が筋になるので上がり率がほんの少し上がっている。一方、他家が遅く見えて、打点が欲しくて、または孤立牌がドラに絡まっている場合なら、辺張外しもあり。

\tehai{126m 1145567p 4789s}[東1局　西家　5巡目　ドラ \mj{1s}][penchan-5]

\tehai{126m 45567p 488s 444z}[東1局　西家　5巡目　ドラ \mj{1s}][penchan-6]

もう一つ辺張を落とすタイミングは6ブロックの手牌を5ブロックに絞るとき。6ブロック理論に関しては第~\ref{subsec:6block} 節参照。

%--------------------------------------------
\subsubsection{嵌張}

\subsubsection*{外嵌張と内嵌張}

外嵌張は両面変化できる受け入れが4枚だけの嵌張、つまり \mj{13m}、\mj{24m}、\mj{68m}、\mj{79m}。外嵌張同士の比較基準は辺張の比較基準と同じ、両面変化の受けにフォローがあれば弱くなり、変化した形がフォローと組むと強くなる\citep{gendai-sotokanchan}。

内嵌張は両面変化できる受け入れが4枚だけの嵌張、つまり \mj{35m}、\mj{46m}、\mj{57m}。35(57)と46の差は両嵌張になる時に出る。46に2を引いても両面変化にできるツモは7だけ。一方35に7を引けば357となりそのときツモ28で両面変化ができる。さらに35に1を引いて135の形は筋にかかったカン2の待ちもできる。よって、35(57)は46より残す\citep{yuusee-kanchan}。

ただし、タンヤオが確定な場面なら、46は28のどちらを引いても有効牌になるのに対して、35(57)は1(9)引いてもあまりうれしくないので、その時は46を残す。例え\ref{mj:kanchan-1}は \mj{57s} を払う。

\tehai{2223344467m 46p 57s}[][kanchan-1]

\begin{table}[ht]
  \centering
  \begin{tabular}{ll}
    強さ & 嵌張搭子 \\
    \hline
    単独 \mj{13m} 以下 & \mj{1346m} \\
    & \mj{1344m} \\
    & \mj{1124m} \\
    & \mj{1335m} \\
    $\downarrow$ & \mj{1334m} \\
    & \mj{1356m} \\
    & \mj{1366m} \\
    & \mj{1367m} \\
    & \mj{1355m}（頭が二つ以上） \\
    単独 \mj{13m} 以上 & \mj{13678m} \\
    & \mj{1355m}（\mj{5m} が唯一の頭） \\
    & \mj{11233m} \\
    & \mj{13345m} \\
    $\downarrow$ & \mj{12234m} \\
    & \mj{22344m} \\
    & \mj{13567m} \\
    & \mj{13456m}
  \end{tabular}
  \caption{嵌張と他の形との複合形の強さの比較。下に行けば行くほど強い}
  \label{tab:kanch}
\end{table}

\subsubsection*{受けにフォローがある場合}

\mj{1344m} の場合に要注意のは、\mj{1m} 引きで頭＋両面対子になるので他のフォローがある嵌張より残す。

\mj{1124m} の場合、\mj{4m} 切ったら受け入れが一番広いが、\mj{14m} を切ったら \mj{5m} 引きの両面変化は逃される。タンヤオを見て \mj{11m} を落とすことも実戦では多い。


\subsubsection*{他の搭子と組む場合}

手牌の中に \mj{1355m} の形がある上、\mj{5m} が唯一の頭である場合、\mj{1m} 引きで受け入れが \mj{12345m} になる。他の頭が既に存在している場合、\mj{355m} の嵌張対子にするか \mj{135m} のリャンカンにするか場況次第。

\mj{11233m} の形は完成すれば一盃口があり、\mj{134m} 引きで1面子＋嵌張対子ができる。

\mj{13345m} の形は \mj{136m} 引きで1面子＋嵌張対子ができ、ツモ \mj{5m} で1面子＋リャンカンができる。

\mj{12234m} の形は \mj{12456m} 引きで1面子＋複合面子ができる。

\mj{13567m} の形は \mj{4m} 引きで三面張、\mj{58m} でリャンカンができる。ですが、両面変化2種・リャンカン変化2種の内嵌張ほどではない。

\mj{13456m} の形は \mj{4m} で両面変化、\mj{5m} で両面嵌張、\mj{7m} で三面張、\mj{8m} で長いリャンカンができるので、内嵌張より強い。


\subsubsection*{嵌張を落とす場合}

\tehai{3456m 137899p 4578s}[東1局　西家　4巡目　ドラ \mj{2m}][kanchan-2]

手牌\ref{mj:kanchan-2}は単純に受け入れの枚数で言うと打 \mj{3m} や \mj{9p} の方が一番広いだが、嵌張以外の搭子は順子と対子、しかも頭ができやすい四連形があるのでピンフが見える。そうするとリーのみに避けてピンフにする方がいい。ここで打 \mj{1p} として嵌張を払う\citep[][p.~34]{uzaku-nankiru}。


%=================================================
\subsection{複合搭子} \label{sec:complex-shape}

複合搭子とは、単独搭子（辺張・嵌張・両面）と孤立牌が複合した形。3枚の牌を使用している複合搭子は表~\ref{tab:fkgtt-1} にあげられている。

\begin{table}[ht]
  \centering
  \begin{tabular}{lll}
    名称 & 形 & 受け入れ \\
    \hline
    辺張対子 & \mj{112m} & \mj{13m} 6枚 \\
    嵌張対子 & \mj{113m} & \mj{12m} 6枚 \\
    両嵌張 & \mj{135m} & \mj{24m} 8枚 \\
    両面対子 & \mj{223m} & \mj{124m} 10枚
  \end{tabular}
  \caption{複合搭子の種類}
  \label{tab:fkgtt-1}
\end{table}

リャンカン以外の複合搭子は対子含みなので、手牌に頭がないときはこれらの形は複合搭子として働かない。逆に対子含みの複合搭子が二つ以上あれば、対子が多すぎて手牌を弱くする。複合搭子同士の比較をする場合は、孤立牌同士や搭子同士を比較する場合と異なって全て外すのではなく、複合搭子から1枚外すことになる。よって、最終的な待ちが良くなりやすいように、強いほうから1枚外すのが原則である\citep{gendai-complex1, gendai-complex2}。対子の数、複合搭子の数と両面対子の数による選択は表~\ref{tab:fkgtt-2} の通り。

\begin{table*}[ht]
  \centering
  \begin{tabular}{llll}
    対子 & 複合搭子 & 両面対子 & 選択 \\
    \hline
    2 & 2 &   & 弱い複合搭子を対子に固定 \\
    3 & 2 & 0 & 内側の嵌張対子から対子をほぐし \\
      &   &   & 外側の辺張対子を対子に固定 \\
    3 & 2 & 1 & 両面固定 \\
    3 & 3 & 0 & 弱い複合搭子を対子に固定 \\
    3 & 3 & 1 & 両面固定 \\
    3 & 3 & 2 & 弱い複合搭子を対子に固定
  \end{tabular}
  \caption{対子の数、複合搭子の数と両面対子の数による選択}
  \label{tab:fkgtt-2}
\end{table*}


\subsubsection*{対子が二つの場合}

\textbf{対子が二つあったら単純搭子にしない}。対子のフォロー牌を働かせると縦引きのため、対子を二つのままで維持しないといけない。これは「2ヘッド理論」と呼ぶ\citep[][p.~28]{uzaku-haikouritsu}。例え\ref{mj:fkgtt-1}の \mj{233p466s} から \mj{3p} を外すと \mj{3p6s} の受け入れが失う上 \mj{4s} が遊び牌になってしまう。頭がないから \mj{5s} 引いても聴牌にならない。打 \mj{2p} と \mj{4s} にすると対子が二つ。両面対子の受け入れが嵌張対子より多いので打 \mj{4s}。受け入れは表~\ref{tab:fkgtt-1} の通り。強い複合搭子から切るではなく、弱い複合搭子を頭にするケース。

\tehai{12334578m 233p 466s}[][fkgtt-1]

\begin{table}[ht]
  \centering
  \begin{tabular}{lll}
    切る牌 & \mj{4s} & \mj{2p} \\
    \hline
    受け入れ & \mj{69m 134p 6s} 20枚 & \mj{69m 3p 56s} 16枚
  \end{tabular}
  \caption{手牌\ref{mj:fkgtt-1}の受け入れ}
  \label{tab:fkgtt-1}
\end{table}

\textbf{二つの両面対子があるなら、いい待ちを残すため、中に寄る方を頭にする}\citep[][pp.~208--209]{suphx}。手牌\ref{mj:fkgtt-1-1}は \mj{4m} と \mj{5s} の二択。\mj{25m} は \mj{47s} よりアガリやすいので \mj{5s} を切る方がいいい。

\tehai{334789m 34p 111566s}[][fkgtt-1-1]

手牌\ref{mj:fkgtt-1-2}も同じ \mj{778m556s} の二つの両面対子の中から中に寄る \mj{556s} を頭にする\citep{shibukawa-twopairs}。\mj{345556s} という複雑そうな形に手を出す抵抗感は多少あるけれど、\mj{345s} と \mj{556s} と認識すると通常の牌効率の通りです。表~\ref{tab:fkgtt-1-2} の通り、打 \mj{5s}。

\tehai{778m 12356p 345556s}[東一局　東家　6巡目　ドラ \mj{6p}][fkgtt-1-2]

\begin{table}[ht]
  \centering
  \begin{tabular}{llll}
    切る牌 & 受け入れ \\
    \hline
    \mj{7m} & \mj{69m47p} 16枚 \\
    \mj{8m} & \mj{7m47p2457s} 22枚 \\
    \mj{6s} & \mj{679m47p25s} 23枚
  \end{tabular}
  \caption{手牌\ref{mj:fkgtt-1-2}の受け入れ}
  \label{tab:fkgtt-1-2}
\end{table}


\subsubsection*{対子が三つ、複合搭子二つ}

手牌\ref{mj:fkgtt-2}に対子が多すぎるので二つにするとすると打 \mj{3p} と打 \mj{6s} の選択。両方の受け入れは同じ5種16枚だが、良形聴牌になる受け入れは表~\ref{tab:fkgtt-2} の通り打 \mj{3p} の方が多いので、強い複合搭子を単独搭子にするケース。

\tehai{12334588m 233p 466s}[][fkgtt-2]

\begin{table}[ht]
  \centering
  \begin{tabular}{lll}
    切る牌 & \mj{3p} & \mj{6s} \\
    \hline
    受け入れ & \mj{8m14p56s} 16枚 & \mj{8m134p5s} 16枚 \\
    両面聴牌ツモ & \mj{8m 56s} 8枚 & \mj{5s} 4枚 \\
  \end{tabular}
  \caption{手牌\ref{mj:fkgtt-2}の受け入れ}
  \label{tab:fkgtt-2}
\end{table}

愚形含み同士なら両面変化が多い方を単独搭子にする。手牌\ref{mj:fkgtt-3}は良形聴牌にならない。よって、\ref{mj:fkgtt-3}は打 \mj{6s} にしてツモ \mj{37s} で両面に渡るのを期待する。

\tehai{12334588m 133p 466s}[][fkgtt-3]

全ての複合搭子の両面変化が乏しくて単独の対子が両面になりやすい場合、一番弱い複合搭子を対子にして、中寄りの対子の両面変化を期待する。よって、\ref{mj:fkgtt-4}は打 \mj{8s} 。

\tehai{123345m 11366p 899s}[][fkgtt-4]

手牌\ref{mj:fkgtt-4-1}は \mj{11333445p} と \mj{778s} の二択。\mj{11333445p} は \mj{11p} + \mj{334p} + \mj{345p} と \mj{11p} + \mj{333p} + \mj{445p} と分析できるので、受け入れは \mj{123456p} となった。受け入れが広い形に手を出さないで、打 \mj{7s} による両面固定の方の受け入れが、表~\ref{tab:fkgtt-4-1} の通り、一番広い\citep[][p.~24]{uzaku-nankiru}。

\tehai{234m11333445p778s}[東1局　北家　7巡目　ドラ \mj{3p}][fkgtt-4-1]

\begin{table}[ht]
  \centering
  \begin{tabular}{llll}
    切る牌 & 受け入れ \\
    \hline
    \mj{7s} & \mj{123456p69s} 24枚 \\
    \mj{4p} & \mj{1236p679s} 21枚 \\
    \mj{3p} & \mj{125p679s} 19枚
  \end{tabular}
  \caption{手牌\ref{mj:fkgtt-4-1}の受け入れ}
  \label{tab:fkgtt-4-1}
\end{table}

手牌\ref{mj:fkgtt-4-2}も両面対子同士の比較なので \mj{5m36s} のどちらが切っても受け入れが同じだが、今回ドラが \mj{3s} なので、ツモドラが困らないように \mj{6s} を残して \mj{5m} を落とすべき。そしてドラ \mj{3s} 引いた時 \mj{22z} を落としてタンヤオへ移行できる\citep[][p.~24]{uzaku-nankiru}。

\tehai{556m 678p 345566s 22z}[東1局　東家　6巡目　ドラ \mj{3s}][fkgtt-4-2]

手牌\ref{mj:fkgtt-ssk}で打 \mj{8m} にと \mj{2s} に一見差がないように見えるが、\mj{7p} を引いて789三色のために \mj{78m} を固定したい。その場合、\mj{78m788p223s} から \mj{3s} を落として \mj{22s} を頭にする。ここで打 \mj{2s} にすると \mj{88m88p} のどちらが頭にならないといけないので、三色の変化が消える\citep{nemata-16}。

\tehai{788m 23488p 223789s}[ドラ \mj{4z}][fkgtt-ssk]

\textbf{対子を崩さない場合は役牌対子がある時}\citep[][p.~72]{uzaku-haikouritsu}。手牌\ref{mj:fkgtt-4-3}の \mj{77z} 対子を暗刻、明刻、または最終形と想定すれば、他の部分がもう一つの対子を求めている。頭固定の方が単独搭子固定より2枚損だが、\mj{77z} が生かせるなら悪くないでしょう。ピンズを引いた形を想定して、打 \mj{8p} の方が打 \mj{2s} より柔軟。よって、ここはこの前の結論を覆して打 \mj{8p}。

\tehai{789m 345668p 244s 77z}[][fkgtt-4-3]

なぜ同じ対子三つ、複合搭子二つのケースでも、対子の性質により打牌の基準が変わるか。それは「頭候補としての対子」と「面子候補としての対子」の差\multicitep{nemata-13; nemata-16}。\ref{mj:fkgtt-2}の \mj{88m} は変化が少ないので「頭候補としての対子」であり、他の複合搭子でバーランスをとるようとする結果は両面固定、そして \mj{4s} がついた \mj{66s} は「面子候補としての対子」。一方、\ref{mj:fkgtt-4}の \mj{66p} は中張牌、すなわち変化が見込む「面子候補としての対子」であり、そうすると他の部分に「頭候補としての対子」が求められる。手牌\ref{mj:fkgtt-4-3}の \mj{77z} も「面子候補としての対子」、だから他の部分を単純対子にしないといけない。


\subsubsection*{対子が三つ、複合搭子三つ（両面対子なし）}

\textbf{対子含みの複合搭子が三つあれば、両面変化が一番乏しい方を頭にする}。そうすると、残りの複合搭子が面子化する時には余剰牌は出ない。手牌\ref{mj:fkgtt-5}では両面にならない \mj{899s} から \mj{8s} を落とす。頭に固定する場合は辺張に固定する場合より2枚損だが、表~\ref{tab:fkgtt-5a} と表~\ref{tab:fkgtt-5b} の通り、\mj{7p} 引いた時受け入れが4枚得。

\tehai{12378m 113668p 899s}[][fkgtt-5]

\tehai{12378m 113678p 99s}[][fkgtt-5-1]

\tehai{12378m 113678p 89s}[][fkgtt-5-2]

\begin{table}[ht]
  \centering
  \begin{tabular}{lll}
    切る牌 & \mj{8s} & \mj{9s} \\
    \hline
    受け入れ & \mj{69m1267p9s} 22枚 & \mj{69m1267p7s} 24枚
  \end{tabular}
  \caption{手牌\ref{mj:fkgtt-5}の受け入れ}
  \label{tab:fkgtt-5a}
\end{table}

\begin{table}[ht]
  \centering
  \begin{tabular}{lll}
    切る牌 & \mj{8s} & \mj{9s} \\
    \hline
    手牌 & \ref{mj:fkgtt-5-1} & \ref{mj:fkgtt-5-2} \\
    受け入れ & \mj{69m12p9s} 16枚 & \mj{69m7s} 12枚
  \end{tabular}
  \caption{手牌\ref{mj:fkgtt-5}から牌を切って \mj{668p} が面子化する時の受け入れ}
  \label{tab:fkgtt-5b}
\end{table}


\subsubsection*{対子が三つ、複合搭子三つ（両面対子一つ）}

\textbf{両面対子が一つあれば、そこを両面固定する}。これは「強いほうから1枚外す」の原則を従っている。手牌\ref{mj:fkgtt-6}から \mj{26p9s} のどちらを切っても受け入れが同じ8種28枚だが、表~\ref{tab:fkgtt-6} の通り、打 \mj{2p} にすると全てのツモで16枚の一向聴になる。一方、打 \mj{6p} ツモ \mj{7p} で完全一向聴になるが、12枚の一向聴になる確率がもっと高い。よって、打 \mj{2p}。

\tehai{12378m 223668p 899s}[][fkgtt-6]

\begin{table}[ht]
  \centering
  \begin{tabular}{lll}
    切る牌 & 次のツモ & 一番広い受け入れの打牌 \\
    \hline
    \mj{2p} & \mj{69m14p} & 打 \mj{8p} $\to$ \mj{146p79s} 16枚 \\
            &          & 打 \mj{8s}   $\to$ \mj{1467p9s} 16枚 \\
            & \mj{6p}  & 打 \mj{8p8s} $\to$ \mj{69m14p} 16枚 \\
            & \mj{7p}  & 打 \mj{6p8s} $\to$ \mj{69m14p} 16枚 \\
            & \mj{7s}  & 打 \mj{8p9s} $\to$ \mj{69m14p} 16枚 \\
            & \mj{9s}  & 打 \mj{8p8s} $\to$ \mj{69m14p} 16枚 \\
    \hline
    \mj{6p} & \mj{69m} & 打 \mj{8s}   $\to$ \mj{1247p9s} 16枚 \\
            & \mj{14p} & 打 \mj{2p8s} $\to$ \mj{69m7p} 12枚 \\
            & \mj{2p}  & 打 \mj{3p8s} $\to$ \mj{69m7p} 12枚 \\
            & \mj{7p}  & 打 \mj{8s}   $\to$ \mj{69m124p9s} 20枚 \\
            & \mj{7s}  & 打 \mj{3p9s} $\to$ \mj{69m7p} 12枚 \\
            & \mj{9s}  & 打 \mj{3p8s} $\to$ \mj{69m7p} 12枚
  \end{tabular}
  \caption{手牌\ref{mj:fkgtt-6}から \mj{2p} と \mj{6p} を切った後のツモと受け入れ}
  \label{tab:fkgtt-6}
\end{table}

実戦でリアタイでこれらの計算をするのは人間にちょっと難しいので、「強いほうから1枚外す」の原則だけ覚えばよし。


\subsubsection*{対子が三つ、複合搭子三つ（両面対子二つ）}

\textbf{両面対子が二つあれば、残りの愚形含みの複合搭子を頭にする}。そうすると対子が減らない上、どちらの両面対子を単独両面にする選択に比べると2枚損になるが、一向聴のとき受け入れの広さを考えると前者の方が有利。表~\ref{tab:fkgtt-7} の通り、\ref{mj:fkgtt-7}で打 \mj{2p} とすると受け入れが9種32枚で、打 \mj{8s} より2枚得だが、ツモ \mj{568p79s} で複合搭子の方が先に面子になると余剰牌が出て、受け入れが両面の両端の16枚だけになる。一方、打 \mj{8s} とするとどちらの牌を引いても完全一向聴になる。よって、打 \mj{8s}。

\tehai{12378m223677p899s}[][fkgtt-7]

\begin{table}[ht]
  \centering
  \begin{tabular}{lll}
    切る牌 & 次のツモ & 一番広い受け入れの打牌 \\
    \hline
    \mj{2p} & \mj{69m14p} & 打 \mj{8s} $\to$ \mj{14578p9s} 20枚 \\  
            & \mj{58p} & 打 \mj{7p8s} $\to$ \mj{69m14p} 16枚 \\
            & \mj{7p9s} & 打 \mj{6p8s} $\to$ \mj{69m14p} 16枚 \\
            & \mj{7s} & 打 \mj{6p9s} $\to$ \mj{69m14p} 16枚 \\
    \hline
    \mj{8s} & \mj{69m} & 打 \mj{2p} $\to$ \mj{14578p9s} 20枚 \\
            &          & 打 \mj{7p} $\to$ \mj{12458p9s} 20枚 \\
            & \mj{14p} & 打 \mj{2p} $\to$ \mj{69m578p9s} 20枚 \\
            & \mj{2p} & 打 \mj{3p} $\to$ \mj{69m578p9s} 20枚 \\
            & \mj{58p} & 打 \mj{7p} $\to$ \mj{69m124p9s} 20枚 \\
            & \mj{7p} & 打 \mj{6p} $\to$ \mj{69m124p9s} 20枚 \\
            & \mj{9s} & 打 \mj{3p} $\to$ \mj{69m2578p} 20枚 \\
            &         & 打 \mj{6p} $\to$ \mj{69m1247p} 20枚
  \end{tabular}
  \caption{手牌\ref{mj:fkgtt-7}から \mj{2p} と \mj{8s} を切った後のツモと受け入れ}
  \label{tab:fkgtt-7}
\end{table}

手牌\ref{mj:fkgtt-8}は同じ複合搭子が三つ、両面対子が二つの牌姿だが、カン \mj{7s} リーのみとタンヤオの可能性を加味すると、\citet[p.~68]{uzaku-haikouritsu}が打 \mj{7p} の答えを出した。場況・ドラなどの情報が一切ないで判断すると、タンヤオを最優先するのは無難。そうすると、\mj{1p} が受け入れられなくて、\mj{223p} の受け入れ枚数が \mj{224p} と同じ6枚となり、前小節の結論に準じると打 \mj{7p} となる。

\tehai{577m 223677p 23468s}[][fkgtt-8]

両向聴の時、三つ目の対子が他の搭子と比べると受け入れが2枚損。対子を二つ持っていけば受け入れが最大となる。手牌\ref{mj:fkgtt-7}で両面固定によって最大の受け入れが32枚。しかし、\ref{mj:fkgtt-7}や\ref{mj:fkgtt-8}はともに対子にフォローがある牌姿で、両面固定するとどちらの対子が面子化可能性があり、一つの対子が面子化すれば他の対子のフォロー牌が不要になって、一向聴の時点の受け入れが狭くなる。


%--------------------------------------------
\subsubsection{両嵌張とその変化形}

\mj{135m} の形は「両嵌張」と呼ばれている。略として「リャンカン」。多くの場合はリャンカンを残して他の複合搭子を崩す。両面対子があればそれを両面にする。両面に変化しやすい嵌張対子があれば（例え \mj{335m} とか \mj{557m} とか）それを嵌張にする。なぜなら、リャンカンを嵌張にすると振聴にならない前提で両面になりにくい\citep[][p. 82]{uzaku-haikouritsu}。


\subsubsection*{両嵌張と他の搭子の選択}

両嵌張は3枚で構成する膨大なブロックで、14枚の手牌から両嵌張を嵌張にするか、もしくは他の搭子から1枚を外すか、というのは永遠の課題。両嵌張を嵌張にすると受け入れが4枚損になるので、打牌の基準はその4枚損のトレードオフ：

\begin{enumerate}
  \item 両嵌張を嵌張にしても受け入れが最大ならそうすべき
  \item 受け入れの差がない時
  \begin{enumerate}
    \item ピンフが見える時に好形率のため両嵌張をキープする
    \item ピンフさえないなら嵌張にして引っ掛けに頼る
  \end{enumerate}
\end{enumerate}

手牌\ref{mj:rk-12}は \mj{135m} と \mj{233345p} の二択。後者は \mj{233p} の両面対子を含んでいて受け入れが \mj{134p} と思われがち、実際 \mj{33345p} から見ると \mj{6p} も受けている。表~\ref{tab:rk-12} の通り、両嵌張のためにピンズに手を出すと受け入れが激減するのでそうすべきではない。嵌張が待ちになる時に出やすさを考えると打 \mj{5m} として引っ掛けを作る\citep[][pp.~44--45]{fkt}。

\tehai{135m 23334599p 789s}[ドラ \mj{9p}][rk-12]

\begin{table}[ht]
  \centering
  \begin{tabular}{lll}
    切る牌 & \mj{5m} & \mj{3p} \\
    \hline
    受け入れ & \mj{2m13469p} 18枚 & \mj{24m14p} 15枚
  \end{tabular}
  \caption{手牌\ref{mj:rk-12}の受け入れ}
  \label{tab:rk-12}
\end{table}

手牌\ref{mj:rk-12a}は\ref{mj:rk-12}と似ている牌姿だが、今回ピンズの部分は複合形ではないので、打 \mj{3p} としても受け入れも減らない。今回の牌姿はツモ \mj{24m} ならメンピンドラドラの7700になって、ツモ \mj{14p} で嵌張待ちになってしまう5200より高いので、好形率のために打 \mj{3p} のが正着\citep[][pp.~46--47]{fkt}。

\tehai{135m 23356799p 789s}[ドラ \mj{9p}][rk-12a]

手牌\ref{mj:rk-13}のようなピンフさえない手牌はどうあがいても5200に過ぎない。それならアガリ率のために打 \mj{5p} \citep[][pp.23--25]{ishibashi-black}\footnote{両嵌張は膨大なので、基本的に完全一向聴にならない。受け入れの差が出ない場合、アガリ率のために両嵌張に手を出すか、好形率のためにキープするか、実はピンフの有無に左右されている。「好形率」の価値は「好形で成立するピンフ」と繋がっているからです。よって、ピンフが見えるか否かを無視して手牌を混同してはいけない。\citet[pp.~46--47]{fkt}の場合はピンフが見えるケースだけに触る。\citet[pp.23--25]{ishibashi-black}の場合はピンフが見えないケースだけに触る。これは誤解を招きやすい。}。

\tehai{23399m 135p 406888s}[南家　4巡目　ドラ \mj{4s}][rk-13]

手牌\ref{mj:rk-1}にはリャンカンと暗刻がある。切る牌の選択の受け入れは表~\ref{tab:rk-1} に挙げられている。頭がないので、搭子が面子になるとき暗刻を頭にすると縦引きで搭子を頭にする二つの可能性があるので、その受け入れは \mj{8s} の頭を固定してリャンカンと両面の完成を期待するの受け入れより広い\citep[][p. 84]{uzaku-haikouritsu}。

\tehai{246m 23467p 45688 -8s}[][rk-1]

\begin{table}[ht]
  \centering
  \begin{tabular}{lll}
    切る牌 & 受け入れ & 枚数 \\
    \hline
    \mj{2m} & \mj{456m5678p} & 7種24枚 \\
    \mj{6m} & \mj{234m5678p} & 7種24枚 \\
    \mj{8s} & \mj{35m58p}    & 4種16枚
  \end{tabular}
  \caption{手牌\ref{mj:rk-1}のそれぞれの切る牌の受け入れ}
  \label{tab:rk-1}
\end{table}

手牌\ref{mj:rk-7}に両面両面の二度受けがあるため、表~\ref{tab:rk-7} の通り両面を落としても嵌張だらけの \mj{24668s} から \mj{268s} のどっちを落としても受け入れが同じ。最終形から考えるとまず両面を落とすことはなく、ツモ \mj{469m} でマンズの両面へ移行することを見込めば \mj{66s} の対子も残したいところ。そして678の三色も見込めば \mj{2s} が一番不要になる\citep[][p.~14]{uzaku-nankiru}。

\tehai{55678m 3467p 24668s}[東1局　東家　5巡目　ドラ \mj{7p}][rk-7]

\begin{table}[ht]
  \centering
  \begin{tabular}{lll}
    切る牌 & 受け入れ & 枚数 \\
    \hline
    \mj{34p} & \mj{5m58p3567s} & 7種24枚 \\
    \mj{67p} & \mj{5m25p3567s} & 7種24枚 \\
    \mj{2s} & \mj{5m258p567s}  & 7種24枚 \\
    \mj{6s} & \mj{258p357s}    & 6種24枚 \\
    \mj{8s} & \mj{5m258p356s}  & 7種24枚
  \end{tabular}
  \caption{手牌\ref{mj:rk-7}の受け入れ}
  \label{tab:rk-7}
\end{table}


\subsubsection*{13r5のリャンカン}

\mj{130p} のリャンカンの赤ドラを使い切れるツモは \mj{4567p} だけなので、\mj{1p} を切ってもこれらのツモによって赤ドラも使い切れるので \mj{1p} がいらなさそう。が、\mj{2p} が埋められたら \mj{0p} が浮いてツモ \mj{34567p} で赤が使えるし、\mj{1p} 引いたら頭＋嵌張という2ブロックができて、搭子が足りない時こそこの \mj{1p} は温存しよう\citep{yuusee-13r5}。

例え\ref{mj:rk-9}から孤立のダブ \mj{1z} を切るべき。実際\ref{mj:rk-9}の局で次に \mj{2p} を引いて \mj{0p} にくっつく形の\ref{mj:rk-9a}になる。そうなると \mj{1s} を切る。

\tehai{345688m 130p 1335s 1z}[東1局　1巡目　ドラ \mj{5z}][rk-9]

\tehai{345688m 130p 1335s -2p}[東1局　2巡目　ドラ \mj{5z}][rk-9a]

似ている牌姿だが、\ref{mj:rk-10}にはもうブロックが整っているので \mj{1p} 切ってタンヤオで行こう。

\tehai{23445688m 130p 335s}[東1局　2巡目　ドラ \mj{5z}][rk-10]

そして\ref{mj:rk-11}には \mj{10m} というリャンカンの先導形がある。ここは \mj{1m} よりオタ風を残すのは損。ただし \mj{46z} と比べると自風・役牌を残すの方がいい。

\tehai{10m 1389p 3666s 246z}[東2局　北家　ドラ \mj{4m}][rk-11]


\subsubsection*{縦引きしたリャンカン（2246、2466など）}

\mj{1135m} や \mj{2246m} の形は5(6)引いたら強い形\footnote{\mjfn{11355m} から1面子＋1頭を作る受け入れは \mjfn{1245m} 4種12枚。\citet{yuusee-11355}によるとこの形は「リャンカンシャンポン」と称される。}になるので離れた対子＋嵌張より強い\citep{chirno-kanchan}。

\mj{1355m} や \mj{2466m} の形はリャンカンシャンポンになる可能性がある上、ツモ \mj{46m} で \mj{134556m} という両面嵌張ができるので \mj{1135m} より強い\citep{chirno-kanchan}。

手牌\ref{mj:rk-2}は打 \mj{4p} でリャンカンを固定するか打 \mj{8p} で頭を二つにするかの選択の比較は表~\ref{tab:rk-2} に挙げられている。単にピンフになる確率で選択すれば、打 \mj{4p} は有利。タンヤオがないため \mj{9m} や \mj{4p} が暗刻になっても嬉しくない。完全一向聴になるツモ牌も \mj{3p} だけ\citep[][p. 86]{uzaku-haikouritsu}。

\tehai{12399m 4468p 3478 -9s}[][rk-2]

\begin{table}[ht]
  \centering
  \begin{tabular}{lll}
    切る牌 & \mj{4p} & \mj{8p} \\
    \hline
    受け入れ & \mj{57p25s} & \mj{9m45p25s} \\
    枚数 & 4種16枚 & 5種16枚 \\
    ピンフがつくツモ & \mj{57p} (50\%) & \mj{5p} (25\%)
  \end{tabular}
  \caption{手牌\ref{mj:rk-2}のそれぞれの切る牌の受け入れ}
  \label{tab:rk-2}
\end{table}

手牌\ref{mj:rk-3}は\ref{mj:rk-2}と違い、その一つの頭が \mj{7m} となり、完全一向聴になるツモ牌（\mj{68m3p}）が増えている。よって、ここは \mj{8p} を切る\citep[][p. 88]{uzaku-haikouritsu}。

\tehai{12377m 4468p 3478 -9s}[][rk-3]

手牌\ref{mj:3577-1}には役がないので立直に向かう。打 \mj{3s} の受け入れが打 \mj{7s} と同じだが \mj{9m7s} が暗刻になっても別に嬉しくないのでここで赤 \mj{0s} を残し \mj{7s} を切る\citep[][p.~26]{uzaku-nankiru}。

\tehai{99m 23468p 3077s 444z}[東1局　東家　7巡目　ドラ \mj{2p}][3577-1]

手牌\ref{mj:3577-2}と\ref{mj:3577-1}の違いは \mj{77z} は鳴ける。\mj{7z} を鳴いた後頭が必要なのでここで \mj{3s} を切る。赤 \mj{5s} が使い切れないだが速度優先なら仕方がない\citep[][p.~26]{uzaku-nankiru}。

\tehai{666789m 34p 3077s 77z}[東1局　西家　6巡目　ドラ \mj{2m}][3577-2]

手牌\ref{mj:3577-3}は \mj{468s} 両嵌張にすると \mj{446s} 嵌張対子にするの二択。対子が \mj{55m77p} あって、両向聴の時点で打 \mj{8s} は打 \mj{4s} より2枚損。一向聴の時、両面＋両嵌張や両面＋嵌張＋双碰なので枚数の差がでない。ピンフの枚数と両向聴の受け入れを見ると打 \mj{4s} が最善となる\citep{zundamon-3head}。両向聴の時、三つ目の対子が他の搭子と比べると受け入れが2枚損。対子を二つ持っていけば受け入れが最大となる

\tehai{23577m 12377p 4468s}[東1局　西家　6巡目　ドラ \mj{1p}][3577-3]


\subsubsection*{中ぶくれリャンカン（1335、2446など）}

1335と2446の形は3種10枚（1335に234、2446に345）の数牌を引いて1面子ができるが、孤立の37と比べると先に3枚の複合搭子にするべき\citep{chirno-kanchan}。一方、\mj{3557m} の形は変化が多いので外嵌張より持っててもいい。中ぶくれリャンカンから数牌を一つ落とすコツは下記のよう：

\begin{itemize}
  \item 両向聴
  \begin{itemize}
    \item 端の方を落として両面変化を狙う
  \end{itemize}
  \item 一向聴
  \begin{itemize}
    \item ピンフが見えるならリャンカンにする
    \item ピンフが見えないなら中の方を切る
  \end{itemize}
\end{itemize}

\begin{table}[ht]
  \centering
  \begin{tabular}{llll}
    切る牌 & \mj{1p} & \mj{3p} & \mj{5p} \\
    \hline
    残った形 & \mj{335p} & \mj{135p} & \mj{133p} \\
    ピンフツモ & \mj{4p} 4枚 & \mj{24p} 8枚 & \mj{2p} 4枚 \\
    両面変化ツモ & \mj{68m6p} 12枚 & （なし） & \mj{68m} 8枚 \\
    最終形 & \mj{3p} 双碰 & \mj{2p} 嵌張 & \minibox{\mj{2p} 嵌張\\\mj{3p} 双碰}
  \end{tabular}
  \caption{\mj{77m1335p}の形から1枚を外す選択肢の比較}
  \label{tab:rk-1335}
\end{table}

表~\ref{tab:rk-1335} の通り、リャンカンにするとピンフができるツモ枚数が一番多い。他の部分が先に埋まっても \mj{5p} 切りリーチでも筋待ちができる。よって\ref{mj:rk-1335a}は \mj{3p} を切る方がいい。一方、手牌\ref{mj:rk-1335b}や\ref{mj:rk-1335c}のようにピンフがない時、最終形カン \mj{2p} のために \mj{5p} を先に切る方がいい。特に\ref{mj:rk-1335b}に \mj{7m3p} のポンテンもある\citep{yuusee-1335}。

\tehai{12377m 1335p 45789s}[東1局　東家　5巡目　ドラ \mj{9s}][rk-1335a]

\tehai{77m 1335p 45789s 555z}[東1局　東家　5巡目　ドラ \mj{9s}][rk-1335b]

\tehai{77m 1335p 45789s 333z}[東1局　東家　5巡目　ドラ \mj{9s}][rk-1335c]

\mj{1p} 切って両面変化が期待できそうだが、一向聴の時は両面変化より待ちと打点を優先すべき。


\subsubsection*{横引きしたリャンカン（2357など）}

\textbf{基本的にはリャンカンにしない}\citep{yuusee-2357}。タンヤオが見えそうでも。同じ8枚の受け入れだが、リャンカンが両面よりもう1枚消費しているので守備力には損。最終形としても嵌張になりやすいので劣っている。\textbf{基本的に5から切るが、頭が欲しい場面は7を切る}。

5から切る原因の一つは、内嵌張同士の両面変化を比較すると、2357で4を引いたら実際両面二つにならない。23のフォローがあるので、できたのは1面子＋1嵌張。ですので、2357の57を両面に変化するツモは8だけ。例として、手牌\ref{mj:rk-2357a}は \mj{57m} と \mj{46p} の嵌張があって、嵌張を両面にするツモは \mj{8m37p}。\mj{5m} 切って \mj{4m} 引いても \mj{23m} に受ける。タンヤオが見えるが、ここで \mj{5m} を切る。

\tehai{2357m 46p 22256788s}[東1局　3巡目　ドラ \mj{4z}][rk-2357a]

両面変化できるツモが1種だけなので、実際に外嵌張にも劣っている。例え手牌\ref{mj:rk-2357b}には \mj{2357m} と \mj{79s} がある。\mj{57m} の部分を落としてよりいい最終形を目指す。

\tehai{2357m 2225699p 79s}[東1局　3巡目　ドラ \mj{4z}][rk-2357b]

頭がない時は7から切る。235の形は2と5引きで1頭＋1搭子ができ、4引いて四連形になって、残り6枚の25で頭ができるのだ。よって、手牌\ref{mj:rk-2357c}は \mj{7m} を切る。

\tehai{2357m 234678p 3456s}[東1局　3巡目　ドラ \mj{4z}][rk-2357c]


\subsubsection*{両面がついたリャンカン（34468）}

34468の形で2457引いて1面子＋1搭子ができ、強い形と思われる\citep{clearrain-34468}。5引いても両面嵌張\footnote{344568という形。詳しくは~\ref{subsec:rmkc} 小節を参照。}ができるので、その時ツモ257で面子が二つできる\citep{yuusee-34468}。浮き牌がある時その浮き牌を先に切る。四連形や中ぶくれという強そうな形があっても。よって、手牌\ref{mj:rk-34468a}は \mj{3p} を切る。

\tehai{34468m 3456p 11789s}[][rk-34468a]

6ブロックの時にも基本的に残す。34468を崩して8を切る場面は他の部分に愚形搭子がない時。手牌\ref{mj:rk-34468b}にはもっと弱い \mj{79s} 嵌張があるのでそこを先に崩す。一方、\ref{mj:rk-34468c}にもっと弱い搭子がないので、\mj{8m} を落とす方がいい。

\tehai{34468m 45789p 2279s}[][rk-34468b]

\tehai{34468m 45789p 2399s}[][rk-34468c]

5ブロックの時、対子があれば34468の方を残す。対子がない時、また巡目が早いなら8を切って向聴数を戻して、34両面を中心に手を進む。対子がない、そして巡目も遅いなら、聴牌に向かって3を切って44の部分を対子にする。よって、同じ手牌\ref{mj:rk-34468d}だが、巡目・ドラ・場況などによって \mj{3m} か \mj{8m} かを切る選択も違う。

\tehai{34468m 123789p 357s}[][rk-34468d]

一向聴の時に \mj{13556m} から \mj{5m} の先切りはオススメしない\citep{yuusee-13556}。手牌\ref{mj:rk-34468e}から \mj{5m} を切ったら両面聴牌になるツモは \mj{2m} の4枚だけ、\mj{47m} の8枚を引いたら愚形聴牌になる。一方、安牌の \mj{4z} を切ったらツモ \mj{24m} の8枚で両面聴牌、\mj{57m3p} の8枚で愚形聴牌になる。つまり打 \mj{4z} の両面聴牌ができる枚数が4枚得。選択肢の比較は表~\ref{tab:rk-34468e} を参照。

\tehai{13556m 33789p 123s 4z}[東1局　西家　6巡目　ドラ \mj{8p}][rk-34468e]

\begin{table}[ht]
  \centering
  \begin{tabular}{lll}
    切る牌 & \mj{5m} & \mj{4z} \\
    \hline
    残った形 & \mj{1356m33p4z} & \mj{13556m33p} \\
    両面聴牌 & \mj{2m} 4枚 & \mj{24m} 8枚 \\
    愚形聴牌 & \mj{47m} 8枚 & \mj{57m3p} 8枚
  \end{tabular}
  \caption{\ref{mj:rk-34468e}の\mj{13556m33p4z}の形から1枚を外す選択肢の比較}
  \label{tab:rk-34468e}
\end{table}

だが、ツモ \mj{6m} ならこのリャンカンを崩す方がいい。\mj{135566m} から \mj{13m} を外して \mj{5566m} が二度受けになるが愚形聴牌になるツモ牌は \mj{56m3p} の3種類6枚になる。つまり2枚減っている。良形聴牌の枚数は変わらない。そしてその余った牌は安牌として働くこともできる。手牌\ref{mj:rk-34468f}は\ref{mj:rk-34468e}に似ているが打牌が違って、ここは打 \mj{3m} にする。選択肢の比較は表~\ref{tab:rk-34468f} を参照。

\tehai{13556m 33789p 123s -6m}[東1局　西家　6巡目　ドラ \mj{8p}][rk-34468f]

\begin{table}[ht]
  \centering
  \begin{tabular}{lll}
    切る牌 & \mj{3m} & \mj{6m} \\
    \hline
    残った形 & \mj{15566m33p} & \mj{13556m33p} \\
    両面聴牌 & \mj{47m} 8枚 & \mj{24m} 8枚 \\
    愚形聴牌 & \mj{56m3p} 6枚 & \mj{57m3p} 8枚
  \end{tabular}
  \caption{\ref{mj:rk-34468f}の \mj{135566m33p} の形から1枚を外す選択肢の比較}
  \label{tab:rk-34468f}
\end{table}

このように並び対子の形にするケースは34468の両面＋両嵌張にも適用できる\citep{yuusee-34468}。手牌\ref{mj:rk-34468fa}と\ref{mj:rk-34468fb}も \mj{6s} 切りがオススメ。

\tehai{444m 33p 678p 334468s}[東1局　西家　4巡目　ドラ \mj{9p}][rk-34468fa]

\tehai{11 789p 334468s 444z}[東1局　西家　4巡目　ドラ \mj{9p}][rk-34468fb]

もう一つ \mj{13556m} を一向聴で崩す場面は3\~7の対子があってツモ \mj{1m} の時。\mj{3m} を \mj{11m} 対子のフォロー牌として \mj{556m} を両面にすると、その瞬間の受け入れは \mj{13556m} と同じ枚数だが、\mj{33p} 対子の両面変化は \mj{135m} より多いので打 \mj{5m} の方がちょっと良さそう。よって、手牌\ref{mj:rk-34468g}は打 \mj{5m}。選択肢の比較は表~\ref{tab:rk-34468g} を参照。

\tehai{13556m 33789p 123s -1m}[東1局　西家　6巡目　ドラ \mj{8p}][rk-34468g]

\begin{table}[ht]
  \centering
  \begin{tabular}{lll}
    切る牌 & \mj{5m} & \mj{1m} \\
    \hline
    残った形 & \mj{11356m33p} & \mj{13556m33p} \\
    両面聴牌 & \mj{12m3p} 8枚 & \mj{24m} 8枚 \\
    愚形聴牌 & \mj{47m} 8枚 & \mj{47m} 8枚 \\
    両面変化 & \mj{24p} 8枚 & \mj{6m} 3枚
  \end{tabular}
  \caption{\ref{mj:rk-34468g}の \mj{113556m33p} の形から1枚を外す選択肢の比較}
  \label{tab:rk-34468g}
\end{table}

3\~7の対子がなくてもツモ1で打5が効く場面は役牌暗刻があるとき。ポン材は優秀だから。手牌\ref{mj:rk-34468h}も打 \mj{5m}。

\tehai{13556m 11789p 555z -1m}[東1局　西家　6巡目　ドラ \mj{8p}][rk-34468h]


\subsubsection*{離れリャンカン（245679）}

「離れリャンカン」という形は\ref{mj:rk-4}のピンズの部分のような \mj{38p} が受け入れになる形。\mj{2p} や \mj{9p} を切ると二つの嵌張のどちらが固定されるが、表~\ref{tab:rk-4} によるとその受け入れの差はない。打 \mj{6m} や打 \mj{2p} にすると三色の可能性が残る。よって、ピンフと三色の確率によりここで打 \mj{6m} にする\citep[][p. 89]{uzaku-haikouritsu}。

\tehai{566m 245679p 45699s}[][rk-4]

\begin{table}[ht]
  \centering
  \begin{tabular}{llll}
    切る牌 & \mj{6m} & \mj{2p} & \mj{9p} \\
    \hline
    受け入れ & \mj{47m38p} & \mj{467m8p9s} & \mj{467m3p9s} \\
    枚数 & 4種16枚 & 5種16枚 & 5種16枚 \\
    ピンフツモ & \mj{38p} (50\%) & \mj{8p} (25\%) & \mj{3p} (25\%) \\
    三色ツモ & \mj{4m8p} (50\%)\tablefootnote{\mjfn{8p} ツモれば \mjfn{47m} の待ちで50\%の確率で三色がつく} & \mj{4m8p} (50\%) & (0\%)
  \end{tabular}
  \caption{手牌\ref{mj:rk-4}のそれぞれの切る牌の受け入れ}
  \label{tab:rk-4}
\end{table}


\subsubsection*{二つのリャンカン}

手牌\ref{mj:rk-5}にはリャンカンが二つで、リャンカンにあと一つの牌がくっついている。\textbf{ここで打} \mj{5p} \textbf{が正解}\citep{hiii}。\ref{tab:rk-5}の通り打 \mj{3p} と打 \mj{5p} の受け入れが同じ12枚だが、\mj{357p} のリャンカンを固定すればツモ \mj{28p} で両面ができるのでここで嵌張対子にしない。打 \mj{2s} と \mj{4s} が余剰牌になるので要注意。

\tehai{123789m 3557p 2466s}[][rk-5]

\begin{table}[ht]
  \centering
  \begin{tabular}{llll}
    切る牌 & \mj{3p} & \mj{5p} & \mj{2s} \\
    \hline
    受け入れ & \mj{56p36s} 12枚 & \mj{46p3s} 12枚 & \mj{46p} 8枚 \\
  \end{tabular}
  \caption{手牌\ref{mj:rk-5}のそれぞれの切る牌の受け入れ}
  \label{tab:rk-5}
\end{table}

\textbf{二つのリャンカンのどちらを崩さなければならないといい待ちの方を残す}\citep[][p.~26]{feng-intermediate}。手牌\ref{mj:rk-6}にもリャンカンが二つで、どちらを嵌張にしなければならない。ここでは \mj{5p} を切ってモロ引っ掛けを作る。

\tehai{246999m 33567s 579p}[][rk-6]


\subsubsection*{リャンカンのそばに暗刻}

手牌\ref{mj:rk-8}から \mj{1s} を切って \mj{0s} が \mj{307s} のリャンカンに使えるだと思いがちだが \mj{3s} を引いたら \mj{0s} が出ないと聴牌が取れない。そう見ると打 \mj{1s} としても \mj{3s} の受け入れはなしと同じ、打 \mj{8s} とは実際受け入れが同じ枚数。よってここで \mj{68p} の嵌張を落とす。

\tehai{23468p 130777s 444z}[東1局　東家　7巡目][rk-8]

\begin{table}[ht]
  \centering
  \begin{tabular}{lll}
    切る牌 & 打 \mj{1s} & 打 \mj{8p} \\
    \hline
    受け入れ & \mj{678p3456s} 24枚 & \mj{6p12456s} 21枚
  \end{tabular}
  \caption{手牌\ref{mj:rk-8}の受け入れ}
  \label{tab:rk-8}
\end{table}

%--------------------------------------------
\subsubsection{134の形} \label{sec:134}

一見に \mj{1m} を切っても \mj{25m} の受け入れが減らないが、実際に残りの形によりこの1が大きな役をたつことができる\citep{yoteru-134}。\mj{134m} の形の有効牌は \mj{1m} から \mj{6m} まで、表~\ref{tab:134} の通りツモ牌によっていい形があり悪い形がある。要すると、

\begin{itemize}
\item くっつきしやすい
\item 頭ができやすい
\end{itemize}

\begin{table}[ht]
  \centering
  \begin{tabular}{llp{4cm}}
    ツモ牌 & 形 & 解説 \\
    \hline
    \mj{1m} & \mj{1134m} & 1雀頭＋1両面、つまり2ブロックというかなりいい形。 \\
    \mj{2m} & \mj{1234m} & ノベタンで頭ができやすいので、他のところに頭がない時残したい。 \\
    \mj{3m} & \mj{1334m} & 嵌 \mj{2m} と \mj{25m} 両面の二度受け、いい形とはいえない。 \\
    \mj{4m} & \mj{1344m} & 1雀頭＋1嵌張、待ちとしてはいい。 \\
    \mj{5m} & \mj{1345m} & 1面子だけに見えるが、この孤立の \mj{1m} は実際にとても優秀。ツモ \mj{2m} で両面ができ、ツモ \mj{3m} で嵌張ができ、ツモ \mj{6m} で \mj{13456m} となり両面に変化しやすい。 \\
    \mj{6m} & \mj{1346m} & 嵌張二つ。いい形とはいえない。
  \end{tabular}
  \caption{\mj{134m} の変化}
  \label{tab:134}
\end{table}

よって、5ブロック揃ってない時に \mj{134m} の \mj{1m} は役牌より残す。揃っている時でも頭がなければ \mj{1m} を残す。他の場合に \mj{1m} を切ってもいい。手牌\ref{mj:134-1}はまだブロックが揃っていないため \mj{56z} から切る。一方、\ref{mj:134-2}は既に揃っているので守備力のため \mj{1m} から落とす。手牌\ref{mj:134-3}もブロックが揃っているが頭がないため、安牌の \mj{2z} より \mj{1m} を残したい。

\tehai{134m 340899p 248s 56z}[東1局　西家][134-1]

\tehai{13488m 3468p 345s 56z}[東1局　西家][134-2]

\tehai{134m 1236789p 888s 2z}[東1局　西家　ドラ \mj{8s}][134-3]


\subsubsection*{手牌がバラバラすぎるなら受け入れより守備力}

\tehai{1348m 4899p 136s 567z}[][134-4]

手牌\ref{mj:134-4}のような門前が厳しい手なら \mj{134m} でどんないい形を作っても大抵他家に追わない。ここは字牌を残して \mj{1m} を落とす。


\subsubsection*{愚形リーのみは打たない}

\tehai{134m 1236789p 888s 2z}[東1局　西家　ドラ \mj{7z}][134-3a]

手牌\ref{mj:134-3a}は\ref{mj:134-3}と同じ牌姿だがドラが無くなった。リーのみなら少なくとも良形にしたいので、ここは \mj{1m} を捨てて \mj{25m} に固定する。


\subsubsection*{孤立の1がドラの時}

\tehai{?67m 134p 11?67s -7'77z}[東1局　西家　ドラ \mj{1p}][134]

\begin{table}[ht]
  \centering
  \begin{tabular}{lll}
    \mj{?} & \mj{1p}切り & \mj{4p}切り \\
    \hline
    \mj{5m} と \mj{5s} & 1000点7枚、2000点1枚 & 2000点 \\
    \mj{5m} と \mj{0s} & 2000点7枚、3900点1枚 & 3900点 \\
    \mj{0m} と \mj{0s} & 3900点7枚、7700点1枚 & 7700点
  \end{tabular}
  \caption{赤ドラの有無で\ref{mj:134}の価値}
  \label{tab:134-dora}
\end{table}

表~\ref{tab:134-dora} の通り、赤ドラがないとき、\mj{2p} の嵌張待ちは \mj{25p} の両面待ちより打点上昇の程度は大きくないため、\mj{1p} 切りはおすすめである。しかし、赤ドラまたは他のドラが1枚以上あれば、\mj{1p} を残し、\mj{1s} ポンで \mj{1p} 単騎に移行することや \mj{1p} を引いて双碰待ちことによる大きい打点上昇は期待できるので\textbf{序盤}では \mj{4p} を切る方がいい\citep{yuusee-134}。

%--------------------------------------------
\subsubsection{578の5が赤ドラ}

\begin{table}[ht]
  \centering
  \begin{tabular}{ll}
    聴牌まで遠い & \mj{078m}を残す \\
    一向聴 & タンヤオがないとき\mj{0m}切り \\
    聴牌 & 基本的に\mj{0m}切り \\
  \end{tabular}
  \caption{\mj{078m} の捌き方}
\end{table}

\tehai{078m 23557p 22478s -7s}[][5-1]

\tehai{078m 23557p 22478s -1p}[][5-2]

\tehai{078m 2456p 688s 555z -2p}[][5-3]

\tehai{078m 22456p 22256s -7s}[][5-4]

\tehai{078m 22567p 22256s -7s}[][5-5]

聴牌まで遠いとき、\mj{078m} を \mj{0m}のくっつきと \mj{78m} の両面と見直すと、\mj{34567m} のどちらを引いて2ブロックができる。ですので、\ref{mj:5-1}の場合 \mj{0m} と \mj{8m} のどちらを切るより \mj{78s} の両面を固定してツモギリしよう。\ref{mj:5-2}のような \mj{1p} を引いてしまうとき立直のみを避けるようになるべく \mj{0m} を使って \mj{4s} を切ろう\citep{yuusee-578}。

一向聴をとるとき、他の役があれば \mj{0m} を切って受け入れを最大限にするべきである。例外として、タンヤオや三色を見込むとき、\mj{8m} を切って仕掛けしても早く聴牌を目指そう。なので、\ref{mj:5-3}のように役牌がもうすでにある場合は \mj{0m} を切ろう。

聴牌を取るとき、基本的には \mj{0m} を切って両面待ちで立直する。なので、\ref{mj:5-4}の場合 \mj{69m} の待ちにして \mj{0m} を切るのがおすすめ。ただし、\ref{mj:5-5}のように567の三色があれば \mj{8m} を切って闇聴しよう。

%--------------------------------------------
\subsubsection{3568の形}

3568という形は頭ができやすい\citep{yuusee-3568}。両面以外の38の縦引きだけではなく、ツモ47で四連形ができるなら尚更。あとで3は孤立牌としても強い。よって、基本的には2ブロックとして残す。切るなら8から切る。手牌の他の部分と比較して3568の捌く基準は表~\ref{tab:3568} の通り。

\begin{table*}
  \centering
  \begin{tabular}{lll}
    手牌のブロック数 & 比較同士 & 3568の捌き \\
    \hline\hline
    4 & & 残す \\
    \hline
    5 & 浮き牌（頭なし） & 残す \\
      & 3456以外の浮き牌 & 残す \\
      & 辺張（唯一の么九牌） & 残す \\
      & 辺張（唯一の么九牌でない） & 崩す \\
      & 3456の浮き牌 & 崩す \\
    \hline
    6 & & 崩す \\
    \hline\hline
  \end{tabular}
  \caption{3568（2457）の形の捌き\citep{yuusee-3568, clearrain-3568}}
  \label{tab:3568}
\end{table*}

手牌\ref{mj:3568-1}は白の暗刻があって、辺張という一のブロック自体の価値が上がる。\mj{3p} から両面ができないうちに辺張は落とさない。よって、ここは打 \mj{8p}。

\tehai{12m 3568p 22788s 555z}[東1局　西家　3巡目　ドラ \mj{1s}][3568-1]

だが、辺張を外してタンヤオがつくなら辺張外しは推奨。\mj{8p} の縦引きや \mj{3p} のくっつきも十分なので。手牌\ref{mj:3568-2}は \mj{1m} を切る。

\tehai{12m 3568p 22455788s}[東1局　西家　3巡目　ドラ \mj{1s}][3568-2]

基本的に3568は3--7の先に切るが、頭がない時は3568の形を優先する。手牌\ref{mj:3568-3}は浮き牌の \mj{7m} と中ぶくれの \mj{6778s} があるが頭がない。よって、ここは \mj{7m} を落とす。

\tehai{237m 3568p 1236778s}[東1局　西家　3巡目　ドラ \mj{1s}][3568-3]

手牌\ref{mj:3568-4}には6ブロックがあって、打牌の選択候補は \mj{2m} と \mj{89p}。\mj{2457m} を嵌張二つにすると「辺張＞嵌張」という通常の落とす順に従って打 \mj{89p} の結論に辿り着く。しかし、\mj{89p} の欠点は下記の通り：

\begin{itemize}
\item \mj{45m} が両面に働きなくなる（\mj{7m} 引かない限り）
\item \mj{7p} の受け入れがなくなる
\end{itemize}

一方、\mj{457m89p} の形を残したら、\mj{3m} 引いたら \mj{45m} あるのでロスにならないし、ツモ \mj{7p} で面子ができるし、そのあと両面が \mj{45m67s} 二つある。強いていうと「辺張＜嵌張」の例。

\tehai{2457m 1189p 667s 444z}[][3568-4]

%--------------------------------------------
\subsubsection{35689の形}

\textbf{基本的には3を切る}。ダイレクトの7が偉いから。浮き牌としてこの3の強さ孤立の3に劣っているのは下記の通り、

\begin{itemize}
\item ツモ2で両面ができるが二度受けになる
\item ツモ4で両面できず34568(9)という面子＋嵌張になる
\end{itemize}

3を切ったらさらに辺張を払う両向聴に戻る場面は下記の通り\citep{yuusee-35689}：

\begin{enumerate}
\item 9切ったらタンヤオがつく
\item 強いの形（四連形など）がある
\end{enumerate}

両向聴の時で浮き牌同士を比較するならやはり35689の3が先に出ていく。手牌\ref{mj:35689-1}は打 \mj{3s}。

\tehai{22378m 1136p 35689s}[東1局　西家　5巡目　ドラ \mj{8m}][35689-1]

手牌\ref{mj:35689-2}は門前一向聴。前述した点でここは打 \mj{3s}。二度受けが嫌い気持ちがわかるが、ツモ \mj{7s} による良形聴牌を逃がさないように大人しく \mj{3s} を手放そう。

\tehai{234789m 115p 35689s}[東1局　西家　5巡目　ドラ \mj{8m}][35689-2]

逆に、\ref{mj:35689-3}でツモった \mj{6m} は先の \mj{3s} より強いのでここで辺張を払う価値がある。

\tehai{2346789m 115p 5689s}[東1局　西家　6巡目　ドラ \mj{8m}][35689-3]

89の辺張を払うとタンヤオがつくなら両向聴に戻っても構わない。手牌\ref{mj:35689-4}は打 \mj{9s}。

\tehai{22367m 2247p 35689s}[東1局　西家　5巡目　ドラ \mj{8m}][35689-4]

%--------------------------------------------
\subsubsection{三嵌張・四嵌張}

要すると「\textbf{偶数三嵌張＜孤立牌＜奇数三嵌張＜外嵌張＜奇数四嵌張＜内嵌張}」\citep{chirno-kanchan}。

奇数三嵌張（\mj{1357m} や \mj{3579m}）は \mj{8m}（\mj{2m}）引いたら両嵌張＋両面なので序盤で頭がない時 \mj{1m}〜\mj{9m} 全て有効牌なので孤立牌より残す。外嵌張があるなら両嵌張に固定する。

奇数四嵌張（\mj{13579m}）は \mj{2468m} 引いたら1面子＋1搭子なので外嵌張より残す。だが両面変化がないので内嵌張あるなら四嵌張を外す。

偶数三嵌張（\mj{2468m}）は1ブロックとしては膨大なスペースに対して効率が低下すぎるので普段は2や8を切る方がいいが、タンヤオを目指すとき24と68のリャンカンと見える、また縦重なりも期待できる。従って、手牌\ref{mj:3-1}から \mj{8s} を切り僅かな三色の可能性を残し、手牌\ref{mj:3-2}から \mj{11m} の対子を落とす\citep{yuusee-2468}。

\tehai{24m 44699p 2468s 77-7z}[][3-1]

\tehai{11357m 4668p 2468s -8p}[][3-2]

%--------------------------------------------
\subsubsection{二度受け}

\subsubsection*{並び対子}

\mj{3344p} のような形は、\mj{33p} と \mj{44p} の二つの対子と、\mj{34p} の順子が二重になる形の両方に解釈できる。ピンフに向かうなら、「愚形＞並び対子＞順子・他の対子」の切る順で行くのがいい。

手牌\ref{mj:narabtt-1}は \mj{34p} のどちらを切って完全ー向聴をとる。ソーズの部分に手を出さない\citep[][pp. 112--113]{uzaku-haikouritsu}。

\tehai{33678m 3344p 23478s}[][narabtt-1]

手牌\ref{mj:narabtt-2}はソーズの部分が嵌張 \mj{68s} になる。打 \mj{34p} の方は2枚得が、表~\ref{tab:narabtt-2} の通り好形聴牌から見ると実は4枚損。ツモ \mj{5s} を見て打 \mj{8s}。

\tehai{33678m 3344p 23468s}[][narabtt-2]

\begin{table}[ht]
  \centering
  \begin{tabular}{lll}
            & 打 \mj{3p} & 打 \mj{8s} \\ 
    \hline
    受け入れ & \mj{3m245p7s} 16枚 & \mj{3m2345p} 14枚 \\
    好形聴牌 & \mj{7s} 4枚 & \mj{25p} 8枚
  \end{tabular}
  \caption{手牌\ref{mj:narabtt-2}の受け入れ}
\end{table}

手牌\ref{mj:narabtt-3}は並び対子が端に寄るので打 \mj{68s} でも打 \mj{89p} でも好形聴牌の枚数が一緒。打 \mj{8p} はツモ \mj{7p} からできる亜両面を見る。今回の手牌は打 \mj{9p} とするとタンヤオがつくので打 \mj{8p} にしない。

\tehai{33678m 8899p 23468s}[][narabtt-3]


\subsubsection*{飛び対子}

\mj{7799p} のような対子が一つ離れた形は「飛び対子」と呼ばれている。順子に解釈できないため、多くの場合では3枚構成の1ブロックに構えるべき。それでも、概ね切る順は並び対子と同じ「愚形＞飛び対子＞順子・他の対子」。飛び対子のどちらから落とすのは牌の数字を考える\citep{yoteru-tobitoitsu}：

\begin{itemize}
\item \mj{1133p}：出アガリを狙うなら打 \mj{3p} として \mj{1p} の双碰を残す。ツモ \mj{4p} の完全一向聴への変化を見ると打 \mj{1p}。
\item \mj{2244p}：他家が早めに \mj{1p} を切ってくると \mj{2p} がまだ山にありそうなら打 \mj{4p} としてツモ \mj{2p} を狙う。ツモ \mj{4p} の完全一向聴への変化を見ると打 \mj{2p}。
\item \mj{3355p}：ツモ \mj{2p} で \mj{14p} の待ちを狙うと打 \mj{5p}。赤ドラを狙うと \mj{3p}。
\item \mj{4466p}：場況次第。
\end{itemize}

手牌\ref{mj:tobitt-1}からは \mj{7m} や \mj{9m} を切る。\mj{9m} の双碰に狙うなら \mj{7m} を切る。ツモ \mj{6m} による完全一向聴に狙うなら \mj{9m} を切る。手牌\ref{mj:tobitt-2}の場合、愚形があるから先に愚形を払う。ツモ \mj{5s} による両面変化を見て \mj{8s} から切る\citep[][pp. 114--117]{uzaku-haikouritsu}。

\tehai{7799m 123456p 5578s}[][tobitt-1]

\tehai{237799m 123456p 68s}[][tobitt-2]


\subsubsection*{2335の形}

例え手牌の中に \mj{2335m} の形があったら、要すると、1ブロックにしたい時は打 \mj{5m}、1面子と一つの頭が欲しいなら打 \mj{2m}。\mj{5m} がなくなると受け入れが \mj{134m} になったが、\mj{134m} のどちらをツモったら \mj{33m} の部分が頭になれない。\mj{2m} を切ったら \mj{33m} が頭になって、受け入れが \mj{5m} のくっつきで広がるが、ツモ \mj{37m} でピンフが消え、ツモ \mj{7m} で嬉しくない愚形が残す。

\tehai{123m 2335p 123789s -9m}[][nssg-1]

\begin{table*}[!tb]
  \centering
	\begin{threeparttable}
  \begin{tabular}{llll}
    切る牌 & \mj{2p} & \mj{5p} & \mj{9m} \\
    \hline
    受け入れ & \mj{789m34567p} & \mj{789m1234p} & \mj{1234567p} \\
    枚数 & 8種28枚 & 7種24枚 & 7種24枚 \\
    ピンフがつくツモ & \mj{46p} (28.6\%) & \mj{9m3p}\tnote{1} (20.8\%) & \mj{3p}\tnote{1} \mj{46p} (41.7\%)
  \end{tabular}
  
  \smallskip
  \footnotesize
  \begin{tablenotes}
  \item[1] 待ちが \mjfn{124p} になるが \mjfn{2p} で和了ればピンフがない
  \end{tablenotes}
  
  \caption{手牌\ref{mj:nssg-1}の受け入れ}
  \label{tab:nssg-1}
  \end{threeparttable}
\end{table*}


\subsubsection*{1224の形}

1224の辺張の部分を払って辺張の受けの \mj{3m} を引いても損にならない\citep{gendai-penchan}。多くの場合で6889の形の中の9を余り牌として切ると大した差は出ないが、\mj{5m} と \mj{7m} を引いたら \mj{9m} が残る方が少し有利（\mj{8m}とあと一つのペアの双碰または \mj{56m} の両面）。ですので、\mj{9m} を残す基準は孤立の么九牌や役牌との比較\citep{yuusee-6889}。

\begin{table}[ht]
  \centering
  \begin{tabular}{lll}
    形 & \mj{6889m} & \mj{688m} \\
    \hline
    枚数 & 6枚 (\mj{7m}$\times4$ + \mj{8m}$\times2$) & 6枚 (\mj{7m}$\times4$ + \mj{8m}$\times2$)
  \end{tabular}
  \caption{1面子になる枚数の比べ}
\end{table}

\begin{table}[ht]
  \centering
  \begin{tabular}{lll}
    形 & \mj{6889m} & \mj{688m} \\
    \hline
    変化 & \mj{567m} + \mj{88m} & \mj{567m} + \mj{88m} \\
     & \mj{56m} + \mj{789m} & 
  \end{tabular}
  \caption{\mj{5m} と \mj{7m} を引いた場合（頭はもう一つできている上に）}
\end{table}

一向聴に絞るとき基本的に \mj{9m} 切りは効率的な方であるが、手牌\ref{mj:2-1}または\ref{mj:2-2}のような切る候補は孤立の么九牌や役牌の場合 \mj{9m} を残す方は少し有利である。

\tehai{6889m1235p167s44z-8s}[東1局　西家　ドラ\mj{7s}][2-1]

\tehai{6889m1235p67s445z-8s}[東1局　西家　ドラ\mj{7s}][2-2]

手牌\ref{mj:2-3}と\ref{mj:2-4}のようなテンパイまで遠い場合は孤立の么九牌やオタ風を切る方がマシである。

\tehai{6889m125p1668s44z-5z}[東1局　西家　ドラ\mj{7s}][2-3]

\tehai{6889m125p668s445z-2z}[東1局　西家　ドラ\mj{7s}][2-4]

手牌\ref{mj:2-5}のように切る候補は \mj{9m} と孤立の役牌の場合、役牌を残し方がいい。しかし、\ref{mj:2-6}のように横引きが期待でき立直を目指す手牌なら孤立の役牌を切る方がいい。

\tehai{6889m125p2668s44z-5z}[東1局　西家　ドラ\mj{7s}][2-5]

\tehai{26889m237p456s44z-5z}[東1局　西家　ドラ\mj{7s}][2-6]


\subsubsection*{1334の形}

孤立牌と比較するとき、タンヤオが見える場合（孤立牌は37）は \mj{1334p} から \mj{1p} を切る。孤立牌は尖張牌なら、孤立牌のくっつきはタンヤオのバトルネックと二度受けより価値が高い\citep{yuusee-1334}。よって、\ref{mj:1334-1}は打 \mj{1s}。

\tehai{2227m 226p 1334678s}[東1局　西家　5巡目　ドラ\mj{8s}][1334-1]

孤立牌が28になると孤立牌と \mj{1334s} は微差。通常は孤立牌を切るが、\mj{2s} が1枚切れになると打 \mj{1s} の方がちょっといい。よって、\ref{mj:1334-2}から \mj{8m} を切るか \mj{1s} を切るかは場況次第。

\tehai{2228m 226p 1334678s}[東1局　西家　5巡目　ドラ\mj{8s}][1334-2]

役牌暗刻の場合、\mj{1334s} の \mj{2s} 待ちが重要になり、通常は孤立牌より残す。しかし、役牌暗刻の他に暗刻があり、三暗刻が見える場合は、\mj{1s} は弱くなって、場に1枚切れなどの場況で覆されることもあり、\mj{6m6p} など赤が使いやすい孤立牌より先に出ていく場面もある（微差だが）。よって、\ref{mj:1334-3}は打 \mj{6p6s}、\ref{mj:1334-4}は打 \mj{1s}。手牌\ref{mj:1334-5}は場況次第：\mj{2s} が1枚切れなら打 \mj{1s}。

\tehai{1236m 6p 133477s 555z}[東1局　西家　5巡目　ドラ\mj{8s}][1334-3]

\tehai{2226m 6p 133477s 555z}[東1局　西家　5巡目　ドラ\mj{8s}][1334-4]

\tehai{2227m 3p 133477s 555z}[東1局　西家　5巡目　ドラ\mj{8s}][1334-5]


%--------------------------------------------
\subsubsection{112246の形}

\begin{table}[ht]
  \centering
  \begin{tabular}{ll}
    切る牌 & 条件 \\
    \hline
    \mj{1m} & タンヤオに向かう \\
    \mj{2m} & ピンフに向かう \\
    \mj{6m} & タンヤオもピンフもならない
  \end{tabular}
  %\caption{キャップション}
\end{table}

\tehai{112246m34p34566-6s}[][1-1]

\tehai{112246m789p3467-8s}[][1-2]

\tehai{112246m999p3467-8s}[][1-3]

\tehai{112246m34678s11-1z}[][1-4]

手牌\ref{mj:1-1}はタンヤオが見えるので、\mj{1m} を切って仕掛けでもしてタンヤオを目指す。手牌\ref{mj:1-2}にはタンヤオはちょっと無理だけど、ピンフがつく可能性が高いなら \mj{2m} を切ってリャンカンを残す。手牌\ref{mj:1-3}のような立直のみを目指す場合は \mj{6m} を切って \mj{12m} の双碰を狙う。手牌\ref{mj:1-4}は役牌に絡んでいるのでポン材を残し \mj{46m} の嵌張を落とす\citep{yuusee-112246}。

%--------------------------------------------
\subsubsection{12235の形}

\tehai{4567m 12235p 67788s}[東1局　西家　7巡目　ドラ \mj{2s}][12235-1]

\textbf{12235の形はほぼ面子固定でいい}\citep[][p.~32]{uzaku-nankiru}。手牌\ref{mj:12235-1}は頭のない形。頭ができやすい四連形や \mj{67788s} に手をかけるより \mj{12235p} から \mj{2p} を外すの方がいい。

%============================================
\subsection{面子に絡まる形}

%--------------------------------------------
\subsubsection{三面張}

\tehai{4567m 345667p 3455s}[東1局　東家　4巡目　ドラ \mj{2p}][sanmenchan-1]

\begin{table}[ht]
  \centering
  \begin{tabular}{lll}
     & 受け入れ & 良形聴牌 \\
    \hline
    打 \mj{5s} & \mj{2m} -- \mj{9m235p} -- \mj{9p} 51枚 & \mj{3m} -- \mj{8m2358p} 26枚 \\
    打 \mj{6p} & \mj{47m258p25s} 23枚 & \mj{47m258p25s} 23枚
  \end{tabular}
  \caption{手牌~\ref{mj:sanmenchan-1} の受け入れ。}
  \label{tab:sanmenchan-1}
\end{table}

手牌\ref{mj:sanmenchan-1}から \mj{6p} を落としたら三面張が固定されてドラの \mj{2p} も受けられるので良さそうに見えるが、表~\ref{tab:sanmenchan-1} の通り打 \mj{5s} の方の受け入れが広い。良形（ノベタン含み）聴牌になる枚数でも打 \mj{5s} の方がいい。

%--------------------------------------------
\subsubsection{四連形}

四連形は面子や頭が作りやすい形で非常に強い。だとしても、タンヤオや他の役を狙う場面で四連形を1面子に固定することも珍しくない。

\tehai{1234m 5678p 122233s}[東1局　西家　7巡目　ドラ \mj{9s}][yonren-1]

手牌\ref{mj:yonren-1}に四連形が２つある。\mj{13s} 切ったら受け入れが \mj{1m}--\mj{6m3p}--\mj{9p} と \mj{2s}（\mj{3s}）になってすごく広いだが嵌張．双碰リーのみになる確率も高い。よって、ここは \mj{13s} からではなく \mj{1234m} の \mj{1m} から切る。四連形とはいえ \mj{1234m} はよほど強いではなく、\mj{122233s} は \mj{123s} + \mj{223s} という形で一盃口も見える\citep[][p.~14]{uzaku-nankiru}。

\tehai{2207m 3466p 225678s}[東1局　東家　6巡目　ドラ \mj{7m}][yonren-2]

手牌\ref{mj:yonren-2}は面子オーバー。弱い嵌張 \mj{07m} とはいえ2飜は大きいので鳴く前提で手牌を進む。四連形の \mj{5678s} が伸びるとタンヤオが崩れやすいので、ここで1面子をする。567の三色も見込めば \mj{8s} を切る\citep[][p.~16]{uzaku-nankiru}。


\subsubsection*{一向聴で両嵌張との比較}

いくら面子作るやすくても、一向聴で素直に四連形を1ブロックにする場面が多い。比較同士が弱い両嵌張でも。手牌\ref{mj:yonren-rkc-1}と\ref{mj:yonren-rkc-2}は役があり、どちらも \mj{3p} を切る\citep{yuusee-yonren-rkc}。

\tehai{246m 3456p 2256677s}[東1局　南家　3巡目　ドラ \mj{1s}][yonren-rkc-1]

\tehai{135m 3456p 2278s 555z}[東1局　南家　3巡目　ドラ \mj{1s}][yonren-rkc-2]

両嵌張に手を出すのを考え始める場合はドラなしピンフ系の手牌。ドラがないとき、両嵌張による愚形リーのみのリスクが浮上、\mj{135m} から \mj{5m} 切りによる筋引っ掛け、\mj{246m} から \mj{2m} を切って赤受けと \mj{7m} 引きの両面変化のメリットが大きくなる。表~\ref{tab:yonren-rkc} の通り、\mj{135m2345p} 以外の場合は両嵌張から一枚を外す。

\begin{table}[ht]
  \centering
  \begin{tabular}{ll|ll}
    両嵌張 & 四連形 & ドラ0枚 & ドラ1枚 \\
    \hline
    \mj{135m} & \mj{2345p} & 微差だが \mj{2p} & \mj{2p} \\
    \mj{135m} & \mj{3456p} & \mj{5m} & 微差 \\
    \mj{246m} & \mj{2345p} & \mj{2m} & \mj{2p} \\
    \mj{246m} & \mj{3456p} & \mj{2m} & \mj{3p}
  \end{tabular}
  \caption{両嵌張・四連形・ドラ枚数別ピンフ系手牌の打牌選択}
  \label{tab:yonren-rkc}
\end{table}

ドラが1枚の場合、\mj{135m3456p} 以外の場合は四連形を崩す。また、この選択の基準は場況や切れ枚数による覆すことも多い。両嵌張の受けが1枚切れていると両嵌張を崩し、巡目が深くなると四連形を崩す、など。


\subsubsection*{3456で聴牌の時}

\textbf{基本的には端側の牌を切る}\citep{yoteru-yonrenkei}。\mj{3456m} の形なら \mj{3m} を切るのが正解の場合は多い。そのメリットは下記のよう、

\begin{itemize}
\item 端の方の放銃率が低い
\item 鳴かれにくい
\item 河の強さ：\mj{6m} が出ると \mj{39m} が通りやすいという情報を与えて、\mj{6m} が4枚出ていてノーチャンスという情報も
\end{itemize}

例外として真ん中から切る場合は下記のよう、

\begin{itemize}
\item ドラ、赤ドラ、三色に絡まる場合
\item 形を改良する機会
\item 真ん中の方がより安全
\item 序盤
\end{itemize}


\subsubsection*{ドラや手役に絡まる場合}

ツモドラや上家の赤ドラに備えるため、\mj{3456m} から \mj{6m} を落とすこともよくある。たとえ\ref{mj:yonren-3}は普通に \mj{3s} を落とすが、ドラが \mj{2s} の場合スライドができるように \mj{3s} は残す。

\tehai{3456s 2256p -2'13m 77'7z}[ドラ \mj{2s}][yonren-3]

手牌\ref{mj:yonren-4}のソーズを \mj{344556s} にすると \mj{0s} がチーできなくなる。\mj{234456s} の形（= \mj{2s} + \mj{345s} + \mj{46s}）なら \mj{0s} チーして \mj{2s} が切れる。だが \mj{0s} がもう捨てられているなら \mj{2s} を落とす。

\tehai{2344556s 3367p -55'5z}[][yonren-4]

手牌\ref{mj:yonren-5}には123や234の三色が見えるので \mj{5p} を落とす。

\tehai{123m 2345p 2399s -55'5z}[][yonren-5]

三色は本当にガチで見落としやすいので気をつけてね。手牌\ref{mj:yonren-5a}には \mj{468p} のリャンカンがあって、三色があれば456または678。一方、四連形のソーズには8がないので678の三色がないのは確実。456のために \mj{7s} を切る。

\tehai{3345667m 468p 4567s}[][yonren-5a]


\subsubsection*{改良の機会があれば一盃口の形にしない}

手牌\ref{mj:yonren-6}は \mj{36p} のどちらを切っても待ちが \mj{47p} だが、ツモ \mj{2p} で三面張にする機会を逃がしないため \mj{6p} を切る。つまり一盃口の形にしない。また\ref{mj:yonren-4}の場合も似ていて一盃口に近くにしない。

\tehai{345566p 23488s -6'66z}[][yonren-6]


\subsubsection*{片筋}

9が現物なら6の方が3より安全。8が現物なら5の方が2より安全。詳しくは第 \ref{sec:read} 節を参照。


\subsubsection*{序盤}

後でもう1度 \mj{6m} を引く時に安全な \mj{3m} を切るため序盤では \mj{6m} を切るべき。


\subsubsection*{1234や2345の場合}

3456の場合と違って、1234や2345の場合なら素直に端牌を切るべき。真ん中の牌を引くときスライドができるようにするため。

%--------------------------------------------
\subsubsection{23445の形}

\textbf{23445という形は頭が非常にできやすい形}。手牌\ref{mj:1-5}は一見に \mj{307s} のリャンカンが \mj{4s} で埋まって \mj{79s} が不要に見えるが、実際 \mj{8s} 対子を崩す方の受け入れが広い。頭がない状態で、\mj{2356m} のどちらを引いたら頭ができで \mj{58p} 待ちになる。\mj{67p} が重なったら \mj{34556m} をノベタンにすることもできる。望ましくないツモは \mj{47m} 2種7枚けれど、打 \mj{9s} による両面＋両面固定より受け入れが18枚多いのでそのリスクを負う価値がある。

\tehai{34556m 67p 3407889s}[][1-5]

\begin{table}[ht]
  \centering
  \begin{tabular}{ll}
    切る牌 & 受け入れ \\
    \hline
    \mj{8s} & \mj{234567m5678p} 33枚 \\
    \mj{5m} & \mj{36m58p8s} 16枚 \\
    \mj{9s} & \mj{47m58s} 15枚
  \end{tabular}
  \caption{手牌\ref{mj:1-5}の受け入れ。}
  \label{tab:1-5}
\end{table}

%--------------------------------------------
\subsubsection{両面嵌張} \label{subsec:rmkc}

\mj{455679p} という形は両面嵌張と呼ばれている。\mj{456p} の順子と \mj{579p} のリャンカン、\mj{45p} の両面と \mj{567p} の順子（と \mj{9p} の余り牌）、と両方に解釈できるので、二つ面子できる受け入れは \mj{368p} 3種11枚で、\mj{9p} 引いたら頭ができる。\mj{134556m} の形も両面嵌張\citep[][p. 118]{uzaku-haikouritsu}。

\tehai{67m 11456p 235677 -8s}[][rmkc-1]

\begin{table}[ht]
  \centering
  \begin{tabular}{ll}
    切る牌 & 受け入れ \\
    \hline
    \mj{2s}   & \mj{58m469s} 19枚 \\
    \mj{578s} & \mj{58m14s} 16枚
  \end{tabular}
  \caption{手牌\ref{mj:rmkc-1}の受け入れ}
  \label{tab:rmkc-1}
\end{table}

手牌\ref{mj:rmkc-1}は搭子オーバーの手牌。\mj{67m} と \mj{23s} と \mj{78s} の中から一つ両面を切らなければならない。表~\ref{tab:rmkc-1} によれば \mj{56778s} の部分を四連形や中ぶくれにすると打 \mj{2s} による両面嵌張の構えより3枚損。

\tehai{67m 23456p 235677 -8s}[][rmkc-2]

\begin{table}[ht]
  \centering
  \begin{tabular}{ll}
    切る牌 & 受け入れ \\
    \hline
    \mj{2s} & \mj{5m} -- \mj{8m1p} -- \mj{7p3s} -- \mj{9s} 18種59枚 \\
    \mj{3s} & \mj{5m} -- \mj{8m1p} -- \mj{7p2s4s} -- \mj{9s} 18種59枚 \\
    \mj{7s} & \mj{5m} -- \mj{8m1p} -- \mj{7p1s} -- \mj{5s8s} 17種57枚 \\
    \mj{6m} & \mj{7m1p}-- \mj{7p1s} -- \mj{9s} 17種55枚
  \end{tabular}
  
  \vspace{\stdbls}
  
  \begin{tabular}{ll}
    切る牌 & 頭ができる受け入れ \\
    \hline
    \mj{7s} & \mj{67m2356p2358s} 10種30枚 \\
    \mj{2s} & \mj{67m2356p378s}  9種26枚 \\
    \mj{3s} & \mj{67m2356p278s}  9種26枚 \\
    \mj{6m} & \mj{7m2356s2378s}  9種26枚
  \end{tabular}
  \caption{手牌\ref{mj:rmkc-2}の受け入れ}
  \label{tab:rmkc-2}
\end{table}

頭がない形ならどうでしょう。手牌\ref{mj:rmkc-2}は\ref{mj:rmkc-1}の \mj{11p} 頭が \mj{23p} 両面になって、向聴数が一つ戻る。単騎待ちを外して、単に頭ができる受け入れによれば、打 \mj{7s} で \mj{5678s} の四連形を取るのが最適であろう。

\tehai{345m34455688p348 -8s}[][rmkc-31]

手牌\ref{mj:rmkc-31}でピンズの部分は一見に \mj{345p} と \mj{456p} の順子と \mj{88p} の対子に見えるが、\mj{344568p} の両面嵌張と \mj{58p} の余剰牌に分析できる。切る牌の候補は対子同士の比較なら、\mj{8p} を切ると（\mj{5p} を余剰牌にして）ピンフが確定なので \mj{8s} を残す。

\tehai{2245667m 22678p 45s}[東1局　4巡目][rmkc-34]

手牌\ref{mj:rmkc-34}は \mj{2m} と \mj{2p} の二択。どちらも対子なので一見に差がなさそうが、実際 \mj{2m} を落とすとツモ \mj{3m} で面子もできる \mj{245667m} の両面嵌張になる。打 \mj{2p} は4枚損\citep{yoteru-joui}。


\subsubsection*{両面嵌張の先導形}

両面嵌張から一枚外しの先導形は見逃されやすい。例え、\mj{356778m} の \mj{5m} がなくなる \mj{36778m} の形は一見で \mj{3m} は \mj{6m} の筋なので切りやすい。例え、\mj{356778m} の \mj{8m} がなくなる \mj{35677m} の形でも \mj{3m} は不要に見えそう。

\tehai{348m 34458p 357788s}[][rmkc-5]

\tehai{348m 44568p 357788s}[][rmkc-6]

手牌\ref{mj:rmkc-5}と\ref{mj:rmkc-6}は同じ \mj{8m} と \mj{8p} の二択。ツモ \mj{3p} や \mj{6p} によって両面嵌張ができるのでどちらでも一番孤立されている \mj{8m} を切るべき。


\subsubsection*{ピンフとタンヤオ} 

\tehai{678m 455679p 2238 -8s}[][rmkc-3]

\textbf{暗刻があるタンヤオで和了すると点数がピンフより高いのでタンヤオにする}\citep[][p. 119]{uzaku-haikouritsu}。手牌\ref{mj:rmkc-3}では打 \mj{9p} の受け入れ枚数は打 \mj{2s} と同じ19枚に過ぎない。打 \mj{9p} でタンヤオ狙いか、打 \mj{2s} でピンフ確定か、という選択。タンヤオのみは出アガリでもツモでもピンフのみより符が多い上、鳴くことができるので、ここで打 \mj{9p} にする。

\tehai{23488m 455679p 556s}[][rmkc-32]

手牌\ref{mj:rmkc-32}も同じ、両面嵌張の形があるが打 \mj{9p} でタンヤオが確定される\citep[][p. 10]{uzaku-nankiru}。

\subsubsection*{頭が不確定なら両面嵌張を崩しても受けは減らない}

\tehai{134556m 67p 307889s}[][rmkc-33]

\begin{table}[ht]
  \centering
  \begin{tabular}{lll}
    切る牌 & 受け入れ \\
    \hline
    \mj{1m} & \sfrac{16}{49}：\mj{1234567m5678p34568s} \\
    \mj{9s} & \sfrac{7}{27}：\mj{247m58p46s} 
  \end{tabular}
  \caption{手牌\ref{mj:rmkc-33}の受け入れ}
  \label{tab:rmkc-33}
\end{table}

頭が確定した前提で \mj{134556m} から \mj{1m} を抜いたら \mj{2m} を引いても面子は完成させない。けれど\ref{mj:rmkc-33}の場合なら、頭が浮いている上に \mj{34556m} という形はツモ \mj{2356m} で1面子＋1雀頭ができるので \mj{1m} を切っても \mj{2m} が役に立つ\citep[][p.~18]{uzaku-nankiru}。

今回リャンカンの \mj{307s} があるのでピンフが確定されていない。さらに頭ができやすい形があるので、頭を両向聴の時に固定して真っ直ぐにいく必要がない。打 \mj{1m} としたら下記のように複数の場面が挙げられている。

\begin{itemize}
\item ツモ \mj{2356m67p35s} で搭子が対子になるので搭子の他の牌と \mj{8s} が不要になる
\item ツモ \mj{47m58p4s} で搭子が面子になるので打 \mj{8s} によってヘッドレス一向聴になる
\item ツモ \mj{6s} なら \mj{39s} が不要になる
\end{itemize}

従って、両向聴の時既に受け入れが打 \mj{9s} の方より広い上に、一向聴の時全ての場面の平均受け入れもより広い。

%--------------------------------------------
\subsubsection{1345の形}

\textbf{1345の1は孤立の28より価値が高い}\citep{ittaku}。ツモ2で両面、ツモ36なら嵌張ができる。孤立の2(8)は同じ3種類の受け入れだがツモ1で辺張ができてしまうので1345より劣っている。よって、\ref{mj:1345-1}は打 \mj{8p}。

\tehai{1345m 2338p 5889s 11z}[東1局　西家　3巡目　ドラ \mj{1p}][1345-1]

%https://www.youtube.com/watch?v=n10soiTBuJs

%--------------------------------------------
\subsubsection{13457の形}

\textbf{13457の1は孤立の28より価値が低い}\citep{ittaku}。先述の1345の形に似ているが7があったら1が不要になる。13457の1を切っても、ツモ6で三面張、ツモ2でも23457という両面ができやすい形になる。手牌\ref{mj:13457-1}は打 \mj{1m}。

\tehai{13457m 2668p 23699s}[東1局　東家　4巡目　ドラ \mj{6p}][13457-1]

%https://www.youtube.com/watch?v=cyPv6MBksyg

%--------------------------------------------
\subsubsection{多面張} \label{sec:tmc}

\begin{table}[ht]
  \centering
  \begin{tabular}{lll}
    形 & 分析 & 待ち \\
    \hline
    \mj{3334568m} & \mj{333m} + \mj{456m} + \mj{8m}   & \mj{8m} \\
                  & \mj{33m}  + \mj{345m} + \mj{68m}  & \mj{7m} \\
    \mj{3335567m} & \mj{333m} + \mj{55m}  + \mj{67m}  & \mj{58m} \\
                  & \mj{33m}  + \mj{35m}  + \mj{567m} & \mj{4m} \\
    \mj{3335678m} & \mj{333m} + \mj{5678m}            & \mj{58m} \\
                  & \mj{33m}  + \mj{35m}  + \mj{678m} & \mj{4m} \\
    \mj{3334556m} & \mj{333m} + \mj{456m} + \mj{5m}   & \mj{5m} \\
                  & \mj{33m}  + \mj{345m} + \mj{56m}  & \mj{47m} \\
    \mj{3335777m} & \mj{333m} + \mj{5m}   + \mj{777m} & \mj{5m} \\
                  & \mj{333m} + \mj{57m}  + \mj{77m}  & \mj{6m} \\
                  & \mj{33m}  + \mj{35m}  + \mj{777m} & \mj{4m} \\
    \mj{3344455m} & \mj{33m}  + \mj{444m} + \mj{55m}  & \mj{35m} \\
                  & \mj{345m} + \mj{345m} + \mj{4m}   & \mj{4m} \\
    \mj{3444567m} & \mj{34567m} + \mj{44m}            & \mj{258m} \\
                  & \mj{3m}   + \mj{444m} + \mj{567m} & \mj{3m} \\
    \mj{3334445m} & \mj{333m} + \mj{444m} + \mj{5m}   & \mj{5m} \\
                  & \mj{333m} + \mj{44m}  + \mj{45m}  & \mj{36m} \\
                  & \mj{345m} + \mj{33m}  + \mj{44m}  & \mj{34m} \\
    \mj{3334455m} & \mj{333m} + \mj{44m}  + \mj{55m}  & \mj{45m} \\
                  & \mj{33m}  + \mj{345m} + \mj{45m}  & \mj{36m} \\
    \mj{3334456m} & \mj{333m} + \mj{44m}  + \mj{56m}  & \mj{47m} \\
                  & \mj{33m}  + \mj{34m}  + \mj{456m} & \mj{25m} \\
    \mj{3334567m} & \mj{333m} + \mj{4567m}            & \mj{47m} \\
                  & \mj{33m}  + \mj{34567m}           & \mj{258m} \\
    \mj{3334555m} & \mj{333m} + \mj{45m}  + \mj{55m}  & \mj{36m} \\
                  & \mj{333m} + \mj{4m}   + \mj{555m} & \mj{4m} \\
                  & \mj{33m}  + \mj{34m}  + \mj{555m} & \mj{25m}
  \end{tabular}
  \caption{7枚の牌を使っている多面張}
  \label{tab:nanatmc}
\end{table}

多面張とは、最小単位の待ち（単騎・辺張・嵌張・双碰・両面）が複数存在している手牌の組み合わせ。\mj{23456m} のような形の三面張は、待ちが \mj{147m} なので多面張に該当し、もっとも知られている多面張である。\citet[pp. 124--125]{uzaku-haikouritsu}が書いた表~\ref{tab:nanatmc} に挙げられている形は全て7枚を使っていて、多面張に該当する条件は下記の通り。

\begin{enumerate}
\item 2面子＋浮き牌
\item 1面子＋嵌張・両面＋頭
\end{enumerate}

「2面子＋浮き牌」の場合、その浮き牌はあと1面子に関連し、\mj{3334568m} の場合を除け、頭ができやすい。よって、13枚の手牌の他6枚の部分には頭がない時、7枚の多面張は十分に機能する。逆に他6枚の部分に頭がすでにできているなら手牌は非常に弱い。

\tehai{3334555m123p3478s}

孤立の \mj{4m} を切って、待ちが \mj{2s} --- \mj{9s} の28枚になる超広い一向聴。暗刻あり頭なし、そして他の搭子は全て両面、有効牌を引いたら必ず良形聴牌になる。

\tehai{3334555m123p3468s}

残りの搭子は全て良形でなくても、最大に6ブロックができているので、\mj{68s} の愚形を払って良形聴牌に期待する。

\tehai{3334555m139p4478s}

頭がもうできているので、向聴戻ししない前提で孤立の \mj{4m} や \mj{9p} を切らなければならない。\mj{9p} 切っても残りの \mj{4m} は浮き牌だけとして働いていて、受け入れが狭い。

\tehai{3334555m179p2266s} % 8p 26s

7枚の多面張を残しても、受け入れが僅かな \mj{8p26s} の8枚で、非常に狭い一向聴。

\tehai{340677m 3456p 5557s}[東1局　東家　7巡目　ドラ \mj{2s}][tamenchan-1]

手牌\ref{mj:tamenchan-1}から \mj{7m} を1枚外して \mj{34067m} の三面張を確定したいという気持ちがわかるけど、そうしたらツモ \mj{27p} でピンズの三面張にならない上、ツモ \mj{7m} で \mj{3406777m} の多面張（\mj{23568m}）にならないので、ここで \mj{7s} を切って多面張待ちになる確率が一番高い\citep[][p.~22]{uzaku-nankiru}。

\subsubsection*{他の多面張}

\tehai{99m11123456p456s}

待ちは \mj{9m147p} 4種10枚。

%--------------------------------------------
\subsubsection{面子固定と頭固定}

ピンフ系の手牌の中にたった一つ対子を含む形は中ぶくれや亜両面なら、面子固定によって受け入れを最大にするか、頭固定によって好形（6枚以上聴牌）を狙うか、\textbf{他の部分・場況・巡目次第です}\multicitep{uzaku-haikouritsu, pp.~46--55; yuusee-1223fix}。要すると：

\begin{itemize}
  \item 頭固定
  \begin{itemize}
    \item 巡目が早くて、他家から中張牌がまだ出てこない階段
    \item 他の部分で縦引きのツモ枚数は横引きや頭ができるツモ枚数より多い
    \item 好形・手役が求められている
  \end{itemize}
  \item 面子固定
  \begin{itemize}
    \item 巡目が深い
    \item 他の部分で頭ができやすい
    \item 単騎を取っても良形聴牌を取り戻しやすい
  \end{itemize}
\end{itemize}

\begin{table*}[t]
  \centering
  \begin{tabular}{l|l|llll}
       & 頭固定 & 面子固定 \\
    形 & 両面聴牌 & 両面聴牌 & 6枚聴牌 & 双碰聴牌 & 単騎聴牌 \\
    \hline
    両面＋両面＋順子     & 16枚 & 12枚 & 0枚 & 0枚 & 16枚 \\
    両面＋三面張        & 19枚 & 18枚 & 8枚 & 0枚 & 11枚 \\
    両面＋ \mj{56789p} & 15枚 & 18枚 & 8枚 & 0枚 & 7枚 \\
    両面＋ \mj{66778p} & 15枚 & 14枚 & 8枚 & 0枚 & 7枚 \\
    両面＋ \mj{77889p} & 15枚 & 10枚 & 0枚 & 8枚 & 7枚 \\
    二度受けの両面＋順子 & 12枚 & 12枚 & 4枚 & 0枚 & 8枚
  \end{tabular}
  \caption{搭子の形による面子固定と頭固定の差}
  \label{tab:mentsu-or-head}
\end{table*}

孤立の両面＋両面なら、\mj{1223p} や \mj{1233p} などの形を面子にするとヘッドレス一向聴になり、受け入れが28枚になって、一番多い。だが、28枚のうち、16枚のツモでどちらの両面が面子になると単騎になってしまうので、実際好形聴牌の枚数は両面両面の縦引きの12枚だけ。しかし、両面と面子の関係により、表~\ref{tab:mentsu-or-head} の通り、好形（6枚以上の聴牌）になる枚数から見ると、面子固定が優れている場面もある。

例え手牌\ref{mj:moh-1}から \mj{2m} を切って面子に固定すると、ツモ \mj{14s} でソーズの部分が先に埋まると、打 \mj{7p} で亜両面で聴牌できる。一方、\ref{mj:moh-2}でツモ \mj{14s} の後、\mj{8p} を切っても亜両面ができないので、双碰や単騎しかできない。ですので、\ref{mj:moh-1}は面子固定がオススメが、\ref{mj:moh-2}は頭固定がオススメ。

\tehai{1223m 66778p 23678s}[][moh-1]

\tehai{1223m 77889p 23678s}[][moh-2]

三面張の場合、三面張が先に面子化しても、一旦ダマして、三面張からできた面子の部分でいい待ちを作ることが期待できる。例えば、\ref{mj:moh-3}で \mj{8p} を引いたら一旦立直しないで、\mj{345678p} の形で \mj{235689p} を引いたら三面張が復活できる。

\tehai{1223m 34567p 23678s}[][moh-3]

「単騎になってしまうけど他の面子でノベタンや亜両面ができたら復活できる」という性質は他の形にも適用するが、単騎になる瞬間から復活するまで巡目がかかっているので、面子固定や頭固定の選択の瞬間の好形受け入れより重要な論点ではない。


\subsubsection*{他の要素による判断}

\begin{itemize}
  \item 頭固定に寄る要素
  \begin{itemize}
    \item 巡目：早い巡目なら、両面待ちのために受け入れを減っても余裕がある
    \item ドラ：ドラが両面に受けられている、もしくは両面搭子に含まれている
    \item 手役：一盃口、タンヤオなどの手役が頭を求めている
    \item 点数状況：役なしノベタン聴牌もうれしくない状況
  \end{itemize}
  \item 面子固定に寄る要素
  \begin{itemize}
    \item 巡目：深い巡目なら聴牌機会を優先する
    \item ドラ：面子化するとドラを引いた時スライドができる
    \item 暗刻があるので、頭は必要ではない
  \end{itemize}
\end{itemize}


\subsubsection*{搭子が両面ではない場合と他の選択肢}

他の搭子は両面両面ではないと、両面聴牌になる受け入れが激減するが、縦引きの枚数が減らないので、面子固定が有利になる。例え\ref{mj:moh-4}で頭固定すると両面聴牌になるツモが \mj{4p} 4枚だけ。一方、面子に固定すると受け入れが24枚になり、両面聴牌になるツモは \mj{35p} 6枚。ツモ \mj{14s} で嵌張聴牌。

\tehai{1223m 35789p 23678s}[][moh-4]

手牌\ref{mj:moh-5}の受け入れは表~\ref{tab:moh-5} の通り、打 \mj{9p} と受け入れが16枚になり、（亜）両面聴牌が確定されている。一方、打 \mj{2m} で面子固定にすると、受け入れが29枚となり、そのうちの20枚を引いたら6枚以上の聴牌が取れる。打 \mj{2m}。\mj{67889p} という形はツモ \mj{5689p} で1面子＋1雀頭ができるので頭がない時価値が高まる。

\tehai{1223m 67889p 23678s}[][moh-5]

\begin{table}[ht]
  \centering
  \begin{tabular}{lll}
    切る牌 & \mj{2m} & \mj{9p} \\
    \hline
    受け入れ & \mj{56789p1234s} 29枚 & \mj{2m58p14s} 16枚 \\
    6枚以上聴牌 & \mj{5689p14s} 20枚 & \mj{2m58p14s} 16枚
  \end{tabular}
  \caption{手牌\ref{mj:moh-5}の受け入れ}
  \label{tab:moh-5}
\end{table}

手牌\ref{mj:moh-6}の受け入れは表~\ref{tab:moh-6} の通り、面子固定の方が亜両面にする方より受け入れも好形聴牌枚数も多いし、嵌張や単騎になる時も一盃口がつくのでダマできる。打 \mj{2m}。

\tehai{1223m 66788p 23678s}[][moh-6]

\begin{table}[ht]
  \centering
  \begin{tabular}{lll}
    切る牌 & \mj{2m} & \mj{6p} \\
    \hline
    受け入れ & \mj{56789p1234s} 29枚 & \mj{2m58p14s} 16枚 \\
    6枚以上聴牌 & \mj{5689p14s} 20枚 & \mj{2m58p14s} 16枚 \\
    嵌張聴牌 & \mj{23s} 6枚 & 0枚 \\
    単騎聴牌 & \mj{7p} 3枚 & 0枚
  \end{tabular}
  \caption{手牌\ref{mj:moh-6}の受け入れ}
  \label{tab:moh-6}
\end{table}

手牌\ref{mj:moh-7}の一盃口の部分が端に寄るので受け入れが表~\ref{tab:moh-7} の通り、面子固定の方も亜両面にする方も6枚以上の聴牌の枚数が同じ16枚。手牌価値から見ると、同じ1飜でも、一盃口は嵌張で待つことができるので符計算のためピンフより点数が高い。一盃口がつく枚数もピンフがつく枚数多いので、打 \mj{2m} の方がちょっと有利。

\tehai{1223m 77899p 23678s}[][moh-7]

\begin{table}[ht]
  \centering
  \begin{tabular}{lll}
    切る牌 & \mj{2m} & \mj{7p} \\
    \hline
    受け入れ & \mj{6789p1234s} 25枚 & \mj{2m69p14s} 16枚 \\
    6枚以上聴牌 & \mj{679p14s} 16枚 & \mj{2m69p14s} 16枚 \\
    嵌張聴牌 & \mj{23s} 6枚 & 0枚 \\
    単騎聴牌 & \mj{8p} 3枚 & 0枚
  \end{tabular}
  \caption{手牌\ref{mj:moh-7}の受け入れ}
  \label{tab:moh-7}
\end{table}

手牌\ref{mj:moh-8}はマンズの部分が亜両面となり、表~\ref{tab:moh-8} の通り、面子固定しなくても好形聴牌が打 \mj{1m} の方より多い。好形聴牌の枚数が上位なので、ここは面子固定せず打 \mj{7p}。

\tehai{1123m 77899p 23678s}[][moh-8]

\begin{table}[ht]
  \centering
  \begin{tabular}{lll}
    切る牌 & \mj{1m} & \mj{7p} \\
    \hline
    受け入れ & \mj{6789p1234s} 25枚 & \mj{14m69p14s} 20枚 \\
    6枚以上聴牌 & \mj{679p14s} 16枚 & \mj{14m69p14s} 20枚 \\
    嵌張聴牌 & \mj{23s} 6枚 & 0枚 \\
    単騎聴牌 & \mj{8p} 3枚 & 0枚
  \end{tabular}
  \caption{手牌\ref{mj:moh-8}の受け入れ}
  \label{tab:moh-8}
\end{table}


%============================================
\subsection{6ブロック理論} \label{subsec:6block}

14枚の絵を描くため4面子＋1雀頭が必要。面子と雀頭を「ブロック」と呼ぶなら手牌の最終形は5ブロック\footnote{国士無双と七対子を除け。}。しかしながら、聴牌より前の階段なら、手牌の中にもう一つのブロックの空間がある。このようなブロックを６つ維持するながら手牌を進める進行は「6ブロック理論」と呼ばれる\citep[][p. 16]{uzaku-haikouritsu}。

6ブロックは搭子の選択が多くて、手役や待ちが自由に選ばれる。同じ両向聴の5ブロックより受け入れも多い。両面固定による先切りが効く場面も多い\citep[][p. 17]{uzaku-haikouritsu}。巡目が進むほど多くなる情報（他家のポンや山読みなど）によって搭子の選択もできる\citep{yuusee-6block-1, yuusee-6block-2}。

一方、ブロックを六つ維持するには両面固定や面子固定が多くて、搭子のフォローや変化は限られている、そのせいで6ブロック打法は一向聴が狭くなるし、特に完全一向聴になりづらい。手牌を目一杯で構えることにより、安牌や打点に直結する孤立牌も持ちづらい。そして、面子手は必ず5ブロックで、一向聴に進むとき搭子を一つ落とさなければならないで、その残った牌は（安全でもない）余剰牌となり、最低でも1巡残る\multicitep{uzaku-haikouritsu, p~18; yuusee-6block-1}。

%\tehai{23 677m 23 99p 23456s}[東1局　南家　6巡目　ドラ \mj{4p}][block-1]
%
%例え手牌\ref{mj:block-1}に6つのブロックがあり、\mj{23m9p} を落とすのは5ブロックの打法、打 \mj{7m} とするのは6ブロックの打法。5ブロックの打法なら、完全一向聴になりやすくて、三色の可能性も残っている。一方、6ブロックの打法なら、一向聴の時受け入れが4種16枚しかない。それが6ブロック打法の最大な欠点\citep[][p. 18]{uzaku-haikouritsu}。

%---------------------------------------------
\subsubsection{6ブロックと5ブロックの基準}

基準として\multicitep{uzaku-haikouritsu, p.~22--23; yoteru-6block; yuusee-6block-1; yuusee-6block-2}、

\begin{itemize}
  \item 6ブロックにする：
  \begin{itemize}
    \item 真っ直ぐに進めば愚形濃厚
    \item 打点・手役に絡まる
    \item 搭子同士に明確な優劣がない
  \end{itemize}
  \item 5ブロックにする：
  \begin{itemize}
    \item 搭子同士に明確な優劣がある
    \item 聴牌速度を優先する
    \item ポン材が求める副露手
    \item 守備力が求める
  \end{itemize}
\end{itemize}

\textbf{真っ直ぐ進めると手が安いまたは愚形濃厚なら6ブロックにする}\citep[][p. 23]{uzaku-haikouritsu}。平場の東場なら\ref{mj:block-3}は完全一向聴にするべき。だが、点数が少ない南場なら、ここで打 \mj{7m} でピンフを確定してツモ \mj{24p} による三色の可能性もある。点数が必要の場面なら孤立のドラや手役に絡まる牌を1ブロックとして残す選択もある。

\tehai{234 677m 3 88p 234 78s}[][block-3]

両面同士に明確な差がない時6ブロックにして手牌を先延ばしのもあり\citep{yoteru-6block}。手牌\ref{mj:block-make6a}に搭子が六つあり、\mj{78m23p45s} のどちらを払って5ブロックにするより、\mj{334m} を両面に固定して6ブロックで進むのはバランスいい。

\tehai{334 78m 23 678p 45 88s}[東1局　南家　3巡目　ドラ \mj{3z}][block-make6a]

手牌\ref{mj:block-make6b}には打 \mj{4s} で6ブロックにすると打 \mj{6z} で5ブロックにする二択\footnote{\mjfn{7s} 切られたら三色が不確定になる以上ピンフが消滅するのでオススメはしない。}。この手に345の三色の可能性があり、\mj{6z} が暗刻になる可能性もあり、\mj{6z} が暗刻明刻にならなくても一枚切られていたら守備にも十分。どちらの両面が先に埋まったら \mj{6z} 落として立直・ピンフも狙える。それぞれの要素で構う6ブロックは早めに手牌を決める5ブロックよりいい。よって、\ref{mj:block-make6b}は打 \mj{4s}\citep{yuusee-6block-1}。

\tehai{345m 34 78p 445 77s 66z}[東1局　南家　3巡目　ドラ \mj{8p}][block-make6b]

手牌\ref{mj:block-make6c}は234の三色と \mj{3m} 先切りによる \mj{1m} の出やすさによると6ブロック打法は有利。メンピンドラドラも悪くないが、今回の場面は第一打なので、6ブロックにする余裕がある。巡目が遅くなるなら打 \mj{67m} で5ブロックにするのもあり\citep[][pp.~42--43, p.~212]{fkt}。

\tehai{23367m 3466p 23440s}[1巡目　ドラ \mj{3p}][block-make6c]

手牌\ref{mj:block-make5a}は好形ばかり、一見に打 \mj{3m} で三面張を固定して良さそうだが、こんなにいい手牌は「真っ直ぐに進めば愚形濃厚」という条件に当たらず「打点・手役に絡まる」にも当たらない。よって、ここは5ブロックにして、\mj{4405s} を1ブロックにするのは効率的\citep[][pp.~40--41]{fkt}。並び対子を2ブロックとするケースは他にも（特に一盃口を狙いたいとき）あるが、今回のような手役の重要度はそんなに高くない状況なら効率的に両面対子にすべき。

\tehai{3340567m 23p 440578s}[3巡目　ドラ \mj{3m}][block-make5a]

手牌\ref{mj:block-make5b}は辺張が明確的に嵌張より劣っているので \mj{89s} を落とす\citep[][p. 21]{ishibashi-black}。

\tehai{11 667m 68p 34 89s 333z}[][block-make5b]

手牌\ref{mj:block-make5c}にある \mj{6667p} は1ブロック（\mj{666p} + \mj{7p}）や2ブロック（\mj{66p} + \mj{67p}）に見える。今回は \mj{6667p} を1ブロックにすればヘッドレスなので受け入れが最大7種24枚だが、そのうちの \mj{789m} 9枚（\mj{8m} の1枚がドラ表示牌）を先に引かなければ待ちが弱い。よって、ここでは \mj{7p} を残って \mj{8m} を切る。\mj{6667p} の形ならツモドラ \mj{9m} でも打 \mj{7p} でドラが残る\citep[][p.~16]{uzaku-nankiru}。

\tehai{23489m 6667p 34789s}[東1局　東家　7巡目　ドラ \mj{9m}][block-make5c]

手牌\ref{mj:block-make5d}はポン \mj{7z} が入ったばかりの場面。\mj{3m} を切るより\textbf{ポン材を大事にする}\citep{yoteru-6block}。そして副露手は面子手と違って、食い伸ばしや嵌張の両面変化ができるので、搭子を目一杯にしなくても構わない。例え\ref{mj:block-make5d}が\ref{mj:block-make5e}に発展してきたら、チー \mj{1245m} で両面両面へか、ツモ \mj{5p} で両面へ変化することができる。

\tehai{33478m 2477p 68s -77'7z}[][block-make5d]

\tehai{233478m 2477p -77'7z}[][block-make5e]

\textbf{辺張は2枚孤立の役牌より先払うべき}\citep[][pp.~154--155]{suphx}。手牌\ref{mj:block-make5f}に6ブロックと字牌2枚ある。\mj{12m} を残すメリットは4枚の \mj{3m} で1面子を完成させることだけ。他の搭子が先に埋まられたらピンフが消滅してうれしくないリーのみになる。一方、6枚の \mj{15z} の一枚を引けば仕掛けができる上、満貫も見込める。一向聴の時 \mj{1m} や \mj{2m} より字牌の方も安全。こう見ると \mj{12m} を残して6ブロックにするのは損。Suphxは\ref{mj:block-make5f}から \mj{12m} を落とした。

\tehai{12m 445578p 67s 1336z}[東1局　親　1巡目　ドラ \mj{7p}][block-make5f]

\textbf{役・打点を見込めば両面も払える}\citep[][pp.~156--157]{suphx}。手牌\ref{mj:block-make5g}には混一色に必要なブロックは既に５つあって、\mj{23p} を残して他の搭子を払うと打点がひどく下がる。よって、Suphxはここで \mj{23p} を落とした。孤立牌の \mj{6s} より落とすのはドラ \mj{7s} の受けを保つため。

\tehai{3357m 23p 6s 1144556z}[東1局　親　4巡目　ドラ \mj{7s}][block-make5g]


\subsubsection*{紙一重の場面}

先に触れられて「搭子同士に明確な優劣がない時は6ブロックにする」けど、比較同士が辺張という弱い搭子になったら諸説になる。例え\ref{mj:block-makewhich1}、\citet{yoteru-6block}は「一番弱い辺張は選択できないで、ダイレクトの \mj{37m} を逃さないように」という理由で \mj{6s} を切って6ブロックにする。一方、\ref{mj:block-makewhich2}の牌姿は\citet{yuusee-6block-2}は「一向聴が狭くなって、ピンズソーズが伸びられない」という理由で辺張を払う。

\tehai{234 78m 12 99p 12 667s}[東1局　南家　3巡目　ドラ \mj{3z}][block-makewhich1]

\tehai{12 89m 22 455p 45678s}[東1局　西家　3巡目　ドラ \mj{9s}][block-makewhich2]

実際に「こんな手牌なら6ブロックにするしかない」でもない。同じ\ref{mj:block-makewhich1}でも、\citet{yoteru-6block}は「ラス目の親番でどうしても連荘したい」という場面で辺張を払って聴牌に真っ直ぐ向かうのもあり、と言及している。

%---------------------------------------------
\subsubsection{搭子の落とし順}

\subsubsection*{違う形の搭子の比較}

\textbf{6ブロックを5ブロックに絞る時に搭子の落とし順は「辺張＞嵌張＞対子＞両面」}\citep[][pp. 21--22]{ishibashi-black}。例え、手牌\ref{mj:block-4}は両面が対子より優れているので \mj{1m} を落とすのは効率が一番良いい。手牌\ref{mj:block-5}なら辺張は全ての搭子の中で一番弱いので \mj{12p} を払う。

\tehai{11 668m 23 56p 34s 333z}[][block-4]

\tehai{11 668m 12 68p 35s 333z}[][block-5]

両面より嵌張を先に落とすのは嵌張の1枚がドラでも変わらない\citep{yuusee-5block}。手牌\ref{mj:block-make5c}の \mj{79s} は \mj{7s} から切って、ドラ \mj{9s} が重なったら \mj{3p} を落とす。

\tehai{234788m 3367p 3479s}[南1局　北家　5巡目　ドラ \mj{9s}][block-make5c]

対子が嵌張より強い理由は下記の通り\citep[][pp. 30--33]{uzaku-haikouritsu}\citep{yuusee-kanchan-toitsu}、

\begin{itemize}
\item 対子を二つ持つことで両面変化に対応できるし・完全一向聴になりうる
\item 待ちとして対子が端に寄ると嵌張より出やすい。対子が字牌になると尚更。
\item 2種類があるので後筋\footnote{立直する後の捨て牌で出来た筋をいう。}になりやすい
\item 待ちが隣の数牌を使っていないのでノーチャンス\footnote{中の方の隣の数牌が四つ捨てられていて、その牌でできる両面が否定されていること。}になりやすい
\end{itemize}

単純に両面変化と複合形変化によると、外嵌張＋対子一つは28対子二つより弱い。対子が二つなら、どちらの両面変化や複合形への変化が効く。それに対して、嵌張＋対子という組み合わせは、手牌を完成させるために対子が必要なので、変化できるのは嵌張だけ。ですので、直接の受け入れは同じだが、両面変化と複合形変化の枚数は対子二つの方が多い。詳しくは表~\ref{tab:24p2288s} を参照。

\begin{table}[ht]
  \centering
  \begin{tabular}{lll}
    形 & \mj{24p22s} & \mj{2288s} \\
    \hline
    受け入れ & \mj{3p} 4枚 & \mj{28s} 4枚 \\
    両面変化 & \mj{5p} 4枚 & \mj{37s} 8枚 \\
    両嵌張への変化 & \mj{6p} 4枚 & \\
    嵌張対子への変化 & & \mj{46s} 8枚
  \end{tabular}
  \caption{外嵌張＋対子二つの打牌選択の比較}
  \label{tab:24p2288s}
\end{table}

\mj{1199s}のような両面変化がない対子といっても、他の部分が先に埋まられたら最終の待ちとして後筋やノーチャンスになりやすい。対子が字牌となると出アガリもかなり期待できる。よって、手牌\ref{mj:block-make5d}も嵌張を落とす。ツモ \mj{3m} でも \mj{11m} と複合できるので裏目が一番小さい \mj{4m} から落とす\citep{nemata-12}。

\tehai{1146m 11 34578p 666s}[ドラ \mj{5p}][block-make5d]

嵌張を残したい場面は三色・タンヤオなどの手役やドラが絡むとき。リーのみになる可能性があれば尚更。手牌\ref{mj:block-make5e}は\ref{mj:block-make5d}のドラを \mj{5m} に変えて、\mj{46m} の価値が高まった手。ここはツモ \mj{2p} とその後の一通を見て、\mj{1p} を浮かせる\citep{nemata-12}。

\tehai{1146m 11 34578p 666s}[ドラ \mj{5m}][block-make5e]

\textbf{6ブロックを5ブロックに絞る時に安牌があれば不要な搭子を先に落とす}\citep[][pp.~152--153]{suphx}。Suphxは\ref{mj:block-6}から \mj{68s} を先に落とすのは(1)ツモ \mj{7s} で面子を完成させても打点に付加価値がないことと(2)一向聴の時 \mj{1z} は \mj{68s} より安全。

\tehai{456m 2245p 68s 1226 -6z}[南2局　親　4巡目　ドラ \mj{6p}][block-6]


\subsubsection*{同じ形の搭子の比較}

同じ形の搭子同士でも隣の搭子や面子で優劣が分かる。同じ形の搭子の比較基準は下記のように\multicitep{suphx, pp.~152--153; yuusee-follow; nemata-13}、

\begin{enumerate}
\item （嵌張・辺張）受け入れの残り枚数
\item （両面・辺張）直接の受け入れが裏目になるときのフォロー
\item （嵌張）両面変化・リャンカン変化が裏目になるときのフォロー
\item （対子）ドラスライドができるか否か
\item （対子）浮かせると変化があるか否か
\end{enumerate}

手牌\ref{mj:block-follow1}に \mj{89m} と \mj{12s} の辺張が二つある。白がポンしたら対子を二つ維持したいのでここは \mj{3m} 切らない。\mj{12s} を落とした後裏目の \mj{3s} が引いたら \mj{57s} の嵌張と両嵌張ができるので、\mj{89m} がなくした後要らなくなる \mj{7m} と比べるとその痛さはよほどないので、ここは \mj{12s} は切る\citep{yuusee-follow}。

\tehai{33489m 057p 1257s 55z}[東1局　北家　3巡目　ドラ \mj{7z}][block-follow1]

手牌\ref{mj:block-follow2}も \mj{89m} と \mj{12s} の二択。今回は \mj{55m} の対子があって、\mj{89m} 切ったら裏目の直接の受け入れの \mj{7m} が引いたら嵌張対子ができるし、裏目の嵌張への変化の \mj{6m} が引いたら両面対子ができるので、打 \mj{12s} より損失が少ない。なんの情報がないなら \mj{89m} を落とす。ちなみに \mj{9m} を切ってすぐに \mj{7m} を引いたら振聴両面は残す。さらに、\mj{3s} が1枚切りになったらこの基準を覆して \mj{12s} を切る方がいい\citep{yuusee-follow}。

\tehai{5589m 123678p 1288s}[東1局　東家　2巡目　ドラ \mj{8s}][block-follow2]

手牌\ref{mj:ist-1-1}は \mj{56m} と \mj{67p} の二択。\mj{56m} を落として \mj{4m} 引いても \mj{123m} の搭子が \mj{234m} へスライドできるので \mj{56m} を落とす\citep{gendai-iishanten-yojouhai}。

\tehai{12356m2267p23456s}[][ist-1-1]

ドイツ同士の比較なら\textbf{ドラへスライドができる方を優先に残す}\citep[][p.~20]{uzaku-nankiru}。手牌\ref{mj:block-9}に6ブロックの中に対子が \mj{99p} と \mj{11s} がある。他の搭子が全て両面なので対子落しにする。今回のドラが \mj{4s} で、\mj{11234s} の形でツモ \mj{4s} の時打 \mj{1s} で \mj{12344s} に入れ替えることができるので、ここで \mj{99p} を落とす。

\tehai{45m 67899p 1123467s}[東1局　西家　7巡目　ドラ \mj{4s}][block-9]

\textbf{牌は繋がっている方を残すのが基本}\citep[][p.~30]{uzaku-nankiru}。手牌\ref{mj:block-10}にピンフが見えて、対子 \mj{3m} と \mj{7s} のどちらを落とすという問題。この巡目は同じ4種16枚だが、\mj{56m78p} が重なった時差が出る。表~\ref{tab:block-10} の通り、\mj{3m} を落とせば \mj{4s} も受けられるので打 \mj{7s} の方より受け入れが多い。

\tehai{3356m 23478p 56777s}[東1局　東家　7巡目　ドラ \mj{8p}][block-10]

\begin{table}[ht]
  \centering
  \begin{tabular}{ll}
    \multicolumn{2}{l}{\mj{356m234788p56777s}} \\[3pt]
    打牌 & 受け入れ \\
    \hline
    \mj{3m} & \mj{47m689p47s} 23枚 \\
    \\
    \multicolumn{2}{l}{\mj{3356m234788p5677s}} \\[3pt]
    打牌 & 受け入れ \\
    \hline
    \mj{7s} & \mj{347m689p} 20枚
  \end{tabular}
  \caption{手牌\ref{mj:block-10}から牌を切ってドラ \mj{8p} を引いた形}
  \label{tab:block-10}
\end{table}

%---------------------------------------------
\subsubsection{複数の対子の捌き方}

前述したように、通常、対子が二つの場合、嵌張などより弱い搭子があるなら対子を優先する。両面しかないなら対子を払う。この小節はどちらの対子を先に払うのを、そして対子が三つあればどうするのを、もっと詳しく説明していく。

\subsubsection*{対子外し}

\textbf{対子が３組の場合は基本的に対子を優先的に外す}\citep{gendai-toitsu}。\textbf{3対子より2対子の方の受け入れが2枚多いから}\citep[p.~34][]{uzaku-haikouritsu}。このようなどちらの対子を全く払うのは「対子外し」という\citep{nemata-13}。対子の優先度は下記の項目によって決まる\multicitep{gendai-toitsu; uzaku-haikouritsu, p.~169; uzaku-nankiru, p.~8}：

\begin{itemize}
\item 変化の多さ
\item 手役との絡み
\item 待ちとしての強さ
\end{itemize}

手牌\ref{mj:toitsu-1}は \mj{15m1p} の三択。ツモ \mj{3m} で \mj{11355m} というウィング形を見込めば、\mj{11p} が一番不要になる。また、場況・ドラによって \mj{11m11p} が待ちとして強いなら \mj{55m} を落とすのもある。

\tehai{1155m11p678p34789s}[][toitsu-1]

手牌\ref{mj:toitsu-2}で \mj{22p} が縦引いたらタンヤオがあって、\mj{13p} 来たら一通があるので一番不要な \mj{11s} を落とす。

\tehai{77m 22456789p 1134s}[][toitsu-2]

手牌\ref{mj:toitsu-3}は打 \mj{1s} と打 \mj{2s} の受け入れは同じ6種18枚、\mj{78p} の両面が先に埋められたら \mj{1s7z} の待ちが \mj{2s7z} よりいい。よって、ここは打 \mj{2s} とする。

\tehai{06778p 1122345s 77z}[][toitsu-3]

手牌\ref{mj:toitsu-4}で \mj{1m} 落として \mj{2p} 引いても \mj{66m} が頭候補になるのでピンズの変化もちゃんと残る。打 \mj{1m}。

\tehai{1166m 11 345 78p 666s}[ドラ \mj{6m}][toitsu-4]

%立体図~\ref{rittai:stts-1} のようなオーラスで、最下位との点差は僅かな1600点。なんの役があっても早く上がりたい状況。\mj{1m}、\mj{4s} と \mj{8s} の選択。一見に打 \mj{1m} だとしたらタンヤオがつきそうだが、表~\ref{tab:stts-1} によれば打 \mj{4s} や打 \mj{8s} でピンフを狙う方がより早い。567の三色もありそうなので \mj{8s} をツモ切りする。
%
%\begin{rittai}[h]
%  \centering
%  
%  \drawrittai{
%    jicha/seat = 南,
%    jicha/hand = 11567m 56p 445678 -8s,
%    jicha/river = 66^z 2z,
%    jicha/point = 14300,
%    %
%    toimen/river = 1s 2p 9s,
%    toimen/point = 39900,
%    %
%    kami/river = 4z 1p 2z 9m,
%    kami/point = 33100,
%    %
%    shimo/river = 4z 9m 4^z,
%    shimo/point = 12700,
%    %
%    kyoku = 南4局,
%    dora = 5z
%  }
%  
%  \caption{オーラスの場面で対子が三つ揃った。}
%  \label{rittai:stts-1}
%\end{rittai}
%
%\begin{table}[ht]
%  \centering
%  \begin{tabular}{lll}
%    切る牌 & 受け入れ & 枚数 \\
%    \hline
%    \mj{1m} & \mj{47p48s} & 4種12枚 \\
%    \mj{4s} & \mj{1m47p3689s}  & 7種23枚 \\
%    \mj{8s} & \mj{1m47p3469s}  & 7種23枚
%  \end{tabular}
%  \caption{立体図~\ref{rittai:stts-1} で自分の手の受け入れ。}
%  \label{tab:stts-1}
%\end{table}
%
%飛び対子や並び対子がある場合、一盃口を狙うなら3対子を維持する。一盃口にしない場合、飛び対子や並び対子の方から1枚を外す。\mj{2244m} の形があって、1面子で良しとすると打 \mj{2m}、もう一つの搭子が欲しければ打 \mj{4m} により \mj{4m} のくっつきを期待する。


\subsubsection*{浮かせ打ち}

浮かせ打ちとは、\textbf{変化・手役のために数牌の対子を1枚外して、孤立した牌へのくっつきを見る}こと\multicitep{uzaku-haikouritsu, p.~168; gendai-toitsu}。例え\ref{mj:toitsu-5}のような双碰待ちのリーのみ手、打 \mj{7m} としてピンズの一通や789の三色を見込む。

\tehai{77m 3456789p 22789s}[東1局　西家　ドラ\mj{2z}][toitsu-5]

手牌\ref{mj:toitsu-6}は\ref{mj:block-make5e}と同じ、ここではなぜ \mj{1m} ではなく \mj{1p} を落す理由をもっと詳しく説明する。ピンズの形で \mj{2p} を引いたら \mj{1234578p} で両面確定が見て、一通も見える。しかし、そうなると他の部分に頭が必要。つまり、この変化に \mj{11m} の頭が必要。一方、\mj{11m} を落とすと \mj{11p} が頭としてしか働きなくて、ツモ \mj{2p} でも変化できない。

\tehai{1146m 11 345 78p 666s}[ドラ \mj{5m}][toitsu-6]

似た牌姿だが、\ref{mj:toitsu-4}と\ref{mj:toitsu-6}の違いは、後者は「頭候補としてのトイツ」を求めている。変化を見て面子化に向かう対子を1枚落とさなければならない。前者の「頭候補としてのトイツ」が \mj{66m} なので、この \mj{11p} が「面子候補としてのトイツ」として十分に働くことができる。そうすると \mj{11m} が一番不要になって、単純の対子落としとなる。

%============================================
\subsection{孤立牌}

孤立牌は、単独で他の搭子や対子や面子と繋がっていない牌。孤立のオタ風 \mj{4z} や役牌 \mj{5z} や数牌 \mj{3p} が該当する。\textbf{孤立牌の重要度はまず有効牌の量を従って、「オタ風 \textless{} 役牌 $\approx$ 19牌 \textless{} 28牌 \textless{} 3--7牌」の順番となる}\citep{clearrain-floater}。

手牌・場況によって役牌の重要度は19牌より高くしたり、同じクラスの数牌の中でも、隣の孤立牌や搭子や面子で弱くしたり強くしたりするので、この後詳しく説明していく。だが、数牌の場合、あるクラスの中でどれだけ強くても、一つ上の数牌より強くことはない。つまり、「19牌 \textless{} 28牌 \textless{} 3--7牌」の順番は変わらない。両面ができるツモの枚数が絶対的だから。

\subsubsection{同じクラスの浮き牌の比較}

\subsubsection*{中張浮き牌の比較}

表~\ref{tab:float-taatsu} の通り、\mj{5p} から嵌張ができたら、その嵌張の両面変化や両嵌張変化が \mj{3467p} からできた嵌張より多いので、聴牌まで遠い階段で \mj{5p} の価値が高い。聴牌に近い階段で、\mj{37p} の浮き牌から搭子ができると端に近いので待ちとして強いので、\mj{456p} より残す。また、タンヤオを狙う場合、\mj{13p} の嵌張ができたら、もしくは \mj{23p} の両面で \mj{1p} を引いたら、タンヤオが崩れるので \mj{456p} より弱い。

\begin{itemize}
\item 聴牌まで遠くて、変化を見る：37 $\approx$ 46 \textless{} 5
\item 聴牌に近くて、待ちを見る：5 \textless{} 46 \textless{} 37
\item タンヤオを狙う：37 \textless{} 46 $\approx$ 5
\end{itemize}

\begin{table}[ht]
  \centering
  \begin{tabular}{llll}
    孤立牌 & \mj{37p} & \mj{46p} & \mj{5p} \\
    \hline
    両面変化 & \mj{23p}、\mj{34p} & \mj{34p}、\mj{45p} & \mj{45p}、\mj{56p} \\
    嵌張変化 & \mj{13p}、\mj{35p} & \mj{24p}、\mj{46p} & \mj{35p}、\mj{57p}
  \end{tabular}
  \caption{孤立牌とそれからできた搭子}
  \label{tab:float-taatsu}
\end{table}


\subsubsection*{浮き牌から３つ離れた浮き牌は弱い}

\mj{25m2p} の形で孤立の \mj{2m2p} の重要度は \mj{2m} \textless{} \mj{2p}。\mj{2m} の有効牌は \mj{1234m}、そしてそのうちの2種類は \mj{5m} の受け入れと被っているので \mj{2p} と比べると切った後のロスは痛くない。


\subsubsection*{浮き牌から４つ離れた浮き牌は強い}

\mj{158m1p} の形で孤立の \mj{18m1p} の重要度は \mj{1p} \textless{} \mj{1m} \textless{} \mj{8m}。ツモ \mj{3m} で両嵌張ができるので、\mj{1p} より残す。ただし、\mj{1m} から両面ができないので、筋である \mj{8m} より弱い\multicitep{clearrain-floater; yuusee-158}。よって、\ref{mj:158}は打 \mj{1p}。

\tehai{4688m 158p 34579s 77z}[東1局　東家　3巡目　ドラ \mj{5s}][158]


\subsubsection*{順子に近い浮き牌は強い}

\mj{3567p3s} の形で孤立の \mj{3p} に、\mj{4p} がついた時三面張ができて、\mj{8p} がつくと嵌張＋面子ができる。そうすると \mj{3567p} の受け入れが \mj{3s} より多い\citep{clearrain-floater}。

\mj{1456m1567p} の形で \mj{3m} を引いたら \mj{13456m} という形の両面変化は \mj{457m}。一方、\mj{3p} を引いたら \mj{13567p} の両面変化は \mj{4p} だけ。よって、この \mj{1p} は \mj{1m} よりちょっと弱い\citep{yoteru-floater}。


\subsubsection*{暗刻に近い浮き牌の強さは頭の有無による}

\mj{3334p} のリャンタン\footnote{両面単騎の略。}や \mj{3335p} のカンタン\footnote{嵌張単騎の略。}は他の部分に頭がない時、頭＋搭子と機能して、待ちとしても最低7枚。ですが、頭がある時、他の孤立牌と比べると、暗刻に近い浮き牌の受け入れは暗刻に消費されているので、もう一つ搭子を作る確率は高くない。


\subsubsection*{搭子に近い浮き牌は弱い}

\mj{356m367p368s} の形で孤立の \mj{3m3p3s} の重要度は \mj{3m} \textless{} \mj{3s} \textless{} \mj{3p}。これらの牌にくっついても現存の搭子にコンボできないので、\textbf{ここは浮き牌を捨てる後裏目の痛さを比較する}。\mj{3m} を切っても \mj{4m} が \mj{56m} に拾えられるので、\mj{3m} の裏目は一番痛くない。\mj{3s} を切った後 \mj{4s} が \mj{68s} と両嵌張になるので、裏目もよほど痛くない。両面を逃す上、直接に面子にならないので \mj{3m} よりちょっと残したい。そして \mj{3p} 切った後 \mj{4p} がただの孤立牌になる。よって、\mj{3p} の裏目は一番痛い。

ですが、\mj{356m} の形が頭がない時は価値が高くなる。詳しくは\ref{sec:134}参照。

%---------------------------------------------
\subsubsection{258の5が赤ドラ}

\mj{208s} の形は \mj{28s} が \mj{0s} の筋になっているので受け入れが被っているが、\mj{28s} のどちらを切ると \mj{37s} の両面変化やその後の変化を逃すのは痛い。まずブロックが揃っていない時は触らない。前述の浮き牌の重要度基準に従うと、他色の28がある時 \mj{208s} から \mj{28s} を切り始めるといい。

手牌\ref{mj:208-1}のように、\mj{4s} を引いて孤立の \mj{28s} から1枚を落とす場面、二度受けを嫌って \mj{2s} を切るか、ツモ \mj{3s} の四連形を狙うか、微妙なところです。諸説あるので、実際の場況で判断する\citep{yuusee-208}。

\tehai{24899m 36678p 2458s}[東1局　南家　3巡目　ドラ \mj{4z}][208-1]

手牌\ref{mj:208-1a}は \mj{5m} が重なった手。もう5ブロックができているので \mj{28m} を処理する。123三色の可能性を残すために \mj{8m} を先に落とす\footnote{どこで拾ったNAGA牌譜：\url{https://naga.dmv.nico/htmls/8f456393c9dc421dc9e62068fe9e07161db5c59cd25b2e5ce967b6e5627950bfv2_2.html?tw=0}}。

\tehai{2508m 13356p 12399s}[東1局　北家　5巡目　ドラ \mj{6p}][208-1a]


\subsubsection*{役牌との比較}

\mj{208s} と孤立の \mj{5z} の比較は場況に左右されやすい。タンヤオや面子手で進めば孤立の役牌を切る。ドラ3など門前より仕掛けの価値が高まる場面、もしくは染め手、七対子など \mj{208s} の横引きの価値が低まる場合なら \mj{28s} を処理する\citep{yuusee-208}。

手牌\ref{mj:208-2}はタンヤオが見える。手牌\ref{mj:208-3a}はドラ1の面子手で、素直に立直に向かう。よって、両方も打 \mj{5z}。

\tehai{222677m 3588p 208s 5z}[東1局　南家　3巡目　ドラ \mj{4p}][208-2]

\tehai{112678m 1257p 208s 5z}[東1局　南家　3巡目　ドラ \mj{4z}][208-3a]

手牌\ref{mj:208-3b}はドラ3なので、\mj{5z} が重なればうれしいので、\mj{28s} を切る。ドラ1の手で役牌が重なってもあまりうれしくないのは、仕掛けて手牌の価値が2600に過ぎない。一方、立直・ドラ1に一発や裏ドラがついたら満貫になる可能性が高い。白・ドラ3ならどう鳴いても満貫あるので、非効率的打牌をしても「打点効率」としては十分。

\tehai{112678m 1257p 208s 5z}[東1局　南家　3巡目　ドラ \mj{1m}][208-3b]

手牌\ref{mj:208-4}は混一色を見て、\ref{mj:208-5}は七対子。よって、両方は打 \mj{28s}。

\tehai{1124678m 5p 208s 445z}[東1局　南家　3巡目　ドラ \mj{4p}][208-4]

\tehai{112677m 88p 208s 445z}[東1局　南家　3巡目　ドラ \mj{4p}][208-5]

%---------------------------------------------
\subsubsection{孤立の役牌}

役牌が重なると仕掛けができるので、「聴牌への急行券」とも思われる。が、役牌は字牌と同じ、横引きができない。それに、生牌字牌が重ならないまま持っていくと、字牌の危険度が巡目が深くなるほど高まる。役牌と19牌（または28牌）の選択はこの後詳しく説明していく。


\subsubsection*{生牌を先に処理するケース}

\textbf{相手がヤバそうに見えるなら生牌から切る}\citep[][pp.~192--193]{suphx}。生牌役牌が重なると仕掛けできるが、僅かな1000点で無理矢理仕掛けて守備力を下げるのは無謀。特に立体図~\ref{rittai:jihai-1} では上家の捨て牌は変則なので、今生牌の \mj{1z} を切らないと巡目が深くなるほど \mj{1z} の危険度が高まる。

\begin{rittai}[h]
  \centering
  
  \drawrittai{
    jicha/seat = 北,
    jicha/hand = 45m 11557p 2667s 17z -1p,
    jicha/river = 1m,
    jicha/point = 22100,
    %
    toimen/river = 6z,
    toimen/point = 25000,
    %
    kami/hand = XXXXXXXXXX -66'6z,
    kami/river = 4s 6m,
    kami/point = 24000,
    %
    shimo/river = 7z,
    shimo/point = 28900,
    %
    kyoku = 東1局,
    honba = 1,
    dora = 7m
  }
  
  \caption{1枚見えの \mj{7z} と生牌の \mj{1z} の二択}
  \label{rittai:jihai-1}
\end{rittai}

立体図~\ref{rittai:jihai-2} では \mj{1z} が重なれば打点が3900から7700まで上昇する。けれどここはトップ目、ブロックも足りているのでSuphxは安全度を優先して \mj{1z} を切った\citep[][p.~194]{suphx}。

\begin{rittai}[h]
  \centering
  
  \drawrittai{
    jicha/seat = 西,
    jicha/hand = 1134488s 134 -4z -555'z,
    jicha/river = 6m 65p,
    jicha/point = 47800,
    %
    toimen/river = 33^s 2^p 9s,
    toimen/point = 7200,
    %
    kami/river = 9^s 3z 8^m 6s,
    kami/point = 13500,
    %
    shimo/river = 3z 2s 7z,
    shimo/point = 31500,
    %
    kyoku = 東4局,
    honba = 2,
    dora = 1m
  }
  
  \caption{2枚見えの \mj{3z} と生牌の \mj{1z} の二択}
  \label{rittai:jihai-2}
\end{rittai}

立体図~\ref{rittai:jihai-3} では重なれば$2000\to7700$もしくは12000になるドラの \mj{4z} だが、一向聴かつすぐ中盤なのでSuphxは \mj{4z} を手放した\citep[][p.~195]{suphx}。これも一向聴の時手牌をスリム化して打点より安全度を優先する打筋。

\begin{rittai}[h]
  \centering
  
  \drawrittai{
    jicha/seat = 北,
    jicha/hand = 34m 2205p 78s 47z -9s -1'11z,
    jicha/river = 1m 1p 2s 8m,
    jicha/point = 25000,
    %
    toimen/river = 1p 17m 3z 7p,
    toimen/point = 25000,
    %
    kami/river = 1m 6z 7m 4s,
    kami/point = 22100,
    %
    shimo/river = 73^z 2m 1z 3s,
    shimo/point = 27900,
    %
    kyoku = 東1局,
    honba = 1,
    dora = 4z
  }
  
  \caption{1枚見えの \mj{7z} と生牌ドラの \mj{4z} の二択}
  \label{rittai:jihai-3}
\end{rittai}

\textbf{打点もないバラバラな手牌は生牌字牌を残さない}\citep[][p~.197]{suphx}。立体図~\ref{rittai:jihai-4} はアガリの可能性がほとんどない手牌。わざと役牌を重ねて上がり率を上げようとすることはない。Suphxは生牌の \mj{7z} を切った。

\begin{rittai}[h]
  \centering
  
  \drawrittai{
    jicha/seat = 北,
    jicha/hand = 23m 3p 1245s 234567z -5p,
    jicha/river = 9m 8p,
    jicha/point = 24400,
    %
    toimen/river = 3z 1m 5z,
    toimen/point = 17400,
    %
    kami/river = 4z 1m 5^z,
    kami/point = 24400,
    %
    shimo/river = 1s 1p 5z,
    shimo/point = 33800,
    %
    kyoku = 東3局,
    dora = 5m
  }
  
  \caption{5向聴かつバラバラの手牌}
  \label{rittai:jihai-4}
\end{rittai}

立体図~\ref{rittai:jihai-5} では仕掛けて7700の決定打。Suphxは2枚切れの \mj{7z} を切らず生牌の \mj{6z} を手放した\citep[][p.~198]{suphx}。いくら一歩も早く上がりたくても、トップ目かつ4向聴の手牌なので、Suphxは安全度を優先した。

\begin{rittai}[h]
  \centering
  
  \drawrittai{
    jicha/seat = 東,
    jicha/hand = 116m 1206p 34s 1567z -1m,
    jicha/river = 9m 23z 8p,
    jicha/point = 34000,
    %
    toimen/river = 9p 9s 1m 7z,
    toimen/point = 22700,
    %
    kami/river = 1s 43z 9^p,
    kami/point = 18300,
    %
    shimo/river = 57z 89^m,
    shimo/point = 25000,
    %
    kyoku = 東4局,
    dora = 5p
  }
  
  \caption{仕掛けて7700の手牌}
  \label{rittai:jihai-5}
\end{rittai}


\subsubsection*{生牌字牌を残すケース}

\citet[pp.~200--201]{suphx}によると

\begin{itemize}
\item 一色手・対子手が見える場合
\item 高打点だが役なしの手牌
\end{itemize}

\citet{yuusee-yakuand1}と\citet{yuusee-yakuand2}による孤立の \mj{12m} と役牌の基準は下記の通り：

\begin{itemize}
  \item 役牌を \mj{1m} より先に残す場面：
  \begin{itemize}
    \item \mj{14m} のような孤立の筋牌が手牌の中にある
    \item ブロックが足りない
    \item ブロックが足りていても愚形がある
    \item ピンフが薄い
    \item \mj{1m} はドラに絡まっていない
  \end{itemize}
  \item 役牌を \mj{2m} より先に残す場面：
  \begin{itemize}
    \item 序盤0面子、かつ、愚形搭子が二つ以上
    \item 端牌がドラ
  \end{itemize}
\end{itemize}

役牌が頭にするとピンフが消滅するので、ピンフが見える場合は役牌を先に切るべき。たとえ\ref{mj:keepyaku-1}の牌姿、数牌が連続して頭がない。\mj{5z} の縦引きより他の数牌の縦引きが嬉しいのでここで \mj{1s} ではなくて \mj{5z} を切る。

\tehai{27m 2348p 145667s 5z -1p}[東2局　北家　ドラ \mj{8m}][keepyaku-1]

立体図~\ref{rittai:keepyaku-2} の場合はドラが2枚あるけどブロックが4つだけで、重なった時に鳴けるというメリットでここで1枚切の \mj{1z} を残して \mj{1s} を切る。

\begin{rittai}[h]
  \centering
  \captionsetup[table]{labelsep=space}
  
  \drawrittai{
    jicha/seat = 東,
    jicha/hand = 34067m 1168p 16s 67z -1z,
    jicha/river = 3z,
    jicha/point = 21300,
    %
    toimen/river = 3z,
    toimen/point = 19800,
    %
    kami/river = 9m,
    kami/point = 42600,
    %
    shimo/river = 1z,
    shimo/point = 16300,
    %
    kyoku = 南1局,
    honba = 1,
    dora = 6s
  }
  
  \caption{}
  \label{rittai:keepyaku-2}
\end{rittai}

同じ牌姿だが、\ref{mj:keepyaku-3}でNAGAの選択がドラによる。ドラが \mj{7s} になったらNAGAが \mj{6z} を切る。でなければNAGAは \mj{9s} を切る。

\tehai{44m 45p 1109s 44467z -7m}[東4局　北家][keepyaku-3]

手牌\ref{mj:keepyaku-4}は0面子で愚形搭子が２つなのでここで \mj{56z} を残して \mj{2m2s} を切る。

\tehai{268m 455899p 279s 5z -6z}[東4局　西家　3巡目　ドラ \mj{1p}][keepyaku-4]

端牌がドラの時、役牌が重なって鳴いてスピードと上がり率を上げる方がいい。よって、\ref{mj:keepyaku-5}は \mj{8s} 切り。

\tehai{2367m 11699p 348s 56z}[東1局　西家　3巡目　ドラ \mj{1p}][keepyaku-5]


\subsubsection*{孤立ドラ役牌の扱い方}

\textbf{他の役があれば先に処理する。でなければ他の字牌すら安易に切らない}\citep{yoteru-lonedorayakuhai}。たとえ\ref{mj:lonedorayaku-1}のようなバラバラの配牌、\mj{5z} が重ならなければほとんど上がらない上、切って鳴かせたら損になる。守備力にも構えば字牌は安牌として残して、打牌候補は \mj{8p125s} などになる。

\tehai{24668m 8p 12578s 125z}[ドラ \mj{5z}][lonedorayaku-1]

逆に\ref{mj:lonedorayaku-2}のようにタンヤオ平和赤などな長所があればドラ \mj{5z} が不要になり、他家の \mj{5z} が重なる前に第一打として切るのは正解。

\tehai{34m 468p 334056s 125z}[ドラ \mj{5z}][lonedorayaku-2]

手牌\ref{mj:lonedorayaku-3}はメンタンピン一向聴だが厄介なものを掴まえてしまった場合。ドラを手放す方がいい。ドラがポンされたら捨て牌や場況で対応できるし、ポンされなくて代わりに立直が来る場面でも安牌があるので戦いやすい。

\tehai{23467m 34588p 56s 3z -5z}[7巡目　ドラ \mj{5z}][lonedorayaku-3]

混一色に絡まる場合なら満貫が既にあったらドラを手放す方がいい。手牌\ref{mj:lonedorayaku-4}はもう満貫確定なので、断ラスなど点数が切実に必要な場面でなければ上がり率を優先する。

\tehai{11340899s 125z -7'77z}[東1局　南家　7巡目　ドラ \mj{5z}][lonedorayaku-4]


%============================================
\subsection{役作り}

ここまでの話の趣旨は「最速聴牌」と「最高の聴牌形」。立直、ドラ、裏ドラ、一発などの役によって他の手役を狙わなくても満貫になる効率がいいので、手役を作る技術の利用頻度が低い。しかし、速度が同じでも、役の選択により手牌の価値が変動する。また、局収支として「最速聴牌」が正解ではない場面もある。特に打点が欲しい場合、または真っ直ぐに進めばドラなしリーのみの一途になる場合、手役を重視して、純粋な牌効率を逆らうことになる。

手役の同士を比較する方法と、どの程度まで手役を狙うために牌効率が逆らえるか、この節で説明していく。

%---------------------------------------------
\subsubsection{タンヤオ}

\subsubsection*{6ブロックから么九牌のブロックを落とす}

\textbf{6ブロックから么九牌のブロックを落としてタンヤオを確定するのがほぼ正解}\multicitep{gendai-tanyao; zero-teyaku, p.~22}。手牌\ref{mj:tanyao-1}は \mj{46s} のどちらを落とせば受け入れが最大16枚になる。しかし、そのうちピンフがつくツモは \mj{5s} だけ。一方、\mj{9s} を落とせば目前の受け入れが12枚になって4枚損だが、タンヤオが確定されて仕掛けも効くので、速度は実際 \mj{46s} の方より早い。打 \mj{9s}。

\tehai{23467m 345p 446699s}[][tanyao-1]

手牌\ref{mj:tanyao-2}も打 \mj{9m}。牌効率の通りなら打 \mj{2p} だが、他の部分に両面が一つだけ、嵌張待ちになる可能性が高い。打 \mj{9m} とすると、仕掛けが効き、嵌張が頭になる道も解放される。

\tehai{35799m 2467p 24888s}[東1局　南家　3巡目　ドラ \mj{8s}][tanyao-2]

手牌\ref{mj:tanyao-3}も打 \mj{9m}。牌効率の通りなら打 \mj{3m} だが、メンピンで終わる可能性が高い。ドラ引きや345の三色を備えて、打 \mj{9m} の方は打点が高い。

\tehai{356789m 346788p 45s}[東1局　南家　3巡目　ドラ \mj{2m}][tanyao-3]

タンヤオを狙わない場面は下記の二つ\citep[][p.~23]{zero-teyaku}：

\begin{itemize}
\item 好形搭子が十分
\item タンヤオにすると愚形聴牌になる可能性が高い
\end{itemize}

手牌\ref{mj:tanyao-4}は打 \mj{1m} にすると、両面が先に埋まるとカン \mj{3s} になってしまう。両面両面一向聴が十分なので、ここは牌効率の通りに打 \mj{2s}。

\tehai{1156m 34678p 24666s}[東1局　南家　3巡目　ドラ \mj{8p}][tanyao-4]

手牌\ref{mj:tanyao-5}は牌効率の通りに打 \mj{48p}。両面搭子が十分なので、リャンカンの部分を払うと門前聴牌が無理なく狙える。一盃口を見て \mj{8p} を切るか危険な \mj{4p} を先に切るかが場況次第。

\tehai{4599m 234468p 3478s}[東1局　南家　3巡目　ドラ \mj{8s}][tanyao-5]

手牌\ref{mj:tanyao-6}は牌効率の通りに打 \mj{3m}。手牌\ref{mj:tanyao-2}に似ていてドラ3の機会手が、ここは大人しく牌効率を従う。

\tehai{35799m 2367p 23888s}[東1局　南家　3巡目　ドラ \mj{8s}][tanyao-6]

手牌\ref{mj:tanyao-7}もタンヤオを狙わない。\mj{13s} を払えばタンヤオが確定されたが、\mj{25s} の両面が消える上ドラ受けもなくなる。愚形の部分は \mj{68m05p} と \mj{0577p} の受け入れが同じ4枚だが、両面変化の枚数は前者は \mj{5m} 4枚だけ。それに対して後者は \mj{468p} 12枚。よって、ここは嵌張の \mj{68m} を払う\citep[][pp.~32--33]{fkt}。

\tehai{23468m 0577p 12334s}[5巡目　ドラ \mj{5s}][tanyao-7]


%---------------------------------------------
\subsubsection{ピンフ}

\subsubsection*{ドラなしピンフの手筋}

\textbf{ドラなしならリーのみを打たないように両面を固定する}\citep[][pp.~200--201]{doi-prob}。手牌\ref{mj:pinfu-1}は\mj{3m8s} を切ればピンフが確定されている。ここで \mj{3z} をツモった。牌効率の通りならこの \mj{3z} がいらない。しかし、5巡目で自分の手の中にドラがない事実から、他家の手牌の中にドラが通常より多いのが推測できて、先切り・守備力・打点などの要素から考えると、リーのみを避けたいためここで牌効率に従わない。打 \mj{3m}。

\tehai{334789m 56p 11788s -3z}[東1局　西家　5巡目　ドラ \mj{8p}][pinfu-1]

似たような牌姿だが、今回\ref{mj:pinfu-2}にタンヤオが見込むので、タンヤオの優先度が安牌より高い。\mj{788s} を最大限にタンヤオ牌からできるように構う前提で、より危険な \mj{3m} を先切るか、安牌の \mj{3z} を切るかが場況次第。

\tehai{334678m 56p 22788s -3z}[東1局　西家　5巡目　ドラ \mj{8p}][pinfu-2]

%---------------------------------------------
\subsubsection{一盃口}

一盃口は中ぶくれや亜両面などの連続形から自然にできることが多くて、無理に狙う場面が少ない。\mj{33445p} のような両面＋面子の形でも半分の確率で一盃口が消える。\mj{33455p} のような一盃口が確定している形は、自分が待ちの牌をすでに1枚持っているので、もう場に1枚切れになると苦しくなる。両面より \mj{33455p} を先に処理する。2枚見えなら他の嵌張より先に切る\citep{gendai-ipk}。

一盃口を意識していい階段は孤立牌を選ぶ階段とリーのみ濃厚の階段。後者の討論は形次第。前者の討論はドラ枚数次第。手牌\ref{mj:ipk-1}がドラなしなら一盃口で打点上昇のために打 \mj{8p}。ドラ3枚ならダマテンに構えられるために一盃口を狙って打 \mj{8p}。ドラが1枚・2枚なら単純に受け入れの残り枚数で打 \mj{2m} \citep[][p.~32]{zero-teyaku}。

\tehai{12236m 2348p 35799s}[\mj{1m} 1枚切れ][ipk-1]

手牌\ref{mj:ipk-2}はヘッドレス。\mj{1223m} を頭に固定する確率が高くて、一盃口は狙いづらい。打 \mj{8p}。

\tehai{12237m 2348p 23079s}[\mj{1m} 1枚切れ　ドラ \mj{9s}][ipk-2]


\subsubsection*{両面嵌張にするかにいにい形にするか}

\mj{3556778p} という形から \mj{3p} を切るとにいにい形\footnote{\citet{chitist-niinii}がつけた名前。構成が「2枚・1枚・2枚・1枚」から。}になり、\mj{5p} を切ると両面嵌張になる。どちらを切ると受け入れが同じだが、手牌の役・牌の危険度によって最善の一打が変わる。

\begin{itemize}
  \item 両面嵌張
  \begin{itemize}
    \item 面子化する時必ず順子＋順子になる
    \item 危険な内側が先に出る
  \end{itemize}
  \item にいにい形
  \begin{itemize}
    \item 一盃口目がある
    \item 2ヘッドで柔軟
  \end{itemize}
\end{itemize}

\textbf{多くの場合はにいにい形にするのは正解}。手牌\ref{mj:ipk-3}は打 \mj{3p} としたらツモ \mj{5p} や \mj{3s} でテンパネして、ツモ \mj{2s} で \mj{69p} と \mj{14s} を選べることができる。\mj{9p} が急にたくさん切られてくるという緊急事態にも柔軟に対応。赤 \mj{5p} の受けも\citep[][p.~35]{zero-teyaku}。

\tehai{23m 3556778p 33s 777z}[][ipk-3]

そういったメリットは手牌\ref{mj:ipk-4}にはない。役もドラもない手牌、なるべくリーのみを避けたいなら、一盃口が確定されていないにいにい形より、ピンフが確定されている両面嵌張の方がいい。さらに \mj{3p} が \mj{1p} より危険。こう見ると、危険な \mj{3p} を先に切る\citep[][p.~35]{zero-teyaku}。

\tehai{23m 1334556p 33789s}[ドラ \mj{7m}][ipk-4]

手牌\ref{mj:ipk-5}も似ている形の牌姿。今回の牌姿は既にドラ1があって、「なんとなく役をつくる」という念頭がないし、\mj{3m5m} も危険度に差がない。そうすると、打点最大値を見て打 \mj{3m} として一盃口を狙う\citep[][pp.~120--121]{uzaku-haikouritsu}。ただし、トップ目でピンフでダマしたい場面などなら打 \mj{5m} もあり。

\tehai{3556778m 45p 12388s}[東家　11巡目　ドラ \mj{2s}][ipk-5]

\begin{table}[ht]
  \centering
  \begin{tabular}{lll}
    切る牌 & \mj{3m} & \mj{5m} \\
    \hline
    受け入れ & \mj{569m36p8s} 19枚 & \mj{469m36p} 19枚 \\
    ピンフツモ & \mj{69m36p} (78.9\%) & \mj{469m36p} (100\%) \\
    一盃口ツモ & \mj{6m} (15.8\%) & (0\%)
  \end{tabular}
  \caption{手牌\ref{mj:ipk-5}の受け入れ}
  \label{tab:ipk-5}
\end{table}


\subsubsection*{三連飛び対子}

\mj{113355p} のように連なっている形は「三連飛び対子」\citep[][p.~35]{zero-teyaku}や「2回飛んだ飛び対子」\citep{yoteru-tobitoitsu}と呼ばれている。多くの場合は、三連飛び対子をもう一つの搭子とまとめてどちらを払うのを考える\citep[][pp. 114--117]{uzaku-haikouritsu}。表~\ref{tab:santobitt-1} の通り、搭子が対子でない限り、三連飛び対子から真ん中を抜く方は受け入れが多い。けれど、そうするとツモ \mj{24m} でも一盃口ができない。他の搭子は愚形で、リーのみになりやすいし待ちとしても強くないならそのデメリットはひどい。

\begin{table}
  \centering
  \begin{tabular}{lll}
    形 & 打 \mj{3m} の受け入れ & 打 \mj{8s} の受け入れ \\
    \hline
    \mj{113355m78s} & \mj{1245m69s} 20枚 & \mj{12345m} 14枚 \\
    \mj{113355m68s} & \mj{1245m7s} 16枚 & \mj{12345m} 14枚 \\
    \mj{113355m89s} & \mj{1245m7s} 16枚 & \mj{12345m} 14枚 \\
    \mj{113355m88s} & \mj{1245m8s} 14枚 & \mj{12345m} 14枚 \\
  \end{tabular}
  \caption{三連飛び対子＋搭子から1枚を外した後の受け入れ}
  \label{tab:santobitt-1}
\end{table}

例え手牌\ref{mj:santobitt-1}は \mj{113355m99p} から \mj{3m} と \mj{9p} の二択。それぞれのメリットは、

\begin{itemize}
  \item 打 \mj{3m}
  \begin{itemize}
    \item \mj{1m9p} の双碰になるとアガリ率が高い
  \end{itemize}
  \item 打 \mj{9p}
  \begin{itemize}
    \item 一盃口目が残る
    \item 連打できるので安牌やフォロー牌が持てる
  \end{itemize}
\end{itemize}

「双碰として強い」というメリットは、\mj{9p} が場に1枚見えになると消える。そうすると \mj{9p} を切って一盃口を狙う\citep[][p.~35]{zero-teyaku}。

\tehai{113355m 123 99p 567s}[][santobitt-1]

手牌\ref{mj:santobitt-2}の三連飛び対子の部分が中寄りになって、\mj{2m} 双碰が待ちとして弱くなる。ここは \mj{1p} 切りに寄る\citep{yoteru-tobitoitsu}。

\tehai{224466m 11p 789s 123s}[][santobitt-2]


\subsubsection*{一盃口クラッシュ}

一盃口は門前役で、鳴いたら一盃口が消える。が、一盃口の形は頭が作りやすくて、仕掛け手順も意外にも存在する。手牌\ref{mj:ipk-crash}から単騎を取らずに \mj{4p} を落とすという手順は「一盃口クラッシュ」と呼ばれている。上がりづらい単騎より、ポン \mj{35p} とチー \mj{2356p} で聴牌できる広い一向聴をとるべき\citep[][p.~36]{zero-teyaku}。

\tehai{67m 334450p 06s -77'7z}[ドラ \mj{5m}][ipk-crash]


%---------------------------------------------
\subsubsection{混一色・清一色}

混一色を狙う場面は下記の通り、

\begin{itemize}
\item 受け入れが少なくて、愚形が多くて、門前聴牌まで遠い\citep{shibukawa-honitsu}
\item 鳴いても満貫が見える\citep{hirasawa-honitsu}
\item 手なりで立直打っても2飜以下\citep{hirasawa-honitsu}
\item 役牌対子が二つある（すなわち役役）\citep[][pp.~96--97]{zero-teyaku}
\item 役牌対子が一つ、他の部分がバラバラ\citep{zero-teyaku}
\end{itemize}

混一色は、たった一つの色の数牌と字牌しか受け入れないので、牌効率としては非効率的。しかし、仕掛けが効く上、字牌を持っていくことで守備力も十分で、打点も十分。特に立直から遠い悪い手牌こそ、混一色で守備力と打点を維持して進むのがオススメ。

\subsubsection*{役役の場合は混一色を狙う}

麻雀は打点効率として満貫を作るのは効率が一番高い。役牌対子が二つある場合、両飜の役がつけば満貫になるので、混一色やトイトイの価値が高まる。そして混一色の方が受け入れが多いので、役役の場合は常に混一色を意識する\citep[][pp.~96--97]{zero-teyaku}。よって、手牌\ref{mj:honitsu-1}は染めていく。打 \mj{8s}。

\tehai{7m 12456p 228s 44557z}[東1局　北家　1巡目　ドラ \mj{3z}][honitsu-1]

手牌\ref{mj:honitsu-2}はソーズの部分が完成している。が、超序盤（4巡目以内）で役牌が出たらポンして混一色を狙ってソーズの面子を落とす\citep[][p.~97]{zero-teyaku}。

\tehai{12456p 123s 44557z}[東1局　北家　配牌　ドラ \mj{3z}][honitsu-2]


\subsubsection*{役牌対子がない場合は基本的に狙わない}

手牌\ref{mj:honitsu-3}はすでに頭と一盃口が整っていて、\mj{8p26s} にくっつけば立直・タンヤオ・一盃口で打点が十分。一方、混一色を狙って仕掛けたら概ね3900。僅かな清一色の可能性を残して、\mj{37z} から切る\citep{shibukawa-honitsu}。

\tehai{144556688m 8p 26s 37z}[南1局　南家　1巡目　ドラ \mj{5p}　ダンラス][honitsu-3]

手牌\ref{mj:honitsu-4}には5枚の孤立字牌があるが、字牌を重なるのに待つより、数牌のくっつきでタンヤオやピンフを狙う方が早い。ここは字牌を処理する\citep{yuusee-honitsu}。

\tehai{48m 23457p 25s 12456z}[東1局　西家　1巡目　ドラ \mj{1m}][honitsu-4]

手牌\ref{mj:honitsu-5}と\ref{mj:honitsu-6}はともに受け入れが少ない愚形だらけの手牌、門前聴牌は遠く見えるので、守備力のために消極的に染めていく\citep{shibukawa-honitsu}。

\tehai{114478m 246s 44567z}[東3局　東家　2巡目　ドラ \mj{2m}　\mj{4m} 1枚切り][honitsu-5]

\tehai{35m 57p 123449s 2244z}[東3局　東家　2巡目　ドラ \mj{8s}　\mj{4z} 1枚切り][honitsu-6]

手牌\ref{mj:honitsu-yakuyaku-1}にも6ブロックがある。対子が四つなので満貫を作るためトイトイも視野に入る。だとすると一番落としたい部分は \mj{56m} の両面と \mj{35z} の孤立牌。\mj{3z} は自風なので \mj{5z} が先。赤ドラがある場合なら赤受けで \mj{6m} が先。手なりの2000で上がる手順を残したいなら \mj{6z}。混一色とトイトイの可能性を最大限に高めるなら \mj{6m} \citep[][pp.~98--103]{tetsusenryaku}。

\tehai{56m 22p 2499s 113566z}[東1局　西家　1巡目　ドラ \mj{8p}][honitsu-yakuyaku-1]

手牌\ref{mj:honitsu-yakuyaku-2}は\ref{mj:honitsu-yakuyaku-1}に似ているが、マンズの部分が嵌張になり、ソーズの下が両嵌張になる。こうなるとマンズの部分が明確に弱くになって、\mj{46m} を切って混一色に行く\citep[][pp.~98--103]{tetsusenryaku}。

\tehai{46m 22p 24699s 11566z}[東1局　西家　1巡目　ドラ \mj{8p}][honitsu-yakuyaku-2]


\subsubsection*{役牌対子が一つの場合は他の部分次第}

役牌対子が一つ、孤立の役牌が二つなら、役牌の重なりの価値が高まる。これは「役役ダッシュ準備状態」と言い、他の色の孤立の1289を切る\citep[][p.~97]{zero-teyaku}。よって、手牌\ref{mj:honitsu-7}は打 \mj{8p}、\ref{mj:honitsu-8}は打 \mj{8s}。

\tehai{1126m 8p 56889s 1566z}[東1局　北家　1巡目　ドラ \mj{2s}][honitsu-7]

\tehai{34m 1455689p 8s 4556z}[東1局　北家　1巡目　ドラ \mj{2m}][honitsu-8]

手牌\ref{mj:honitsu-9}も打 \mj{28s} でピンズの混一色を狙う\citep{hirasawa-honitsu}。

\tehai{347m 34479p 28s 1157z}[ドラ \mj{1p}][honitsu-9]

手牌\ref{mj:honitsu-10}は混一色を狙わない。孤立の役牌が一つ、そしてマンズの部分は伸びにくい。ここは素直に \mj{6z} を切る\citep[][p.~98]{zero-teyaku}。

\tehai{35999m 379p 48s 5556z}[東1局　北家　1巡目][honitsu-10]

手牌\ref{mj:honitsu-11}も混一色を狙わない。他の色の部分が全て両面なので、立直が近くなる。

\tehai{67m 23p 223056s 5567z}[東1局　北家　1巡目 ドラ \mj{1s}][honitsu-11]

手牌\ref{mj:honitsu-12}は他の色に搭子がない。ピンズの部分に一盃口の形があるが、立直しても2600が現実的。混一色に行くなら最低3900、自風の \mj{4z} か重なったら満貫が見える。オタ風の \mj{3z} もワンチャントイトイの種にもなる。赤受けを考えると打 \mj{7s} \citep[][pp.~104--109]{tetsusenryaku}。

\tehai{5m 22333447p 7s 3477z}[東1局　北家　1巡目 ドラ \mj{6z}][honitsu-12]


\subsubsection*{2番目に多い色から切る}

手牌\ref{mj:honitsu-13}は親からの第一打の \mj{4z} をポンしかける場面。\mj{6m46p} のどちらも不要だが、\mj{6m} を遅らせて他家にソーズ寄せを悟られないようにするのは混一色のコツ。この後マンズを引いたらツモ切り\citep[][pp.~118--119]{suzuki}。

\tehai{6m 46p 14467s 33446z -4z}[東場　北家　1巡目　ドラ \mj{2m}][honitsu-13]

手牌\ref{mj:honitsu-14}はダブ東をポンしたところ。ピンズとソーズがどちらも2枚ずつ。枚数が同じだが、外側切ると河が手なりに見えるのでここは \mj{8p} を切る\citep[][pp.~120--121]{suzuki}。

\tehai{1179m 68p 36s 447z -111'z}[東場　東家　2巡目][honitsu-14]


%---------------------------------------------
\subsubsection{役役の進行}

\textbf{役牌の対子が二つあるなら満貫を目指せ}\citep{yoteru-yakuyaku}。「役役」は手牌の中に二つの役牌の対子のことを指し、最低の得点の2000を7700や満貫にすれば四倍になるので打点上昇のスィートスポット。役役にさらに2飜をつく方法は(1)混一色(2)トイトイ。

例え\ref{mj:yy-1}の手牌は嵌張だらけで、\mj{34s} が唯一の両面。\mj{34s} の主線で進行すればただの2000点。うまくいけなくて役牌対子は一つしか刻子にならない場合は1000点。どうせアガりづらくて安いの手牌ならアガリ率を下げて手牌を染めにいくの方が得。よって \mj{68p34s} の全てが打牌候補。

\tehai{2579m 68p 34s 135577z}[東1局　南家　1巡目　ドラなし][yy-1]

手牌\ref{mj:yy-2}ならトイトイを目指す。トイトイなら中張牌が重なっても嬉しくないのでここで \mj{3m} を切る。

\tehai{223m 119p 28s 135577z}[東1局　南家　1巡目　ドラなし][yy-2]


\subsubsection*{ドラが1枚持っている場合}

\tehai{67m 1330p 46s 135577z}[東1局　南家　1巡目　ドラ \mj{2m}][yy-3]

\textbf{ドラが1枚持っているなら手なりで打つのが基本}。無理矢理染めても3900 $\to$ 8000は二倍に過ぎないので効率的に良くない。よって、手牌\ref{mj:yy-3}は染めないで、素直に \mj{13z} から切る。

\tehai{137m 3799p 124s 5577z}[東1局　南家　1巡目　ドラ \mj{1m}][yy-4]

そういっても満貫が見込めるケースは\ref{mj:yy-4}のようにチャンタがついている手牌。他に三色で満貫にするケースも。

\tehai{13m 67s 5577z -6'0799'9m}[東1局　南家　7巡目　ドラ \mj{8p}][yy-5]

鳴いたけど役牌が一つしか刻子にならない牌姿なら同じ、できれば満貫にする。手牌\ref{mj:yy-5}は混一色が見えるので \mj{67s} の両面は割愛。

\tehai{405s 25577z -1'11p888'm}[東1局　南家　7巡目　ドラ \mj{8p}][yy-6]

手牌\ref{mj:yy-6}の場合も \mj{4s} を切ってトイトイで満貫を目指す。もちろん場況・点棒状況によって素直に \mj{2z} を切るケースもあるけど、平場の東1局なら満貫を取る。


\subsubsection*{役役で真っ直ぐに行くケース}

\tehai{2467m 223p 56s 25577z}[東1局　南家　1巡目　ドラ \mj{1m}][yy-7]

\textbf{両面が三つあるなら素直に行く}。手牌\ref{mj:yy-7}のように両面がいっぱいの手牌を無理矢理染めると効率が悪過ぎて打点がいくら高くてもアガらないと意味がない。けれどやはり良形と愚形がともにある手牌は染めるか否かのは紙一重。


%---------------------------------------------
\subsubsection{三色同順}

「配牌を見たら三色を探せ」は、阿佐田哲也氏が残した言葉。昭和時代の麻雀は一発・赤ドラや裏ドラなどの役がまだ普及されていないため、三色などの手役で打点を作るのが大事。平成になると、麻雀ルールがインフレしてきたため、「最速聴牌」が優先され、手役の重要度が落ちて、「三色は偶然役」という意見も。

そう言っても、ここまでの話もしばしば言及したように、全ての場面で「棒テン即リー」は、科学的に局収支から見ると必ずしも正しくない。よりいい結果に辿り着くため、牌効率を逆らって、もしくはもっと聴牌から遠い階段で工夫を入れて手役の可能性を最大限にする必要がある。そうすると三色を狙う技術と基準が重要。三色を狙う基準は\citet{gendai-ssk}によると：

\begin{itemize}
\item 確定しやすさ
\item 確定しづらい場合、三色を狙うことによる受け入れのロス
\item 三色の２翻がつくことによる打点面のメリット
\end{itemize}

手牌にドラはないなら打点を作るため三色が重要になる。そして手牌の価値が非常高いが門前聴牌まで遠い場合、三色による仕掛けのメリットが高い。

手牌が一定の程度で整っているなら三色は狙いづらい。例え手牌\ref{mj:ssk-1}に123の三色部分が7枚もあるが、頭がないので実は123の順子が三つ揃いづらい。この手牌はピンフ狙いの効率が高い。打 \mj{1m}。手牌\ref{mj:ssk-2}も三色部分以外で面子が二つ完成しているので三色を狙わない。打 \mj{1p} \citep[][p.~67]{zero-teyaku}。

\tehai{122m 12357p 236778s}[南家　5巡目　ドラ \mj{2p}][ssk-1]

\tehai{235678m 13567p 123s}[南家　5巡目　ドラ \mj{8m}][ssk-2]


\subsubsection*{リーのみの場合は受け入れより手役}

手牌\ref{mj:ssk-3}と\ref{mj:ssk-4}は打 \mj{7p} で完全一向聴になる。手牌\ref{mj:ssk-3}はそうしべきだが、手牌\ref{mj:ssk-4}はタンヤオが付かないので完全一向聴にするとリーのみになってしまう可能性がある。打 \mj{6m} によって \mj{16m} の受けが消えるが、ツモ \mj{7m6p} で三色が見える。期待値として打 \mj{6m} の方がいい\citep{clearrain-ssk}。

\tehai{22566m 457p 222567s}[6巡目　ドラなし][ssk-3]

\tehai{11566m 457p 222567s}[6巡目　ドラなし][ssk-4]

手牌\ref{mj:ssk-5}は複合搭子二つ、両面対子が一つの形。打 \mj{4m} と打 \mj{6s} の受け入れが同じ。表~\ref{tab:fkgtt-2} の通りなら、両面聴牌の枚数のために打 \mj{6s}。が、この手牌はリーのみの上、打 \mj{4m} と234の三色が確定するので、ここは打 \mj{4m}。両面リーのみより役が大事\citep[][p.~65]{zero-teyaku}。

\tehai{244m 234p 234667s 44z}[南家　3巡目　ドラ \mj{3z}][ssk-5]

手牌\ref{mj:ssk-6}はと同じだが、ドラが2枚なので素直に良形聴牌率で打 \mj{6s} \citep[][p.~65]{zero-teyaku}。

\tehai{244m 234p 234667s 44z}[南家　3巡目　ドラ \mj{4z}][ssk-6]


\subsubsection*{悪い配牌こそ手役を狙うべき}

手牌\ref{mj:ssk-7}は4向聴でバラバラすぎる悪配牌。上がりに速度が必要なので、仕掛けが効く役といえばタンヤオ、役牌、そして三色。今回 \mj{1s} があって123の三色が見込む。打 \mj{9s} \citep[][p.~64]{zero-teyaku}。

\tehai{1356m 124p 15899s 56z}[南家　1巡目][ssk-7]


\subsubsection*{ブロックより三色の種}

手牌\ref{mj:ssk-8}に6ブロックがある。\mj{8s} は孤立牌なので、牌効率通りなら打 \mj{8s}。が、\mj{2244p33z} が3ブロックとして効率が悪いので、ここは打 \mj{2p} として \mj{244p} の1ブロックにして三色の種を残す\citep[][p.~64]{zero-teyaku}。

\tehai{79m 2244789p 458s 33z}[南家　3巡目　ドラ \mj{9p}][ssk-8]

手牌\ref{mj:ssk-9}も \mj{2446m} が2ブロックとして面子を作るのが効率的に悪い。打 \mj{2m} によって1ブロックにして \mj{7p} を残して789の三色を狙う\citep{gendai-ssk}。

\tehai{2446789m 7p 113578s}[][ssk-9]


\subsubsection*{6ブロックの選択}

手牌\ref{mj:ssk-10}は789の三色部分に7枚があるが両面両面なので実は確定していない。そして \mj{8m} を落とすと \mj{233457m} という両面嵌張ができて、牌効率通りの打法。よって、どっちを切るのはドラ枚数次第。ドラがないなら打 \mj{3m} として三色と一通を見て、ドラが1枚以上なら牌効率の通りに打 \mj{8m} \citep{gendai-ssk}。

\tehai{2334578m 789p 5578s}[][ssk-10]

手牌\ref{mj:ssk-11}は6ブロックがある。6ブロック理論によると \mj{4688p} は2ブロックにしないで、\mj{688p} の1ブロックにする方が効率的だが、\mj{34m} 両面固定によって上がり率と三色の可能性が高まる。打 \mj{3m} \citep[][p.~65]{zero-teyaku}。

\tehai{334m 2344688p 3488s}[][ssk-11]

手牌\ref{mj:ssk-12}は三色不確定の形。両面と対子しかないなので、ここは三色に拘らない。受けが被っている \mj{67p} を落とす\citep{gendai-ssk}。

\tehai{06m 34467p 1112367s}[][ssk-12]

%手牌\ref{mj:ssk-13}には567の三色が確定しているが、タンヤオがあるので（平場なら）三色に拘る理由はない。牌効率の通りに打 \mj{5m} \citep{gendai-ssk}。
%ssk-11もタンヤオがあるのになぜssk-11は三色狙う？
%\tehai{5788m 223457p 3357s}[][ssk-13]


\subsubsection*{孤立牌同士の比較}

手牌\ref{mj:ssk-14}は一向聴、すでに \mj{789m789p} があって、\mj{9s} にくっつけば三色が確定する。が、他の孤立牌は両面できやすい \mj{3m3p}。よって、この手牌の打牌は巡目とドラ枚数次第。ドラがないなら \mj{3m3p} を切って789の三色を狙う。ドラが2枚以上あるなら牌効率の通りに打 \mj{9s}。ドラが1枚ならどちらもいい。巡目が早ければ3色を狙い、遅ければ牌効率に従う\citep{yoteru-ssk}。

\tehai{3789m 3789p 229s 333z}[南家][ssk-14]

手牌\ref{mj:ssk-15}は\ref{mj:ssk-14}より聴牌まで遠い手牌。同じ \mj{9s} にくっつけば三色が確定する。ここでドラがないと三色を狙うという方針が\ref{mj:ssk-14}と同じだが、ドラが3枚以上なら、仕掛けを効かせるために三色を狙う価値が高まるので打 \mj{3m} \citep{yoteru-ssk}。

\tehai{3789m 3379p 3559s 33z}[南家][ssk-15]

手牌\ref{mj:ssk-16}にもう4ブロックがある。牌効率として \mj{7s} から両面ができるので \mj{1s} より強いが、\mj{3p} を先に引かなければ嵌張リーのみになってしまう。よって、ここは \mj{1s} を残して打 \mj{7s} \citep{gendai-ssk}。

\tehai{123567m 1255p 17s56z}[][ssk-16]


\subsubsection*{三色の種にも強さの差がある}

\textbf{真ん中の種は弱い}\citep[][pp.~65--66]{zero-teyaku}。手牌\ref{mj:ssk-17}の \mj{3p} は下の三色の種、\mj{7p} はドラなのでどっちも切らない。そうすると \mj{14m} のどっちを切る問題になる。打 \mj{1m} と打 \mj{4m} の比較は表~\ref{tab:ssk-1} の通り、\mj{3p} に \mj{124p} がくっついたら差がでて、特にツモ \mj{1p} の先に打 \mj{1m} の場合は役なしとなり、打 \mj{4m} の場合は三色が確定するなのでカバーできる。打 \mj{4m} で123の三色を狙う。

\tehai{123405m 37p 122334s}[南家　5巡目　ドラ \mj{7p}][ssk-17]

\begin{table}[h]
  \centering
  \begin{tabular}{lll}
    役 & 打 \mj{1m} & 打 \mj{4m} \\
    \hline
    ツモ \mj{1p} & なし & 三色 \\
    ツモ \mj{2p} & (三色)・ピンフ & (三色)・ピンフ \\
    ツモ \mj{4p} & (三色)・ピンフ & ピンフ
  \end{tabular}
  \caption{手牌から \mj{14m} のどちらを切る選択肢の比較。括弧は三色が確定しない場合。}
  \label{tab:ssk-1}
\end{table}

手牌\ref{mj:ssk-18}と\ref{mj:ssk-19}も三色の種がある牌姿。この種を切ると愚形立直ドラ1になる。先のように比較をすると、\mj{7m} にピンフが付かないと三色もない。一方、\mj{6m} にピンフがつくなくても \mj{8m} 引いたら三色が確定する。よって、手牌は打 \mj{7m} として立直。手牌は \mj{13p} を落とす\citep[][p.~66]{zero-teyaku}。

\tehai{1237m 13678p 67778s}[南家　3巡目　ドラ \mj{1m}][ssk-18]

\tehai{1236m 13678p 67778s}[南家　3巡目　ドラ \mj{1m}][ssk-19]


\subsubsection*{両天秤}

「両天秤」は複合しない二つの役の準備状態を指す。どっちの役になるのかがツモ次第。三色と一通の両天秤があり、異なる三色の両天秤もある。手牌\ref{mj:ssk-20}は打 \mj{4m} にすると、\mj{5678m67p67s} の形で \mj{5p} 引いたら打 \mj{8m} で567の三色が見えて、\mj{8s} 引いたら打 \mj{5m} で678の三色が見える。つまり、567と678の三色の両天秤。牌効率の通りなら打 \mj{67p} として受け入れが最大22枚だが、打 \mj{4m} にしても受け入れも19枚あるのでその3枚損のリターンが大きい\citep{rb1}\footnote{もちろん、点棒状況・巡目・ドラ枚数によって三色狙わない場面もある。}。

\tehai{45678m 3367p 34567s}[][ssk-20]

手牌\ref{mj:ssk-21}も123と234の三色の両天秤。打 \mj{3p} は打 \mj{14s} に対して4枚損だが、立直・ドラ・三色の方が立直・ドラ・ピンフより高いので相当なリターン\citep{clearrain-ssk}。

\tehai{23067m 23399p 1234s}[6巡目][ssk-21]

手牌\ref{mj:ssk-22}も345と456の三色の両天秤。今回はリーのみになる可能性があるので、三色の両天秤を狙う価値がもっと高まる\citep{gendai-ssk}。打 \mj{4p}。

\tehai{3456m 445 999p 45 88s}[][ssk-22]

手牌\ref{mj:ssk-23}は打 \mj{2m3p} にすると567と678の三色の両天秤になるが、打 \mj{58s} に対して実は11枚損。暗刻ありヘッドレスの一向聴なので、良形メンタンが確定している。よって、この11枚の損失は大きすぎて、三色の両天秤にしない。567と678の三色のどちらを選ぶと、8の方がでやすくて、678の方が完成しやすいので打 \mj{5s} \citep{clearrain-ssk}。

\tehai{22267m 33367p 5678s}[6巡目][ssk-23]

手牌\ref{mj:ssk-24}にはもうメンタンピンドラ1（＝満貫）なので三色の両天秤にしない。牌効率の通りに打 \mj{8m} \citep{gendai-ssk}。

\tehai{566678m 306778p 67s}[][ssk-24]


%---------------------------------------------
\subsubsection{一気通貫}

一通は三色と同じ、門前で2飜、喰い下がりで1飜。が、一通の完成形に一色の1から9までの数牌が必要なので、組み合わせの数（=3）は三色（=7）より少ない。従って、一通の出現確率は三色より低い。


\subsubsection*{一通が狙える時}

聴牌まで遠い階段、またはドラなしの手牌は、一色の牌は5種類あるなら一通を視野に入れる。孤立牌や愚形搭子を比較する時、一通のために牌効率を逆らうことが多い。手牌\ref{mj:ittsu-1}は孤立の \mj{7m9p8s} の比較。牌効率の通りなら \mj{9p} だが、この手牌はピンフもタンヤオもドラもない。ドラが \mj{6m} なのでドラを使うため \mj{7m} は切らない。そして、\mj{2345679p} に \mj{18p} 引いたら一通が見えるので \mj{8s} を切る\citep[][p.~68]{zero-teyaku}。

\tehai{3457m 2345679p 118s}[ドラ \mj{6m}][ittsu-1]

手牌\ref{mj:ittsu-2}は \mj{1m1s} の比較。マンズの部分は5種類があって、一通にはちょっと遠いが、それでもより孤立の \mj{1s} は先に切るべき。手牌\ref{mj:ittsu-3}も同じ理由で打 \mj{2p} \citep[][p.~69]{zero-teyaku}。

\tehai{14578m 1136p 1899s 5z}[ドラ \mj{5m}][ittsu-2]

\tehai{267899m 2678p 1399s}[][ittsu-3]

手牌\ref{mj:ittsu-4}は辺張 \mj{12m} と嵌張 \mj{13s} の比較。牌効率通りなら、ツモ \mj{4s} の両面変化を見て辺張を落とす。が、マンズの部分が6種類があって、ツモ \mj{45m} による一通チャンスを逃さないようにここは \mj{13s} を落とす\citep[][pp.~68--69]{zero-teyaku}。

\tehai{1267899m 678p 1399s}[][ittsu-4]

手牌\ref{mj:ittsu-5}は辺張 \mj{12m} と辺張 \mj{89m} と両面 \mj{45s} の比較。\textbf{一通確定愚形愚形は、何の役も確定していない愚形＋両面より強い}\citep{gendai-ittsu}。よって、打 \mj{45s}。

\tehai{1245689m 99p 45567s}[][ittsu-5]


\subsubsection*{一通を見切る時}

一色の部分に7種類があっても、1枚・2枚を払えばタンヤオがつく場合、もしくはドラはすでに1枚ある場合、一通を狙うことは滅多にない。例えば、手牌\ref{mj:ittsu-6}にマンズの部分が7種類。が、\mj{9m} を切っても \mj{27m} の受け入れが減らない上、さらに \mj{1m} を切るとタンヤオがつく。正着は打 \mj{9m}。手牌の形を限定する一通は、タンヤオ・ドラ1より弱い\citep[][p.~69]{zero-teyaku}。

\tehai{1345689m 33468p 06s}[][ittsu-6]

手牌\ref{mj:ittsu-7}のようなドラなしタンヤオにもならない手牌なら \mj{26p} を切って一通を見る。ドラが1枚あるならやはり打 \mj{9m} \citep[][p.~70]{zero-teyaku}。

\tehai{1345689m 11246p 78s}[][ittsu-7]

手牌\ref{mj:ittsu-8}から \mj{19m} を外すとタンヤオがあるので一通を見ない。打 \mj{19m} \citep[][p.~70]{zero-teyaku}。

\tehai{1345679m 3557p 455s}[][ittsu-8]

手牌\ref{mj:ittsu-9}から辺張 \mj{12m} を落とすとピンフがつくので一通は見ない。打 \mj{1m} \citep{gendai-ittsu}。

\tehai{1256789m 99p 45567s}[][ittsu-9]


\subsubsection*{黄金の一向聴}

くっつき一向聴の中で、三色と一通の両天秤の牌姿は「黄金の一向聴」と呼ばれている。例え手牌\ref{mj:golden-1}は打 \mj{9s} にすると、ツモ \mj{16p} で一通が見えて、ツモ \mj{79s} で789の三色が見える。打 \mj{8m} のピンフより価値が高いので、ここは打 \mj{9s} \citep{rb1}。

\tehai{2345789p 556789s 8m}[][golden-1]

手牌\ref{mj:golden-2}も黄金の一向聴。黄金一向聴は必ず\ref{mj:golden-1}のように明確に浮いている \mj{8m} があるわけではない。くっつき一向聴なら\ref{mj:golden-2}のような四連形＋両面対子の形にもなりうる。

\tehai{2346789p 234m 344s}[][golden-2]

手牌\ref{mj:golden-3}も打 \mj{6s} による黄金の一向聴。打 \mj{3m} による四連形＋四連形のくっつきの受け入れが53枚で一番広いが、そのうち愚形リーのみになってしまうツモも多い\citep{shibukawa-ssk}。

\tehai{344m 2345789p 3456s}[東1局　東家　5巡目　ドラ \mj{2z}][golden-3]

手牌\ref{mj:golden-4}は打 \mj{6m} とすると345の三色とソーズの一通の両天秤。が、四連形にピンフツモは \mj{2457m} 4種類があって、\mj{3p} へのくっつきの \mj{24p} より多い。この手牌は三色や一通を特に見なくてもメンタンピンも十分なので打 \mj{3p} \citep{gendai-ittsu}。

\tehai{345688m 3p 2345678s}[][golden-4]

ですが、対子の部分を \mj{9m} に変えるとタンヤオが消えるので三色や一通の重要さが大きくなって、打 \mj{6m} の方が優勢になる\footnote{シミュレーションによる結論}。

\tehai{345699m 3p 2345678s}[][golden-5]


\subsubsection*{6ブロックから一通と三色の決断}

三色と一通がともに3ブロックあって、両天秤にならないためどちらを選ぶときは\textbf{三色にすると有利になることが多い}\citep{gendai-ittsu}。一通は確定しづらい点と、タンヤオと複合しやすい点からです。

手牌\ref{mj:ittsu-10}は789の三色とピンズの一通がともに見える。が、この手牌は6ブロック、\mj{78m} と \mj{23p} と \mj{56p} のどちらを落とさなければならない。この手から一通を成立させるには \mj{14p} 両方が必要なので、まず一通は確定しづらい。\mj{78m} を落として \mj{1p} を引いても待ちが \mj{47p} になって、実は弱い。よって、この手牌は一通を狙わない。打 \mj{5p} とすると、ツモ \mj{6s} で打 \mj{9s9p} として678の三色へ移行することにより、タンヤオも見える。

\tehai{3378m 2356789p 789s}[東場　北家　7巡目　ドラ \mj{8s}][ittsu-10]

手牌\ref{mj:ittsu-11}はピンズの一通が確定している。この手から両面を落とす後、他の両面が先に埋まってしまうとカン \mj{8p} の5200になる。一方、打 \mj{9p} により123と234の三色の両天秤にすることができるし、ツモ \mj{4m4p} ならタンピン三色まで狙える。この手牌も一通を狙わないで、打 \mj{9p}。

\tehai{23m 12345679p 2377s}[東場　西家　8巡目　ドラ \mj{3p}][ittsu-11]

%---------------------------------------------
\subsubsection{トイトイ}

トイトイに喰い下がりはないため、ドラや他の役に複合すると5200や8000になりやすい。高打点を作るために相性がいい役だが、立直を受けている時自分の手牌の種類は限定されていて守備力は低下\citep{yoteru-toitoi}。門前や七対子の準備状態からトイトイの移行は「鳴き判断」一章を参照。

\subsubsection*{1300から5200へ}

手牌\ref{mj:toitoi-1}は打 \mj{6m} とすると打点が1300しかない。トイトイにして、聴牌枚数が半分になるが打点が四倍になる。よって、ここは打 \mj{5m} \citep[][p.~57]{zero-teyaku}。

\tehai{566m 22288p 222s -666'z}[ドラ \mj{1z}][toitoi-1]

2副露でもトイトイに受けた方がいい。手牌\ref{mj:toitoi-2}も打 \mj{5m} \citep[][p.~57]{zero-teyaku}。

\tehai{566m 88p 222s -22'2p666'z}[ドラ \mj{1z}][toitoi-2]

両面聴牌に満貫があるならトイトイで跳満にする必要はない。聴牌枚数を半分にしても打点が$8000\to12000$、つまり1.5倍しかないので、ここは打 \mj{6m} \citep[][p.~57]{zero-teyaku}。

\tehai{566m 88p 222s -22'2p666'z}[ドラ \mj{2p}][toitoi-3]

手牌\ref{mj:toitoi-4}はツモ \mj{1s} で4つの対子ができて、トイトイにすると5200になる。ここは \mj{3s7p7m} の順番に切っていく\citep[][p.~57]{zero-teyaku}。

\tehai{788m 799p 113s 33z -666'z}[ドラ \mj{8s}][toitoi-4]

手牌\ref{mj:toitoi-5}はツモ \mj{2p} の場面。トイトイを意識しないと早い巡目でツモ切ってしまうがち。ここは打点のために打 \mj{3p}。4巡目以降は打 \mj{2z} としてトイトイと手なりの道を残す\citep[][p.~58]{zero-teyaku}。

\tehai{8899m 1223p 233z -666'z}[東場　南家　3巡目　ドラ \mj{8s}][toitoi-5]

手牌\ref{mj:toitoi-6}はトイトイとして対子が一つ足りていない。手なりで \mj{2m} を切りたいが、もう一つの対子ができたら5200に近くなるので、まず対子に手をかけない。打 \mj{7p} にすると \mj{8p} の受け入れがなくなるが、\mj{9p} が出やすくなる。\mj{3z} としてもトイトイへの道は残っているが、安全度と重なった時の鳴きやすさが優れているので残したい。よって、ここは打 \mj{7p}。もう一つの対子ができたら \mj{45s} を払っていく\citep[][p.~58]{zero-teyaku}。

\tehai{223m 11799p 45s 3z -77'7z}[ドラ \mj{8s}][toitoi-6]


\subsubsection*{ゴーニー理論}

\textbf{待ちが苦しい8000より両面5200}、というのは福地誠氏が提唱するゴーニー理論\citep{yoteru-toitoi}。例え手牌\ref{mj:toitoi-7}はドラドラ・中・トイトイで満貫だが、待ちが \mj{5p5s} 双碰なのであまり出てこない。\mj{47s} にすると打点が5200になっても待ちが8枚になる。ゴーニー理論によるとここは打 \mj{5s}。

\tehai{111m 05p 056s -999'm 77'7z}[][toitoi-7]

一方、手牌\ref{mj:toitoi-8}もドラ・中・トイトイで満貫。両面にすると打点が2600にすぎないので、ここはトイトイを狙う。打 \mj{6s}。

\tehai{111m 05p 556s -88'8s 77'7z}[][toitoi-8]


\subsubsection*{トイトイの決め打ち}

もっと遠くて、対子も５つ揃っていない階段でトイトイを狙う基準は\citep[][p.~59]{zero-teyaku}、

\begin{itemize}
\item 打点が5200以上、つまり役牌などの役に複合している
\item 対子が４つ、そのうち尖張牌の対子が１つ以下
\end{itemize}

手牌\ref{mj:toitoi-9}は三向聴だが両面が一切なしで実は受け入れがひどく狭い三向聴。真っ直ぐに行けば \mj{7z} ノミで1000にすぎない。よって、ここはトイトイを狙う。そうすると、重なったら鳴きやすいと高い牌を優先する。つまり、尖張牌＞28＞19＞字牌の順番に牌を切る。手牌からポンできる牌を全部鳴いて \mj{46p} を切る。

\tehai{2288m 1146p 18s 5677z}[ドラ \mj{6s}][toitoi-9]

手牌\ref{mj:toitoi-10}もトイトイを狙う。対子の隣の牌は早ければ早く切るほどいい。よって、ここは打 \mj{7m}

\tehai{22799m 1179p 1s 5677z}[ドラ \mj{6s}][toitoi-10]

手牌\ref{mj:toitoi-11}は \mj{9m} をポンしたばかりの場面。トイトイに向かう際で考えるのは自分の河がどう見えている。ポン \mj{9m} した後 \mj{3p} を切ることで相手を混一色を思わせて、\mj{1p} が出やすくなる。よって、ここは打 \mj{17m} ではなく打 \mj{3p}。\mj{9m} ではなく \mj{1p} からポンすると打 \mj{7m} になる。

\tehai{1227m 1139p 677z -99'9m}[ドラ \mj{6s}][toitoi-11]


%---------------------------------------------
\subsubsection{七対子}

\subsubsection*{三対子から七対子へ}

七対子は一向聴の時受け入れが9枚しかないので、通常は狙うものではないし、決め打ちもしない。対子が３つ持っていても七対子のルートを残す場面は手牌\ref{mj:chit-1}のような嵌張だらけのバラバラ牌姿。対子が３つなので七対子三向聴で受け入れが21枚。面子手でも \mj{3p5z} どちらを切ると受け入れが \mj{15m8p456s3z} 22枚になる。肝心なのは、面子手で進むと愚形立直・ドラ1で済む可能性が高い。それに対して、七対子にすると最大立直・ツモ・七対子・裏裏の倍満まで見える。よって、ここは対子をほぐさず孤立している \mj{3p} を切る\citep[][pp.~82--83]{zero-teyaku}。

\tehai{1146m 379p 3557s 335z}[東場　北家　3巡目　ドラ \mj{9p}][chit-1]

手牌\ref{mj:chit-1}で次巡で \mj{9p} をツモって\ref{mj:chit-2}になる。ここでも打 \mj{37s} として七対子に決め打ってはいけない。面子手ルートも並行して進めるので、打 \mj{4m} \citep[][p.~83]{zero-teyaku}。

\tehai{1146m 799p 3557s 335z}[東場　北家　4巡目　ドラ \mj{9p}][chit-2]


\subsubsection*{好形ばかりなら七対子は基本的無視}

手牌\ref{mj:chit-3}のような対子ばかりの手、面子手両向聴で七対子一向聴。ですが、対子にフォロー牌があって、面子手にしても聴牌までそんなに遠くない。対子のフォロー牌を外すか、対子をほぐすか、というのはこの小節の課題。

手牌\ref{mj:chit-3}には役牌暗刻があって、仕掛けが効いている。対子が多くて受け入れは多少狭くて苦しいが、まずわざと七対子のために役牌暗刻を外すことはない。役牌暗刻以外の部分には対子と両面対子しかいない場合、七対子を無視して両面を固定する。手牌\ref{mj:chit-4}のように嵌張対子があったらそこからフォロー牌を外して七対子手と面子手の並行をみる。\mj{1p} を切ると4枚の \mj{2p} がロスになるのに対して \mj{7m6p} 6枚で七対子聴牌するので、受け入れは実は増える\citep[][p.~83]{zero-teyaku}。

\tehai{788m 334677p 05s 777z}[東家　5巡目　ドラ \mj{2s}][chit-3]

\tehai{788m 133677p 05s 777z}[東家　5巡目　ドラ \mj{2s}][chit-4]

手牌\ref{mj:chit-5}は二度受けの並び対子があって、序盤なら七対子を無視して打 \mj{4m} とする。中盤以降なら打 \mj{78m} として面子手と対子手を並行させる\citep[][p.~84]{zero-teyaku}。

\tehai{334478m 22388s 777z}[ドラ \mj{1p}][chit-5]

手牌\ref{mj:chit-6}のようにマンズの上が嵌張になると巡目に関わらず打 \mj{7m} として面子手と対子手を並行させる\citep[][p.~84]{zero-teyaku}。

\tehai{334478m 22388s 777z}[ドラ \mj{1p}][chit-6]

手牌\ref{mj:chit-7}は表~\ref{tab:chit-7}の通り \mj{8p6s} のどちらを切っても受け入れが同じ6種17枚。\mj{8p} 切ったら \mj{2233p} の並び対子が両面＋両面にならないので \mj{1p} がロスになるが、\mj{1p} 引いたらタンヤオが崩すのでさほど痛くない。\mj{8p} 切りにより七対子の一向聴も残る。ツモ次第でタンヤオへ移行することもできる\citep[][p.~18]{uzaku-nankiru}。

\tehai{2233556678p5667s}[東1局　東家　5巡目][chit-7]

\begin{table}[ht]
  \centering
  \begin{tabular}{ll}
    切る牌 & 受け入れ \\
    \hline
    \mj{8p} & \mj{2347p57s} 17枚 \\
    \mj{6s} & \mj{123457p} 17枚
  \end{tabular}
  \caption{手牌\ref{mj:chit-7}の手の受け入れ。}
  \label{tab:chit-7}
\end{table}

手牌\ref{mj:chit-8}にあるドラダブ東を鳴きたいだが、もう9巡目、そして \mj{1z} を鳴いても向聴数が進まない。よって、ここは面子手と対子手を並行させる。打 \mj{56p} という決め打ちはしない。そうすると、\mj{3m} と \mj{7p} の比較は \mj{1z} を鳴いた後の形に基づく。表~\ref{tab:chit-8} の通り、打 \mj{7p} の方の受け入れが2枚多い。打 \mj{7p} とすると \mj{258p} の三面張が崩れたけれど、両面対子＋両面対子＋対子の方は二度受けの両面両面＋対子＋対子の方より受け入れが広い\citep[][p.~20]{uzaku-nankiru}。

\tehai{344m 3344567p 55s 11z}[東1局　東家　9巡目　ドラ \mj{1z}][chit-8]

\begin{table}[ht]
  \centering
  \begin{tabular}{ll}
    \multicolumn{2}{l}{\mj{44m 3344567p 55s -1'11z}} \\[3pt]
    打牌 & 受け入れ \\
    \hline
    \mj{4m} & \mj{23458p5s} 17枚 \\
    \mj{3p} & \mj{4m2458p5s} 17枚 \\
    \mj{4p} & \mj{4m2358p5s} 17枚 \\
    \mj{5s} & \mj{4m23458p} 17枚 \\
    \\
    \multicolumn{2}{l}{\mj{344m 334456p 55s -1'11z}} \\[3pt]
    打牌 & 受け入れ \\
    \hline
    \mj{4m} & \mj{25m235p5s} 19枚 \\
    \mj{3p} & \mj{245m25p5s} 19枚 \\
    \mj{6p} & \mj{245m25p5s} 19枚 
  \end{tabular}
  \caption{手牌\ref{mj:chit-8}の \mj{1z} を鳴いた形}
  \label{tab:chit-8}
\end{table}


\subsubsection*{一段目の河に工夫をする}

手牌\ref{mj:chit-9}は手なりなら打 \mj{8m}。面子手と七対子の両天秤にすると打 \mj{59s}。が、ここはまだ一段目なので、七対子一向聴の階段でアガリ率を最大化するため打 \mj{7s} の選択肢も浮上してくる。

\tehai{8899m 33445p 579s 33z}[南家　4巡目　ドラ \mj{1p}][chit-9]

手牌\ref{mj:chit-10}も打 \mj{3m} $\to$ \mj{6m} の手順で \mj{1m} が安全に見える河を作っていく。

\tehai{136m 88p 6699s 22334z}[東家　4巡目　ドラ \mj{1s}][chit-10]

そういっても、二段目に入ると河に中張牌たくさん切られていてやや不自然にになってしまうので、中盤以降は河のことを考えなくてもいい。


\subsubsection*{消極的七対子}

七対子は先手を取られると非常に弱くなる。特に中盤以降、誰かから立直がかかってくるのもおかしくない階段、守備力がある孤立牌を残すのは「消極的七対子」と呼ばれている。手牌\ref{mj:chit-11}からは打 \mj{2z} ではなく、立直を受けた後一向聴を維持するため、中張牌の \mj{6p} にすべき。振聴の \mj{2z} だが、生牌の \mj{6p} と2枚切れの \mj{2z} は重なる確率はそう変わらないし、後手を踏んだ時 \mj{6p} は切りづらい\citep[][p.~85]{zero-teyaku}。

\tehai{499m 67799p 11779s 2z}[東家　10巡目　ドラ \mj{4m}　\mj{2z} は2枚切れで振聴][chit-11]


\subsubsection*{七対子とトイトイの分岐点}

\citet{gendai-toitoi}によると、

\begin{itemize}
  \item トイトイ有利
  \begin{itemize}
    \item 暗刻がある
    \item 役牌や混一色などの役と複合している
    \item 端牌対子で仕掛けやすい
    \item オーラスで僅差でアガリトップ
  \end{itemize}
  \item 七対子有利
  \begin{itemize}
    \item 2枚切れている刻子にならない対子がある
    \item 暗刻がない
    \item 他に役がない
    \item 孤立のドラがある
    \item 対子が鳴きにくい
    \item 余裕のあるトップ目等守備を重視する場合
  \end{itemize}
\end{itemize}

手牌\ref{mj:chit-toi-1}から \mj{8s} を切るのは牌効率通りの思考で、完全一向聴も見える。が、打点が1000にすぎない。アガリトップなどの場面以外なら打 \mj{34p} として七対子やトイトイの可能性を最大限にする。6枚損だが、両向聴なので打点上昇のためにいいトレードオフと思う。ちなみに打 \mj{4m} も七対子とトイトイの道が残っているが、両面対子が消えていて一向聴の形は悪くなる\citep{gendai-toitoi}。

\tehai{33477889m 34p 88s 77z}[][chit-toi-1]

手牌\ref{mj:chit-toi-2}はトイトイ有利の条件と七対子有利の条件が満たされている。七対子ドラ単騎が上がりにくいという点から考えると、トップ目ではないと打 \mj{8s} としてトイトイに向かう\citep{gendai-toitoi}。

\tehai{33388m 1155p 8s 4566z}[ドラ \mj{4z}][chit-toi-2]

\begin{itemize}
  \item トイトイ有利
  \begin{itemize}
    \item 暗刻がある
    \item 役牌対子がある
    \item 尖張牌対子が１つ
  \end{itemize}
  \item 七対子有利
  \begin{itemize}
    \item 孤立のドラがある
  \end{itemize}
\end{itemize}

2暗刻＋4対子の形も、役牌と複合しているか、七対子で即リーできる牌の枚数で判断する。手牌\ref{mj:chit-toi-3}は、\mj{3m8s} とすると \mj{8m15p6z} の他の牌がツモったら七対子聴牌になる。打 \mj{5p} とするとポンテンできるし四暗刻まで見える。今回は \mj{66z} の対子があるので、トイトイに向かう方がいい。

\tehai{33388m 1155p 888s 66z}[ドラ \mj{4z}][chit-toi-3]


\subsubsection{三暗刻}

三暗刻聴牌の立直判断に関しては\ref{sec:riichi-snk}を参照。

\subsubsection*{対子３つ＋暗刻１つなら三暗刻を視野に}

手牌\ref{mj:snk-1}は中ぶくれの \mj{2334m} で真ん中の \mj{3m} を引いてしまう結果。\mj{3m} ではなければ \mj{89p} を払ってピンフ一盃口が見えるが、\mj{3m} なら対子系への距離は順子系より近い。ツモ \mj{9m8p9p} でツモり三暗刻、他の部分に対子がもう一つできたら七対子やトイトイも見える。よって、ここは \mj{4m} $\to$ \mj{2m} で切っていく\citep[][p.~73]{zero-teyaku}。

\tehai{2333499m 5678899p}[南家　5巡目　ドラ \mj{3z}][snk-1]

手牌\ref{mj:snk-2}は手なりでいくと \mj{4m8s} が、愚形リーのみがほぼ確定。ここも手なりで進行せず、打 \mj{2m} で三暗刻を狙う。打 \mj{9s} ではないのは安全度のため。手なりで進んでもいいのはドラが2枚以上の時\citep[][p.~74]{zero-teyaku}。

\tehai{244m 12388p 111889s}[南家　5巡目　ドラ \mj{3z}][snk-2]

手牌\ref{mj:snk-3}は節\ref{sec:complex-shape}に触れられたケースと似ていて、両面対子ばかりの手牌は対子のフォロー牌を外して一向聴の広さを瞬間的受け入れより優先する。縦引きで成り立つ三暗刻から考えるとフォロー牌外しも有利。ここは安全度と赤受けを考えて打 \mj{4p} とする\citep[][p.~74]{zero-teyaku}。

\tehai{233m 455p 23677s 777z}[ドラ \mj{3z}][snk-3]


%============================================
\subsection{4ブロック打法}

\textbf{聴牌まで遠い序盤なら辺張・嵌張より役・打点に関連する孤立牌を残すのもあり}\citep[][pp.~158--159]{suphx}。このような5ブロックを4ブロックにするのは「向聴戻し」ともいう。序盤で向聴戻しの有効局面は下記のパターンの一つ\citep{gendai-shantenback}：

\begin{enumerate}
\item 強い形（4連形、中ぶくれ）$\times$2
\item 強い形＋打点増加孤立牌
\item 打点増加孤立牌$\times$２
\item 愚形にフォローがある
\item 愚形が枯れている（2枚切れ）
\item 一向聴に取ると愚形安手確定
\end{enumerate}

手牌\ref{mj:rst-1}は混一色が見えるのでここで外嵌張の \mj{79m} を外す。

\tehai{79m4456789p33445z}[東1局　西家　6巡目　ドラ \mj{1z}][rst-1]

手牌\ref{mj:rst-2}はチャンタや三色が見えるのでここで \mj{77s} を外す。

\tehai{112379m9p77789s44z}[東1局　西家　6巡目　ドラ \mj{1z}][rst-2]

手牌\ref{mj:rst-3}で \mj{46m} には \mj{33m} のフォローがあるので \mj{6m} を切っても \mj{5m} の受けは逃さない。

\tehai{334699m36778p567s}[東1局　西家　6巡目　ドラ \mj{1z}][rst-3]

手牌\ref{mj:rst-4}ではそのまま一向聴をとると外嵌張や辺張待ちになるのでここで \mj{89m} を落とす。

\tehai{1389m224566p4567s}[東1局　西家　6巡目　ドラ \mj{1z}][rst-4]

手牌\ref{mj:rst-5}は \mj{7m} が枯れているなら \mj{89m} を落とす。\mj{7m} まだ4枚あるならそのままで \mj{6p} を切る。

\tehai{23489m11566p4567s}[東1局　西家　6巡目　ドラ \mj{1z}][rst-5]

手牌\ref{mj:kanchan-1a}にはタンヤオがつかないため、辺張を落としても打点はよほどよくならなさそう。聴牌のときの最終形の想定すれば、辺 \mj{7p} を避けるもっといい待ちを作る\citep[][pp. 48--55]{tetsusenryaku}。

\tehai{4678m 3589p 11067 -7s}[東1局　西家　3巡目　ドラ \mj{2m}][kanchan-1a]

ドラ3で打点はもう十分である場合は、素直に聴牌に向かうべき。ですが、ソウズの部分が \mj{114067s} になったら伸びやすいので辺張を外すのはおすすめ。

\tehai{4678m 3589p 11067 -7s}[東1局　西家　3巡目　ドラ \mj{1s}]

手牌\ref{mj:block-11}に愚形がばかり。そのままで立直すれば大体2600に過ぎない。点数が欲しい場面ならチャンタを見込んで、\mj{5s} の対子を落とし、\mj{2z} 重なりや \mj{4z} 槓で打点を最大限にする\citep[][pp. 74--79]{tetsusenryaku}。

\tehai{123 89m 79p 155s 244 -4z}[南4局　南家　1巡目　ドラ \mj{1s}][block-11]

Suphxは\ref{mj:block-9a}から \mj{13m} を落とした\citep[][pp.~158--159]{suphx}。これらのブロックをキープしたら、仕掛けるとほぼ2飜、門前も結構厳しめ。\mj{0p} のくっつきにより1飜の打点上昇や、\mj{48s7z} のくっつきにより混一色への転移は十分見込める

\tehai{137m 0p 2248s 22667z -9m}[東2局　北家　1巡目　ドラ \mj{8m}][block-9a]

\textbf{打点上昇を諦めて愚形確定かつ低打点の一向聴を取ることはない}\citep[][pp.~160--161]{suphx}。手牌\ref{mj:block-10a}は一向聴だがそのままで聴牌を取るとドラ \mj{5m} や \mj{0s} が出ていくのは確定。Suphxはここで \mj{5p} を落とした。ただし、\mj{0s} が \mj{5s} となる場合、向聴を戻しても打点が必ず上昇するわけではないため、向聴を戻さず立直に寄る。

\tehai{588m 1135p 130s 777z -2s}[南2局　北家　5巡目　ドラ \mj{5m}][block-10a]

麻雀は4面子1雀頭を作るゲーム。打点を上げるためや、よりいい形を求めるために瞬間のブロック不足は状況により許容されるが、三向聴・四向聴の手牌の受けがどんなに広いでも聴牌までのツモ巡数は必ず増えていく。前述の4ブロックにする打法はいずれも打点・手役に関連するので、いつもドラを貪って愚形を拒否するわけではない。

%============================================
\subsection{一向聴(1)：種類}

麻雀は14枚の牌で「4面子＋1頭」をできるゲームなので、聴牌の時は「頭なし4面子」（単騎・延べ単）と「3面子＋1頭＋1搭子」（辺張・嵌張・両面）もしくは多面張。一向聴は、前述の形をもう一回崩して、その種類は下記の通り\citep{tomo-iishanten}。

\begin{table}[ht]
  \centering
  \begin{tabular}{lllll}
    & 面子 & 対子 & 他の搭子 & 浮き牌 \\
    \hline
    余剰牌一向聴 & 2 & 1 & 2 & 1 \\
    完全一向聴 & 2 & 1 & 2 & 1 \\
    ヘッドレス一向聴 & 3 & 0 & 2 & 0 \\
    くっつき一向聴 & 3 & 1 & 0 & 2
  \end{tabular}
  \caption{一向聴の種類}
\end{table}


\subsubsection*{余剰牌一向聴}

余剰牌がある一向聴は、面子が二つ、搭子（対子含み）が二つがある上、あと一つの孤立牌が浮いている状態を指す。その余剰牌の機能は下記の通り三つある：

\begin{enumerate}
  \item 安全牌
  \item 良形変化
  \item 打点上昇
\end{enumerate}

\tehai{12389m12366p78s-?}[][ist-1]

手牌\ref{mj:ist-1}の \mj{?} が \mj{4z} の場合、誰にも不要な字牌を抱えて守備へ寄っている。\mj{4s} の場合、\mj{35s} 引いたら良形ができ \mj{89m} の辺張を払う。\mj{2s} の場合、\mj{13s} 引いたら下の三色が見える。余剰牌が出ているため、受け入れが狭い。手牌\ref{mj:ist-1}の受け入れは \mj{7m69s} 最大3種12枚。\mj{89m} が \mj{78m} になっても受け入れが4種16枚に過ぎない。

手牌が良形確定なら余剰牌を安牌としてを抱えていく。愚形があるなら自分の手で優先して良形変化への余剰牌を持っていく。


\subsubsection*{完全一向聴} 

\tehai{123789m23566p88s}[][ist-2]

\tehai{123788m23456p88s}[][ist-3]

完全一向聴は\ref{mj:ist-2}のように、その浮き牌がある両面搭子に重なってできる両面対子がある形。その受け入れが、ある両面搭子が両面対子のともに二度受けになるという最低の場合でも5種16枚。両面搭子が順子に繋がって三面張ができる場合、例え\ref{mj:ist-3}のように、受け入れが \mj{689m147p8s} 7種23枚となる。

\tehai{22m566788p34678s}[][ist-3-1]

\tehai{688m22p567p34678s}[][ist-3-2]

手牌\ref{mj:ist-3-1}が含んでいる \mj{566788p} の形があれば完全一向聴になる。受け入れは \mj{2478p25s} 6種19枚。この形は \mj{567p}＋\mj{688p} で分析できるが、19枚の受け入れになる「順子＋嵌張対子」のはこの形だけ。例え\ref{mj:ist-3-2}のように \mj{688p} が \mj{688m} になると受け入れが \mj{78m2p25s} 5種16枚となる。

\tehai{12378m223456p88s}[][ist-3-3]

手牌\ref{mj:ist-3-3}のように、ピンズの部分が \mj{223456p} か \mj{233456p} か、両面対子が順子に繋がった場合受け入れが7種23枚となる。

\subsubsection*{ヘッドレス一向聴} 

\tehai{123789m23567p78s}[][ist-4]

\tehai{111789m23567p78s}[][ist-5]

\tehai{111789m34567p78s}[][ist-6]

ヘッドレス一向聴は\ref{mj:ist-4}のように全ての搭子には対子がない形。\mj{23p78s} のどちらが先に重なったら良形聴牌になるが、\mj{14p69s} 先に引いたら単騎になってしまうリスクがある。従って、手牌\ref{mj:ist-4}の受け入れが \mj{1234p6789s} 8種28枚となるが、良形聴牌できるツモ牌はただ4種12枚。

暗刻が含まれているヘッドレス一向聴は必ず良形聴牌になる。例え\ref{mj:ist-5}のように、\mj{14p69s} 引いても \mj{1m} 切って残りの両面が十分に働く。よって、\ref{mj:ist-4}と\ref{mj:ist-5}の受け入れが同じ8種28枚であるが、良形聴牌率によると後の方が圧倒的に有利。そして\ref{mj:ist-6}のようにある両面が順子につながって三面張ができたら受け入れが一気に11種37枚になるので、非常に強い形である。


\subsubsection*{くっつき一向聴}

\tehai{12355789m3789p3s}[][ist-7]

\tehai{12355789m1789p1s}[][ist-8]

くっつき一向聴は\ref{mj:ist-7}のように、二つの浮き牌と頭に関連する牌が引いたら聴牌できる形。その受け入れが11種40枚（\mj{5m} + \mj{12345p} + \mj{12345s}）あり、一番聴牌になりやすい形であるが、良形聴牌できる牌は \mj{24p24s} 4種16枚だけ。さらに\ref{mj:ist-8}のように浮き牌が端牌になる場合、受け入れが7種24枚になり、有効牌引いても必ず愚形聴牌になる。

%============================================
\subsection{一向聴(2)：選択}

ここまでの話は全ての搭子に愚形がないという前提の上で成立する。愚形があったら受け入れがさらに減って、不十分と思える形。聴牌になりにくくて、聴牌になっても高い確率で愚形聴牌になるので、\textbf{不十分の一向聴の場合で考えるのは良形変化と向聴戻し}\citep{manboo-iishanten}。

\tehai{122357m5678p2245s}[][ist-9]

手牌\ref{mj:ist-9}では \mj{2m58p} どちら切っても受け入れが \mj{6m36s} 3種12枚になるが、\mj{58p} を残して \mj{4679p} を引いたら良形になるのでここで \mj{2m} を切る。

\tehai{1223467p2345778s}[][ist-10]

手牌\ref{mj:ist-10}では \mj{258s} どちら切っても受け入れが \mj{358p} 3種11枚になる上、一枚の \mj{3p} が既に手牌にあるので、巡目がまだ深くないなら打 \mj{2p} で両向聴に戻ってよりいい一向聴を取るのは得。

\tehai{4578999m45p11145s}[][ist-11]

手牌\ref{mj:ist-11}では \mj{45m}・\mj{45p}・\mj{45s} の三択。どちらを切っても受け入れに差が出ないけれど、\mj{45m} を落とした場合 \mj{6m} を引いたら暗刻ありヘッドなしの広い一向聴になるのでそっちの方がより有利。

\tehai{4668m34445p23777s}[][ist-14]

手牌\ref{mj:ist-14}では、打 \mj{6m} の受け入れ（\mj{57m14s} 4種16枚）が打 \mj{8m} の受け入れ（\mj{56m4p14s} 5種15枚）により1枚多いけれども、打 \mj{8m} の後に \mj{2356p} 引いて \mj{4m} 切りで完全一向聴ができるので、受け入れの枚数に差が大きくない場合良形変化を優先させる。

\tehai{1234m 5678p 122234s}[][ist-28]

\begin{table}[ht]
  \centering
  \begin{tabular}{ll}
    打牌 & 受け入れ \\
    \hline
    \mj{1s} & \sfrac{14}{46}：\mj{123456m3456789p3s} \\
    \mj{1m} & \sfrac{7}{20}：\mj{3456789p12345s}
  \end{tabular}
  \caption{手牌\ref{mj:ist-28}の受け入れ}
  \label{tab:ist-28}
\end{table}

手牌\ref{mj:ist-28}は表~\ref{tab:ist-28} の通り打 \mj{1s} の受け入れが最も広いが、ピンフや一盃口に向かうなら打 \mj{1s} にする\footnote{\mjfn{2s} がドラなら打 \mjfn{1s} にして一番早い速度で聴牌をとる。}。\mj{122233s} の形は \mj{123s} + \mj{223s} に見えるのでピンフ・一盃口になりやすい形\citep[][p.~14]{uzaku-nankiru}。

\tehai{678m 123306p 12378s}[][ist-29]

\begin{table}[ht]
  \centering
  \begin{tabular}{ll}
    打牌 & 受け入れ \\
    \hline
    \mj{3p} & \sfrac{8}{28}：\mj{4567p6789s} \\
    \mj{12p} & \sfrac{4}{16}：\mj{47p69s}
  \end{tabular}
  \caption{手牌\ref{mj:ist-29}の受け入れ}
  \label{tab:ist-29}
\end{table}

手牌\ref{mj:ist-29}は表~\ref{tab:ist-29} の通り打 \mj{3p} とするヘッドレス一向聴の受け入れが最も広いが、半分以上の確率で望ましくない単騎待ちになる。\mj{3p} を対子にして \mj{12p} を落としても \mj{47p69s} の受け入れも消えない上、ツモ \mj{6p} で \mj{0p} が出ていくのを避けるためにここでヘッドレス一向聴にしない\citep[][p.~16]{uzaku-nankiru}。


\subsubsection*{強い待ちを残す}

\tehai{22356799p234556s}[][ist-12]

手牌\ref{mj:ist-12}では \mj{2p} と \mj{5s} の二択。どちらの受け入れも同じだけど、\mj{556s} に \mj{5s} の受け入れを残した場合 \mj{5s9p} の双碰受けに引いたら三面張が崩れてしまうのでここで打 \mj{5s} として三面張を固定する。

\tehai{223m22789p445s -7'77z}[][ist-31]

手牌\ref{mj:ist-31}では \mj{2m} と \mj{4s} の二択。\mj{2m} の方が仕掛けやすい（聴牌できる確率がより高い）が、\mj{14m} 待ちの和了できる率がより高い。どっちにせよ、聴牌率は常に上がり率より高いので、上がり率を重視してここで \mj{2m} を切る。


\subsubsection*{役による打牌判断}

選択肢同士の受け入れに差が出ない場合、役で判断する。

\tehai{67m45688p3345556s}[][ist-13]

\begin{table}[ht]
  \centering
  \begin{tabular}{ll}
    打牌 & 受け入れ \\
    \hline
    \mj{3s} & \sfrac{7}{22}：\mj{58m8p2457s} \\
    \mj{6s} & \sfrac{7}{20}：\mj{58m8p2345s} \\
    \mj{5s} & \sfrac{6}{19}：\mj{58m8p347s}
  \end{tabular}
  \caption{手牌\ref{mj:ist-13}の受け入れ}
\end{table}

手牌\ref{mj:ist-13}の場合、\mj{5s} を切って \mj{334455s} の一盃口になる可能性があるが打 \mj{3s} の受け入れと良形聴牌率で比べると不利なのでここで一盃口を見切って打 \mj{3s} とする。

\tehai{12355m1334556p23s}[][ist-15]

手牌\ref{mj:ist-15}では打 \mj{1p} と打 \mj{3p} の受け入れが同じだが、打 \mj{1p} の方に一盃口チャンスがあり、一方打 \mj{3p} の方に三色とピンフの可能性がある。よって、ここで打 \mj{3p} とする。

\tehai{11155m1334556p23s}[][ist-16]

手牌\ref{mj:ist-16}は\ref{mj:ist-15}の変化形、三色もピンフもつかない。この状況なら \mj{1p} を切って、ツモ \mj{5p} で完全一向聴を期待する。

\tehai{1222m112378p5567s}[][ist-19]

手牌\ref{mj:ist-19}に打 \mj{1m} に受け入れが20枚、打 \mj{1p} に21枚、打 \mj{5s} にも21枚ある。打 \mj{1m} とするとピンフにならない上 \mj{13m} のような良い待ちを失うのでここで論外。亜両面は内側の方に良形変化の枚数が多いのでここで \mj{1p} を切る。

%\tehai{12357889m067p556s}[][ist-32]
%
%\begin{table}[ht]
%  \centering
%  \begin{tabular}{ll}
%    打牌 & 受け入れ \\
%    \hline
%    \mj{8m} & \sfrac{10}{34}：\mj{34567m45678s} \\
%    \mj{9m} & \sfrac{5}{16}：\mj{68m457s}
%  \end{tabular}
%  \caption{手牌\ref{mj:ist-32}の受け入れ}
%  \label{tab:ist-32}
%\end{table}
%
%手牌\ref{mj:ist-32}は一通（\mj{5m} + \mj{789m}、\mj{8m} が不要）と三色（\mj{57m} + \mj{88m}、\mj{9m5s} が不要）の選択\footnote{そう言ってもどちらの役はまだ確定されていない。}。\mj{8m} を切ったら \mj{5m6s} のくっつき一向聴なので表~\ref{tab:ist-32} 打 \mj{9m} の方より受け入れが広い。よって、\citet[p.~18]{uzaku-nankiru}は打 \mj{8m} にする。
%
%\begin{table}[ht]
%  \centering
%  \begin{tabular}{lll}
%     & 価値 & 確率 \\
%    \hline
%    打 \mj{8m} & 2600 & 50\%（ツモ \mj{37m48s}） \\
%     & 3900 & \minibox{38.2\%（ツモ \mj{5m57s} \\ + ツモ \mj{46m} $\times$ 50\%）}\\
%     & 7700 & 11.8\%（ツモ \mj{46m} $\times$ 50\%）\\
%    打 \mj{9m} & 2600 & 50\%（ツモ \mj{8m45s}）\\
%     & 3900 & 12.5\%（ツモ \mj{6m} $\times$ 50\%）\\
%     & 8000 & \minibox{37.5\%（ツモ \mj{7s} \\ + ツモ \mj{6m} $\times$ 50\%）}
%  \end{tabular}
%  \caption{手牌\ref{mj:ist-32}の価値の期待値}
%  \label{tab:ist-32-ev}
%\end{table}
%
%実際それぞれのツモとアガリ形の価値の期待値を求めると表~\ref{tab:ist-32-ev} の通り三色の方の期待値が高いのがわかる。問題になるのは受け入れを34枚を16枚にすることによるツモ巡数の増加による局収支はどれくらい下がるか。役が確定されていない上、\mj{8m} を切って両面聴牌になっても待ちが打 \mj{9m} の方に明らかに優れていないので、やはり場面によるか\footnote{東1局自分西家7巡目平場の場面なら打 \mj{8m} にする、それより点数が欲しい場面なら打 \mj{9m}、と。}。


\subsubsection*{愚形含み完全一向聴}

普通には安牌を持つためにも良形完全一向聴を余剰牌一向聴にはしない\footnote{受け入れが枯れている場合なら安牌を残す。}\citep{gendai-iishanten-choice-1}。単独搭子が愚形である場合、選択肢が下記のように三つある。

\begin{enumerate}
\item 複合搭子を両面にして余剰牌一向聴を取る
\item 浮き牌を切って愚形含み完全一向聴を取る
\item 愚形搭子を払って両向聴に戻る
\end{enumerate}

\tehai{2347m5589p123788s}[][ist-201]

手牌\ref{mj:ist-201}は \mj{7m} を切ると「愚形含み」完全一向聴になる。\mj{8s} 切りで余剰牌一向聴になる。\mj{7m} がドラそばの場合、\mj{8s} 切りでドラ絡みの嵌張・両面を期待する。そうでなければ \mj{7m} 切りで素直に愚形含み完全一向聴を取る方がいい。

\tehai{2347m4699p123788s}[][ist-21]

手牌\ref{mj:ist-21}なら愚形が内嵌張になって両面変化ができやすいので \mj{7m} がよほど重要ではないのでここで切る。

\tehai{2345m5589p123788s}[][ist-22]

手牌\ref{mj:ist-22}なら孤立牌が順子についているので良形ができやすい。よって、ここで \mj{8s} を切って6ブロックにする。

\tehai{2345m6689p123788s}[][ist-23]

手牌\ref{mj:ist-23}でさらに頭が \mj{6p} になって \mj{7p} のフォローができるので、直接に \mj{89p} を落としても \mj{7p} がロースにならないので向聴戻ししてもいい。

\tehai{2345m4699p123788s}[][ist-24]

四連形と内嵌張の組み合わせ。極稀の２３４三色のために打 \mj{5m} にする。打 \mj{8s} で6ブロックにしても良いか \mj{37p} 引きで完全一向聴にならないのは損。

\tehai{234m45589p123788s}[][ist-25]

手牌\ref{mj:ist-25}に両面対子が二つあるが、愚形が対子になる変化が期待できるので両面対子を対子にしない。\mj{0p} の受け入れと \mj{69s} という強い待ちを考えたら \mj{8s} を切る。

\tehai{33455m 668p 345667s}[][ist-26]

手牌\ref{mj:ist-26}は打 \mj{6s} ツモ \mj{7p} でヘッドレス一向聴になる可能性があるが、打 \mj{6s} の受け入れが打 \mj{8p} より4枚損。親なら、もしくは巡目が既に深いなら素直に打 \mj{8p} で愚形含み完全一向聴にする\citep[][p. 10]{uzaku-nankiru}。

\tehai{135m 12399p 123667s}[][ist-27]

手牌\ref{mj:ist-27}は打 \mj{5m} と打 \mj{6s} の受け入れは同じ16枚。今回123の三色があるので、打 \mj{5m} により三色を確定しても受け入れが減らないのでここで打 \mj{5m} にする。さらに先切り \mj{5m} によって \mj{2m} が出やすくなる\citep[][p. 12]{uzaku-nankiru}。

\subsubsection*{愚形含みヘッドレス一向聴} \tehai{234888m57p112367s}[][ist-17]

手牌\ref{mj:ist-17}で打 \mj{1s} で広いヘッドレス一向聴ができるが、ソーズの \mj{5678s} を引いたら必ずカン \mj{6p} 待ちになる。よって、巡目がまだ早いなら愚形の崩してよりいい形を待つのは有利。ただし、聴牌を遅らせる余裕がないなら素直に \mj{1s} を切るべき。

\tehai{234888m57p112389s}[][ist-18]

手牌\ref{mj:ist-18}には良形聴牌ができないので、素直に \mj{1s} を切って聴牌を目指す。その時立直するか黙聴するか考えればいい。

%============================================
\subsection{手牌のスリム化}

速度・守備力が重視されるトップ目、役なし愚形濃厚などの場面なら、手牌の受け入れを少し減らしてもあまりロスにはならない。こういった手狭に構えるのは手牌のスリム化ともいう\citep[][p. 29]{suphx}。


\subsubsection*{序盤のスリム化}

\begin{rittai}[h]
  \centering
  \captionsetup[table]{labelsep=space}
  
  \drawrittai{
    jicha/seat = 南,
    jicha/hand = 1205678m 13p 223s 2 -6z,
    jicha/river = 3z,
    jicha/point = 37500,
    %
    toimen/river = 1z,
    toimen/point = 9800,
    %
    kami/river = 9m 2z,
    kami/point = 23800,
    %
    shimo/river = 9m,
    shimo/point = 28900,
    %
    kyoku = 南2局,
    dora = 2m
  }
  
  \caption{}
  \label{rittai:slim-1}
\end{rittai}

\textbf{トップ目なら鳴ける手役の重要度は高まる}。フラットの状況なら一般の牌効率を従って字牌を先に処理するだろう。立体図~\ref{rittai:slim-1} は南2局のトップ目なので、三色と役牌の重要度と速度は愚形リーのみより高いので、Suphxはここで \mj{2s} を切った\citep[][pp.~202--203]{suphx}。

\begin{rittai}[h]
  \centering
  \captionsetup[table]{labelsep=space}
  
  \drawrittai{
    jicha/seat = 東,
    jicha/hand = 49m 788999p 1239s 1 -3z,
    jicha/river = 46^5z,
    jicha/point = 41000,
    %
    toimen/river = 6z 1m 3z,
    toimen/point = 17800,
    %
    kami/river = 47z 1^p,
    kami/point = 15000,
    %
    shimo/river = 6z 1m 9^p,
    shimo/point = 26200,
    %
    kyoku = 東4局,
    dora = 3z
  }
  
  \caption{}
  \label{rittai:slim-2}
\end{rittai}

\textbf{遠いリーのみは追わない}\citep[][pp.~204--205]{suphx}。立体図~\ref{rittai:slim-2} の手牌は数牌の横引きが狭いので、立直までは遠い上、三色・チャンタ・役牌など崩れてしまう可能性も高い。序盤にも関わらずSuphxは \mj{4m} を切った。

\begin{rittai}[h]
  \centering
  \captionsetup[table]{labelsep=space}
  
  \drawrittai{
    jicha/seat = 南,
    jicha/hand = 45579m 2206p 345s 4z -6m,
    jicha/river = 57z 8p 3z,
    jicha/point = 25000,
    %
    toimen/hand = XXXXXXXXXX -X66Xz,
    toimen/river = 9^6s 11m,
    toimen/point = 30700,
    %
    kami/river = 4z 1p 3s 5p 4^s,
    kami/point = 18300,
    %
    shimo/hand = XXXXXXXXXX -055's,
    shimo/river = 1s 43z 9s 8^m,
    shimo/point = 25000,
    %
    kyoku = 東3局,
    dora = 3p9m,
    honba = 1,
    kyoutaku = 1
  }
  
  \caption{北家は \mj{6s} の次巡で \mj{6z} カンツモ切り \mj{5s}}
  \label{rittai:slim-3}
\end{rittai}

\textbf{明確に高打点かつ早そうな相手に守備を意識すること}\citep[][pp.~206--207]{suphx}。面子手に対する「捨て牌が早そう」はなかなか当たらない。特に序盤が七対子・混一色などの手順により、そうでないパターンが多いだから。早そうに見えたとしても判断に影響するほどならない。

立体図~\ref{rittai:slim-3} で満貫聴牌の手牌ならなおさら。けれど相手が役牌を暗槓して強搭子（中張牌の搭子とかポン材として優れているの \mj{1m} 対子とか）を落とすとなればほぼ早いのは確定される。ここでSuphxは \mj{8m} の受け入れ\footnote{\mjfn{455679m} は両面嵌張。節\ref{subsec:rmkc}を参照。}を捨てて \mj{9m} を落とした。


\subsubsection*{両面対子を崩すケース}

\begin{rittai}[h]
  \centering
  
  \drawrittai{
    jicha/seat = 西,
    jicha/hand = 34789m 4p 111566s 3z,
    jicha/river = 675z,
    jicha/point = 25000,
    %
    toimen/river = 32z 8p 2^z,
    toimen/point = 25000,
    %
    kami/river = 675z 1m,
    kami/point = 25000,
    %
    shimo/river = 6z 9p 4z,
    shimo/point = 25000,
    %
    kyoku = 東1局,
    dora = 9p,
  }
  
  \caption{}
  \label{rittai:slim-4}
\end{rittai}

\textbf{両面対子でも手牌により両面が崩れやすいケースは要注意}\citep[][pp.~208--209]{suphx}。立体図~\ref{rittai:slim-4} は恵まれた手牌だが、Suphxは \mj{5s} を先に切った。頭がないため同じ \mj{25m25p} 4種16枚の一向聴だが、\mj{566s} の形を保って \mj{47s} をツモれば暗刻あり頭なしの超広い一向聴になる。

だとしても、ツモ \mj{34m34p} なら \mj{47s} 待ちはより弱い待ちなので \mj{566s} を頭にしなければならない。ツモ \mj{25m25p} なら頭がないため \mj{5s} が出ていかなければならない。つまり、直接の受け入れにも限らず、有効変化 \mj{34m34p47s} の中でも多くの場合は \mj{5s} はいらない。場に一枚切れの \mj{3z} に比べると先に切るがいい。


\subsubsection*{完全一向聴に構えないケース}

完全一向聴は普通の両面＋両面より4枚の受け入れがあるので多くの人は迷わず完全一向聴に取るのだろう。実際20枚受け入れの完全一向聴の聴牌までの平均巡目は両面＋両面より0.35巡だけ減っている\footnote{作者のシミュレーションによる数値。20枚の受け入れの一向聴は平均的に聴牌までが3.84巡かかる。一方16枚の受け入れは4.19巡。}上、手役・打点・守備力・点棒状況などの要素を加えたら「完全一向聴に取るべき」ということを過信すべきではないのがわかる。\citet[]{yuusee-complete-ist}によると、下記の条件が２つ当てはまれば完全一向聴を外す：

\begin{itemize}
\item 双碰受けが2枚以上切れている
\item 切る牌が34567
\item ピンフや三色などの手役がついている
\item 中盤
\end{itemize}

\tehai{123367m 225p 3788s -8m}[東1局　東家　4巡目　ドラ \mj{4m}][slim-5]

\textbf{聴牌まで遠いなら完全一向聴に構えない}\citep[][pp.~214--215]{suphx}。手牌\ref{mj:slim-5}にはいくつの浮き牌があって、それらの昨日は下記の通り：

\begin{itemize}
\item \mj{3m}：一盃口やドラの受け
\item \mj{5p}：\mj{3467p}の受け
\item \mj{3s}：\mj{1245s}の受け、最終形の待ちとして \mj{5p} より優れている
\item \mj{8s}：一向聴の時 \mj{788s} ＋対子の完全一向聴の構え
\end{itemize}

古来の完全一向聴主義なら \mj{788s} の両面搭子を保って \mj{5p} を切ることをよしとする。\mj{5p} を切ると \mj{3467p} の受けがなくなる上、中盤で完全一向聴をとるとすると安牌がない可能性が高くて、暗刻をツモってリーのみになる可能性もある。そのためSuphxは \mj{8s} を外した。

\tehai{223789mm 99p 12356s 4z}[南家　8巡目　ドラ \mj{3m}][ist-20]

手牌\ref{mj:ist-20}は \mj{4z} を切ると完全一向聴になるが、\mj{2m} はドラ表示牌として既に1枚減っていて、点数が欲しい場況ならリーのみを避けるためここで打 \mj{2m} とする\citep[][p. 19]{ishibashi-black}。







\end{document}